\documentclass[11pt,reqno]{amsart}
\usepackage[T1]{fontenc}
\usepackage{lmodern,amsmath,amssymb,mathtools,microtype,booktabs,array,longtable}
\usepackage[margin=1in]{geometry}
\usepackage{graphicx}
\usepackage[hidelinks]{hyperref}
\ifdefined\pdfinfoomitdate\pdfinfoomitdate=1\fi
\ifdefined\pdftrailerid\pdftrailerid{}\fi
\ifdefined\pdfsuppressptexinfo\pdfsuppressptexinfo=15\fi
\numberwithin{equation}{section}
\newtheorem{theorem}{Theorem}[section]
\newtheorem{proposition}[theorem]{Proposition}
\newtheorem{lemma}[theorem]{Lemma}
\newtheorem{corollary}[theorem]{Corollary}
\theoremstyle{definition}\newtheorem{definition}[theorem]{Definition}
\newtheorem{example}[theorem]{Example}
\theoremstyle{remark}\newtheorem{remark}[theorem]{Remark}
\newcommand{\E}{\mathbb E}\newcommand{\Prb}{\mathbb P}
\newcommand{\R}{\mathbb R}\newcommand{\one}{\mathbf1}
\newcommand{\Q}{\mathcal Q}\newcommand{\TV}{\mathrm{TV}}
\newcommand{\cp}{\mathrm{cp}}\newcommand{\ev}{\mathrm{ev}}\newcommand{\st}{\mathrm{st}}
\DeclareMathOperator{\tr}{tr}\DeclareMathOperator{\dist}{dist}
\DeclareMathOperator{\law}{Law}\DeclareMathOperator{\rank}{rank}
\DeclareMathOperator{\supp}{supp}
\emergencystretch=2em

\hypersetup{pdftitle={General Theta Foundations I: Acquired Geometry and Causal Resource Transfer},pdfauthor={Qian Qi}}
\title[Acquired geometry and causal resource transfer]{General Theta Foundations I:\\Acquired Geometry and Causal Resource Transfer}
\author{Qian Qi}\date{October 2026}
\subjclass[2020]{Primary 60G40, 62C20; Secondary 47A65, 81P15, 90C22, 94A29}
\keywords{causal experiment, predictive quotient, finite memory, task comparison, stopped prediction, quantum instrument, principal angles, minimax}
\begin{document}
\begin{abstract}
We study prediction from information physically acquired by a causal experiment under expected-cost and finite-memory constraints. For compact continuous families of finite-dimensional classical and quantum instruments, we derive a common proper fallback from a finite-depth dark-subspace test. Finite model covers and common realization plans then yield explicit lower and upper bounds on optimal bounded task risk at the same physical budget. The comparison is controlled by the tested policy's actual pre-call occupation, while a separately bounded cost excess is removed by a counted initial mixture. A preliminary cost hierarchy constructs such a fallback at every strictly interior feasible budget. Decision approximation preserves output width, and compact algebraic primitives admit terminating finite-description synthesis. The resulting robustness modulus concerns the constrained optimum, not just a given controller; exact examples show why testing only model vertices, omitting remaining-horizon costs, or removing budget slack is invalid. A complementary task-order theorem transfers acquired geometric profiles through actual stopped occupation, separating relative calibration from acquisition-amplified additive defects. Exact width defects, an atomic zero-risk spectrum, finite-obstruction duality, reachable-support potentials and report-resolved continuation criteria identify further structural boundaries. Finite reset instruments and two projective detectors supply optimal strategies and explicit spectral memory laws. The fixed-law predictive quotient, post-command quantum realizations and nonintervening sensing retain their respective hypotheses. Neither the compact-family synthesis nor a special spectral formula asserts that every causal experiment has the same geometric profile or the same optimal computational resources.
\end{abstract}
\maketitle
\section{Introduction}
A causal experiment acquires reports through physical preparations and interventions. A prediction must be fixed before its future audit, and failed or rejected observations still consume resources. The central problem of General Theta Foundations I is to connect the predictive information actually acquired by such an experiment to the performance of an implementable finite-resource observer. This requires more than quantizing a prescribed probability law: one must establish an all-strategy converse, construct one program that works in every admitted model, and count its persistent memory without supplying an unobserved state or a free clock.

This paper develops two complementary routes. The \emph{geometric route} identifies an interpretable resource curve from unnormalized task comparisons on actual histories. The \emph{primitive optimization route} starts with the physical instruments and constructs lower and upper approximations to their constrained optimal risk, without assuming that a geometric comparison profile is exact. The former applies to broadly prepared causal kernels with stated certificates. The latter admits a compact continuum of unknown fixed models on a finite-dimensional, finite command/report interface. Neither route silently supplies the assumptions of the other.

\subsection*{Primitive synthesis with an unknown operating model}
Consider a compact continuous family of time-homogeneous classical or quantum instruments with common legal commands and productive reports, bounded nonnegative terminal loss, and call costs between two fixed positive bounds. The initial physical state and the fixed model are unknown. A common dark subspace in any one model prevents proper all-initial-state operation. Conversely, its absence throughout the compact family supplies a single full-support randomized command row with a uniform finite cost, by a quantitative finite-depth test. This uses established quantum reachability geometry \cite{Ying14pre}; compact continuity and the explicit spectral margin make its model-uniform cost certificate available.

A finite set of candidate models is not enough by itself. Even when every primitive entry is affine in a parameter on an interval, testing only its endpoints can give a zero-risk certificate for a problem whose true robust risk is $1/2$; Example~\ref{ex:vertex-risk} proves this for next-bit prediction with a fixed two-value decoder. A robust certificate needs both the finite test and a bound on what was omitted.

Theorem~\ref{thm:compact-risk} supplies that bound. A depth-$T$ common realization plan optimized on a certified model cover has value $q_{F,C,T}(N)$. Its constraints charge the continuing mass for at least one more physical call. Completing the plan by a derived fallback and mixing once at initialization yields a finite-label observer at exactly the original expected-cost budget $N$, simultaneously for every model and initial state. The enclosure is
\[
 q_{F,C,T}(N)\le R_{\Theta,C}(N)
 \le q_{F,C,T}(N)+z_{F,T}
       +\frac{Le_{F,T}}{N-\kappa+e_{F,T}},
 \qquad N>\kappa.
\]
Here $\kappa$ is the actual uniform fallback cost, $z_{F,T}$ is a bounded-task error, and $e_{F,T}$ is the completion and model-induced cost excess. Their explicit formulas are given in \eqref{eq:compact-budget-error}--\eqref{eq:compact-risk-error}. The denominator records the real cost slack needed to restore the budget; it is not a formal rescaling of the loss.

The key comparison, Lemma~\ref{lem:occupation-perturbation}, integrates one-step instrument discrepancies against the tested controller's actual pre-call occupation. At cost $N$, bounded terminal loss incurs at most $L(N/c_*)\epsilon_I+\epsilon_F$ model error, rather than a generic $T\epsilon_I$ error. Cost still contains a remaining-horizon factor, and Example~\ref{ex:cost-horizon} realizes it exactly. Joint refinement with $T\epsilon_I\to0$ therefore gives convergent same-budget risk bounds. A fixed interior slack gives an $O(\sqrt\epsilon)$ upper continuity modulus; its optimal exponent is not asserted. The scalar boundary example proves that continuity can fail at the least feasible budget even when both perturbed instances remain feasible.

The preliminary cost enclosure in Theorem~\ref{thm:compact-cost} converges from both sides to the least feasible uniform cost and constructs a cheaper finite fallback whenever the requested budget is strictly larger. Thus the risk result covers \emph{every interior feasible budget}, not only those above an arbitrary exploratory policy's cost. Finite decision nets keep the same output-width constraint. For compact semialgebraic parameter and decision sets with polynomial primitive entries, Corollary~\ref{cor:compact-algebraic} gives a terminating exact-description synthesis procedure and proves computability of the least feasible cost. This is not a polynomial-time claim or a finite-memory attainment theorem at the boundary.

These conclusions concern the constrained optimum, not only the evaluation of a supplied policy. Their construction does not use a balancing group, a physical reset, a model-identification oracle, or a normalized-filter contraction. The returned observer sees only its actual reports and counted labels. A model fixed for the whole run is essential: allowing an adversary to change it after every call defines a different experiment, as Example~\ref{ex:compact-limits} shows.

\subsection*{Actual-history transfer and acquired geometry}
For a genuinely prepared comparison law, let $f(x)=\E_0[Z\mid x]$, $w(x)=\E_0[\|Z-f(x)\|^2\mid x]$, and let $\nu$ be its report law. The marked profile
\[
 \Q_m(b)=\min_{|C|\le m,\ 0\le h\le b,\ \int h\,d\nu=1}
            \int[w+\dist(f,C)^2]h\,d\nu
\]
is population trimmed quantization with retained mass $1/b$ and a residual-loss mark \cite{CGM97,GG99}. Neither this functional nor ordinary rejection sampling is claimed as a new principle. Its transfer to controlled experiments requires a commandwise full-current-report comparison in one actual operating model, an all-program effort lower, actual-occupation defect control, and a common all-model implementation with counted live and export roles.

Lemma~\ref{lem:comparison} constructs the accepted comparison measure by replacing only the current row inside the competitor's actual history law. It never runs that competitor under a fictitious reference experiment where it could fail to terminate. Theorem~\ref{thm:transfer} then gives the converse and common conditional-hazard upper without a minimax exchange. Relative task order preserves exact zero-task directions and all resolution scales; absolute one-call error can instead grow linearly with the number of acquired samples. The task is fixed in this comparison, not changed between its two directions.

Exact finite-width defects in Theorems~\ref{thm:width-defect} and~\ref{thm:quadratic-defect} avoid charging a finite decoder for an unrestricted family of readouts. They distinguish ordinary squared prediction from the larger affine-square matrix family; equality of their zero-defect orders is not an unqualified Blackwell comparison. The width amplification $\min(m,k)/(5k)$ is sharp in its explicit dimension-$k$ family. Theorem~\ref{thm:width-transfer} extends the transfer to bounded decision-continuous losses, including Bregman prediction. Theorem~\ref{thm:zero-spectrum} determines the exact finite-width zero-risk thresholds from the largest atomic masses of the zero-residual predictive measure, correcting the temptation to use one favorable atom as a necessary threshold when other zero-loss atoms are present.

\subsection*{Structure, memory and exact realizations}
The full-history finite-obstruction theorem and its dual characterize common bounded-score synthesis on finite-branching interfaces. The reachable-support potential theorem gives necessary and sufficient cost and score certificates at a fixed classical memory budget. Reusable task identities require a different object, the report-resolved Hermitian span, not the average returned state or every state on its Hilbert support. The finite-depth continuation criterion includes standard Borel reports and a physically realizable sharp depth $d^2-1$ for one initial preparation. Existing Hermitian reachability methods are explicitly credited \cite{Tzeng92,Fu24}; neither their dimension count nor generic quantum control is newly invented here.

The finite reset-instrument theorem derives Bellman continuation, a spectral minimax preparation and equality-based memory certificates from the declared failure branches. It includes any finite family of rank-one detectors in arbitrary finite dimension. Separately, two arbitrary-rank projective detectors have a complete integer memory frontier: for $P,Q\ne0$ and $c=\|PQ\|<1$, the full-history value is $1+[2(1-c)]^{-1}$, the one-live-label value is $1+(1-c)^{-1}$, and the exact two-label value undergoes the transition in Theorem~\ref{thm:two-label-exact}. Three fixed-start live labels attain the unrestricted value. The two-label minimax gap shows why a sharp least-favorable preparation need not survive a nonconvex memory restriction. Principal-angle geometry and first-detection theory have their own established antecedents \cite{Halmos69,Kessler21,DelVecchio26}.

The two stopping interfaces remain distinct. An event-stopped observer can reuse a live label for terminal output because the public final event ends computation. A state-stopped observer has an intrinsic fixed halt map and disjoint live/output types. Their counts are respectively $\max(\ell,t)$ and $\ell+t$. A physical port, continuing phase, retained random seed or stopping clock is not an extra free input. Program data, scratch, numerical precision and physical time are separately charged; a finite-label existence bound does not optimize those resources.

The article also includes non-power, regular, stratified and atomic acquired geometry, singular task matrices, noncommuting post-command realizations, calibrated composition, and a fixed-law compact-exhausted predictive quotient. The complete accompanying mathematical companion contains nonintervening directional sensing and earlier exact stopping and absolute-calibration proofs \cite{Companion}. Independent finite-order measurement, two-stage reset communication, rank-constrained discrimination and spectral/dual developments keep their original source identities. They are not inferred from the present stopping-risk approximation. In particular a common executable predictive quotient for an arbitrary unknown-model family, a general characterization of when the geometric profile is exact, and matching program/precision/workspace/time lower bounds remain separate research questions. The current theorems advance the original Foundations problem without identifying any one special realization with its completion.

\subsection*{Direct comparison with causal sampling and sequential experiment order}
The density cap and conditional-hazard construction have direct causal-sampling antecedents, not only a connection to trimming. For i.i.d. proposals of law $\nu$ and a stopping index $K$ returning the current sample with law $\lambda$, the standard predictable-count argument gives
\[
 \lambda(A)\le\E K\,\nu(A),\qquad
 \E K\ge\|d\lambda/d\nu\|_\infty
             =2^{D_\infty(\lambda\|\nu)}.
\]
The last equality uses base-two divergence. This is the exact runtime lower bound in Goc--Flamich \cite[Theorem III.1]{GocFlamich24}; it is not new here. If $d\lambda/d\nu\le b$, the ordinary acceptance probability $(d\lambda/d\nu)/b$ realizes $\lambda$ at mean index $b$. Consequently the i.i.d. fixed-dictionary special case of $\Q_m^\ell(b)$ is optimization of the task loss over this classical feasible sampling set. Liu--Verd\'u \cite{LiuVerdu18} study moments of causal and noncausal selected indices and sampling--covering duality. The profile alone, its selection threshold and its i.i.d. density-cap converse are therefore not asserted as independent new principles.

The distinction in the transfer theorem is its full-current-report, actual-history comparison measure for an arbitrary competing controller under one actual operating model, coupled to a different common implementation under every model. This measure need not be the law of any executable reference stopping rule. Decoder cardinality, retained live labels and code length of a shared-randomness index are different resources. Goc--Flamich's channel-simulation divergence bounds index entropy; neither $m$ nor the total number of controller labels in this article is that entropy. Arbitrary noncausal selection of a past proposal is not admitted merely because its final output distribution agrees.

The report cites the first version of Li's binary-experiment preprint. Its checked successor is Li--Liu, \emph{Ranking Experiments under Sequential Sampling}, arXiv:2608.19897v2, dated 28 August 2026 \cite{LiLiu26}. Their Theorem 1 treats generic pairs of finite, full-support, pairwise identified \emph{multistate} experiments; it equates the stated large-budget stopping and small-cost decision orders with strict inequalities of every directed pairwise KL divergence. Theorem 2 gives exact reproduction of arbitrary finite-output stopping policies with a statewise KL-ratio leading cost and a uniform square-root remainder. Restricting that current comparison to binary experiments would be inaccurate.

That theorem ranks repeatable i.i.d. sources and reproduces entire terminal experiments. Our localized fixed-loss task order is not Blackwell order, Le Cam deficiency, their generic sequential order, or a complete quantum sufficiency criterion. Theorems~\ref{thm:width-defect} and~\ref{thm:quadratic-defect} identify exact defects for their specified tasks, while Theorem~\ref{thm:primitive-risk} gives a finite-budget approximation for controlled state-changing instruments with counted finite implementations. It does not prove a competing KL characterization or the full sequential-conversion theorem.

The invariant-subspace obstruction used for the new fallback is also established quantum reachability machinery \cite[Theorem 3.6]{Ying14pre}. The new statement is the resulting explicit same-budget risk sandwich with a derived fallback, a common policy across finitely many unknown models and a fixed output width. Its proof separates the positive survival mass bound from the generally unavailable uniform expected cost-tail bound. This specifies a mathematical assertion for independent priority and significance assessment, not a claim to originate compactness, sequence form or quantum reachability.


Model perturbation is itself an established problem. Kara and Y\"uksel \cite{KaraYuksel20} prove continuity of optimal discounted costs and robustness of policies under appropriate transition-kernel convergence, including total-variation estimates. Demirci, Kara and Y\"uksel \cite{Demirci25} quantify perturbations of transition and observation kernels, filter kernels and discounted or average-cost control. We use their version 3, whose title and formulation differ from the original preprint. Neither a policy-uniform comparison nor avoidance of an explicit belief-space computation is claimed here as a new general idea.

The additional issue in Theorem~\ref{thm:compact-risk} is simultaneous preservation of an undiscounted terminal-task risk and the \emph{same} uniform expected stopping-cost budget under compact model uncertainty. A finite grid alone is not a feasibility certificate. Our bounds use the tested plan's actual source occupation, restore its budget by a charged common fallback, and retain one finite dictionary. The algebraic corollary removes externally supplied model grids for its specified polynomial class, but uses classical real-algebraic elimination rather than a new efficient elimination method. The explicit boundary discontinuity and affine-primitive vertex counterexample explain why discounted continuity or a convex parameter set alone does not settle this constrained stopping question. These are precise distinctions for priority assessment, not an assertion that all neighboring control theory is subsumed.

\subsection*{Two complementary routes and their assumptions}
The central object is the resource-risk theorem joining an all-strategy comparison on actual histories to one common finite-role implementation, not the static trimming functional in isolation. The inherited geometric route and the new primitive optimization route are separated in Figure~\ref{fig:assumptions}. The latter computes the actual constrained risk for a compact, continuously parameterized, time-homogeneous, globally legal instrument family without supplying the reference or a reset. Its output is a sequence of bounded-budget policies and risk intervals, not a closed geometric exponent for every such family.

The constructive predictive quotient in Appendix~\ref{sec:quotient} concerns realization of the declared executable future-test equivalence under its compact-exhaustion hypotheses. Neither the proof of the task-order transfer nor the finite common realization plans use a quotient oracle or store a model-dependent posterior. The appendix is necessary for that separate predictive-representation conclusion, not an unstated premise of the risk bounds. Likewise the state space used to compute analytic operator certificates is not counted as a controller-owned quantum memory.

\begin{figure}[p]
\centering\small
\setlength{\fboxsep}{8pt}
\fbox{\parbox{0.87\linewidth}{\centering\textbf{Prepared physical kernels and paid calls}\\
Legal commands; full current reports; prospective audit; one common observer; declared halt semantics.}}
\par\medskip $\Downarrow$ \par\medskip
\begin{minipage}[t]{0.46\linewidth}
\centering\textbf{Geometric transfer route}\\[5pt]
\fbox{\parbox{0.88\linewidth}{One actual model, fixed before the competitor; full-report mass/task comparisons on every reached history; all-competitor effort and occupation bounds.}}
\par\medskip $+$ \par\medskip
\fbox{\parbox{0.88\linewidth}{One common acquisition routine on every operating history; whole-attempt cost; task and mass comparison; distinct export values and counted live roles.}}
\par\medskip $\Downarrow$ \par\medskip
\fbox{\parbox{0.88\linewidth}{Theorems~\ref{thm:transfer}, \ref{thm:width-transfer}: two-sided task-risk transfer. Relative errors preserve zero tasks; absolute defects accumulate on actual occupation.}}
\par\medskip $\Downarrow$ \par\medskip
\fbox{\parbox{0.88\linewidth}{Actual acquired mass and sublevel covers; regular, mixed and atomic laws. Geometry gives a rate only when these mass estimates hold.}}
\end{minipage}\hfill
\begin{minipage}[t]{0.46\linewidth}
\centering\textbf{Primitive compact-model route}\\[5pt]
\fbox{\parbox{0.88\linewidth}{Compact model family; continuous raw instruments; common finite commands/reports and physical dimension; bounded losses and positive bounded call costs.}}
\par\medskip $\Downarrow$ \par\medskip
\fbox{\parbox{0.88\linewidth}{Lemma~\ref{lem:compact-fallback}: no dark subspace in any model $\Longleftrightarrow$ a uniform one-row fallback. Certify its margin on a cover; refine its cost by Theorem~\ref{thm:compact-cost}.}}
\par\medskip $\Downarrow$ \par\medskip
\fbox{\parbox{0.88\linewidth}{Common finite sequence plan and open-path cost reserve. Certified model cover; perturbation on actual source occupation. At $N>N_{\min}$ find $\kappa<N$ and restore the budget by a counted initial mixture.}}
\par\medskip $\Downarrow$ \par\medskip
\fbox{\parbox{0.88\linewidth}{Theorem~\ref{thm:compact-risk}: same-budget risk enclosure with explicit cutoff and model errors; one finite decoder and a counted finite controller. Algebraic primitives admit terminating synthesis.}}
\end{minipage}
\par\bigskip
\fbox{\parbox{0.87\linewidth}{\centering\textbf{Additional effective layer, not a consequence of Borel existence}\\
Effective cover and arithmetic certificates, or compact algebraic primitives as in Corollary~\ref{cor:compact-algebraic}. Physical samplers, program bits, workspace and time remain separately charged; no optimal complexity bound is asserted.}}
\caption{Assumption flow. The two routes are complementary, not a claim that all primitive instruments possess the geometric reference profile. The predictive quotient and the independent ordered-measurement and communication companions retain their separate conclusions.}\label{fig:assumptions}
\end{figure}
\clearpage

\subsection*{Organization of the proof}
The prepared experiment and its three resource interfaces come first. The finite-model primitive construction is then extended to compact uncertainty using the occupation comparison, same-budget completion and feasibility threshold. The geometric route follows with its full actual-history proof, exact task defects and zero-risk spectrum. Subsequent sections develop synthesis, repeated continuation, reset control, spectral memory, acquired geometry and physical realizations. The technical supplement supplies legal-domain and preparation-face arguments, audit/return compilation, numerical certificate slack and explicit size bounds. Complete predecessor mathematics and the independent companion remain source-bound parts of the review packet.

\section{Prepared kernels, acquired task laws and observers}\label{sec:data}
All hidden-state, action and report spaces are standard Borel. A fixed unknown index $\theta\in\Theta$ specifies an actual preparation and positive normalized causal kernels. Positivity means nonnegative measures, not an everywhere strictly positive density. Dependence on the entire physical past may be included in the hidden state. A legal policy sees past visible reports and its admitted randomization, but not $\theta$ or a future audit. Iteration of these kernels defines its finite-prefix laws; the consistent countable path law is their kernel extension.

Each call has a cost $c_i\ge c_*>0$. The decision to make the next call and its productive flag are predictable. Commitment is permitted only after a productive report, and occurs exactly once almost surely in each model. At that time a specified future audit $Z\in H$ is scored, where $H\subset\R^k$ is nonempty, compact and convex. The prediction is committed before that audit is observed. A prospective joint report--audit kernel is defined at every productive call, including rejected calls. The stopping coin is independent of the audit conditional on the entire current visible report and the past. Write $L=\operatorname{diam}(H)^2$; the case $L=0$ is trivial whenever a legal rule exists.

The core cost $C_\tau$ includes initial preparation charged to the core, failed calls and all productive calls through commitment. A separately executed final audit must be added to the raw budget. This convention is fixed before taking a risk infimum. Physical time, stored program constants and temporary computation are not identified with call count.

\subsection*{Three resource interfaces}
A checkpoint procedure may use its whole precommit history, but must output one element of a fixed dictionary of at most $M$ values. This counts its output width only; it is not a claim of $\log_2 M$ total working bits.

An autonomous observer has an alphabet $S$, $|S|\le M$, and one fixed parameter-independent program. At an intercall boundary its only persistent observer memory is its label $s\in S$. Its next command distribution is a fixed function of $s$. A physical port name, elapsed time, subprogram index, random seed or former report is not an extra free input. During a call its current command, current report and fresh coins are available. Whatever persists to another call must be encoded in $S$. A fixed decoder $r:S\to H$ reads only the committed label.

In the \emph{event-stopped} interface a post-report kernel emits $(s',e)\in S\times\{0,1\}$. An event $e=1$ ends the experiment and retains $s'$ for the decoder; no subsequent command is executed. This public commitment event is not a stored bit used in continuing computation. The same label may be used live on one path and terminal on another. In the \emph{state-stopped} interface a fixed map $h_*:S\to\{0,1\}$ determines halt status. The initial label is live, and type-one labels are globally absorbing terminal labels with no executed command rows. Entry into a terminal label is allowed only after a productive report. Thus live and terminal types are disjoint in this interface, and only in this interface by definition.

Write $R_i^\Theta(N,M)$, $i\in\{\cp,\ev,\st\}$, for the infimum of worst-case expected squared loss over programs with robust almost-sure commitment and $\sup_\theta\E_\theta C_\tau\le N$. An empty legal class has value $+\infty$. A program chosen once at random is not a free persistent seed. The ideal measurable class may have fixed real constants; effective implementation is stated separately in Section~\ref{sec:composition}.

\subsection*{A physically acquired comparison law}
Execute a declared prepared comparison experiment and obtain its joint report--audit probability $K_0(dx,dz)$ on $X\times H$. Its report marginal is $\nu$. The law is genuinely preparable and its specification has its own resource account. It need not be a member of $\Theta$. In particular no finite set of calibration data is silently promoted to a uniform certificate. Set
\begin{equation}\label{eq:derived}
 f(x)=\int zK_0(dz\mid x),\quad
 w(x)=\int\|z-f(x)\|^2K_0(dz\mid x),\quad
 d_r(x)=w(x)+\|f(x)-r\|^2.
\end{equation}
Disintegration gives Borel versions with $f\in H$; null sets may be completed arbitrarily. The profile is
\begin{equation}\label{eq:profile}
 \Q_m(b)=\min_{\substack{C\subset H,\ 1\le|C|\le m\\
                  0\le h\le b,\ \int h\,d\nu=1}}
        \int [w+\dist(f,C)^2]h\,d\nu,
 \qquad b\ge1,\ m\ge1.
\end{equation}
We use $+\infty$ for an infeasible cap or width. The quantities $f,w$ are task coordinates of acquired predictive information, not a complete substitute for the predictive quotient. Their derivation and the quotient construction in Appendix~\ref{sec:quotient} precede any use of local dimension.

\subsection*{Lower certificates on actual histories}\label{sec:lower-cert}
Fix \emph{one} actual model $\theta_*\in\Theta$, independent of the competing program. At each actual reached productive pre-call history $h$ and legal command $a$, write $P_{h,a}(dy,dz)$ for its conditional full current report--audit law. Let $P^0_{h,a}$ be a measurable comparison row on the same report space, with a Borel mark $x=\pi_{h,a}(y)$ of law $\nu$ and
\begin{equation}\label{eq:full-moments}
 \E^0[Z\mid y]=f(\pi_{h,a}y),\qquad
 \E^0[\|Z\|^2\mid y]=w(\pi_{h,a}y)+\|f(\pi_{h,a}y)\|^2.
\end{equation}
These are conditional identities given the \emph{full} report, not just its mark. Versions and inequalities are required on actual reachable histories, not only almost everywhere under an unrelated comparison path law.

Fix $0<l_-\le1\le u_-<\infty$. For all Borel $q:Y\to[0,1]$ and all finite-image Borel $r:Y\to H$, suppose predictable nonnegative defects $a_-(h,a),e_-(h,a)$ obey
\begin{align}
 \int q\,d(P_{h,a})_Y
 &\le u_-\int q\,d(P^0_{h,a})_Y+a_-(h,a),\label{eq:lower-mass}\\
 \int q(y)\|z-r(y)\|^2P_{h,a}(dy,dz)
 &\ge l_-\int q(y)\|z-r(y)\|^2P^0_{h,a}(dy,dz)-e_-(h,a).
 \label{eq:lower-task}
\end{align}
Independent current randomization can be adjoined to both rows: integrating inequalities for each fixed random value proves the same inequalities with the same defects. All command-conditioned lower inequalities hold before averaging the chosen command.

Let $A_\tau$ count productive calls. At this same actual model require, for every legal competing procedure,
\begin{equation}\label{eq:effort}
 \E_* C_\tau\ge\kappa_-\E_* A_\tau,\qquad\kappa_->0.
\end{equation}
At budget $N$ suppose, uniformly over all robustly feasible competitors,
\begin{equation}\label{eq:occupation}
 \E_*\sum_{j\le\tau:\ j\ {
m productive}}a_-(H_{j-1},a_j)\le A_N,
 \qquad
 \E_*\sum_{j\le\tau:\ j\ {
m productive}}e_-(H_{j-1},a_j)\le E_N.
\end{equation}
These are actual occupation certificates, not properties merely of the proposed upper program. Uniform row defects give $A_N\le(N/\kappa_-)\bar a_-$ and $E_N\le(N/\kappa_-)\bar e_-$. Lemma~\ref{lem:drift} supplies a nonuniform alternative.

\subsection*{One common upper construction}\label{sec:upper-cert}
Independently, let $\mathcal G$ be a parameter-independent program which, from each of its start histories, reaches one productive report. Its conditional expected attempt cost is at most $\kappa_+>0$ on every surviving history in every $\theta\in\Theta$. It uses $B$ live roles and returns to its start label on rejection; no reset of the physical state is assumed. It forms a virtual mark $x$, a current choice $J$ and a permitted export $T_Jr$ of its task readout. For every finite dictionary $C$ assume
\begin{equation}\label{eq:exports}
 |\{T_jr:j\in J,\ r\in C\}|\le D_0|C|,
 \qquad D_0\in\mathbb N.
\end{equation}
Every continuing phase and coin needed for this construction is included in the $B$ live roles. For fixed $0<l_+\le1\le u_+<\infty$ and $a_+,e_+\ge0$, assume for all Borel $q:X\to[0,1]$ and finite-image Borel $r:X\to H$, conditionally on every such history,
\begin{align}
 \E[q(x)]&\ge l_+\int q\,d\nu-a_+,\label{eq:upper-mass}\\
 \E[q(x)\|Z-T_Jr(x)\|^2]
 &\le u_+\int q(x)d_{r(x)}(x)\nu(dx)+e_+.
 \label{eq:upper-task}
\end{align}
The expectation here includes the within-attempt choice $J$. Unlike the lower, the upper is not required to hold for each choice separately. The choices retained after this call are nevertheless charged through the actual export or live role. These hypotheses express one common causal construction, not modelwise optimizers.

\section{A constructive finite-budget frontier from primitive instruments}\label{sec:primitive-risk}
The geometric transfer theorem identifies a geometric curve when its converse and construction certificates are available. We now give a different, intrinsic route for finite time-homogeneous experiments: derive a proper fallback from the raw instruments and compute converging bounds on the optimal task risk itself. Neither a reset state, a balancing group, a supplied reference member nor a posterior stored by the controller is assumed. The finite and all-initial-state qualifications are essential. The qualitative invariant-subspace obstruction has a direct antecedent in quantum reachability \cite[Theorem 3.6]{Ying14pre}; it is proved below to obtain the quantitative fallback required by the new budget-preserving approximation.

\subsection{Primitive interface and universal termination}
The unknown model belongs to a nonempty finite set $\Theta$. Model $\theta$ has physical space $\mathcal H_\theta$ of dimension $d_\theta$. The nonempty finite command set $\mathcal A$ is globally legal. A command has completely positive report maps $\mathcal I_{\theta,a,y}$ with normalized sum. Reports range over finite sets, with a declared productive subset $Y_a^+$. After a productive report a controller may commit a decision from a fixed dictionary $C$ or reject; after any rejected or nonproductive report the post-report state is the next state. No additional return or reset is available. An interface requiring extra physical return calls must first represent and charge those calls explicitly; the theorem is not applied by hiding them in an unpriced transition. Commitment occurs exactly once, before a future audit is observed. The audit loss for decision $r\in C$ is represented on the post-report state by
\[
 0\preccurlyeq F_{\theta,a,y,r}\preccurlyeq LI,
 \qquad 0<L<\infty.
\]
Nonnegative call costs have observables $c_*I\preccurlyeq G_{\theta,a}\preccurlyeq c^*I$, with $0<c_*\le c^*<\infty$. These are the precommit core costs. A final audit with one fixed separately charged cost is added to every total budget after the theorem; an action-dependent unmodelled terminal charge is not included by assertion. Policies share one set of command and decision rules across models and unknown initial density matrices. Let $R_C(N)$ be their least worst-model, worst-initial-state expected terminal loss at uniform expected cost at most $N$. Their precommit memory is unrestricted for this value; a finite-label implementation is derived, not assumed.

Fix any command distribution $\pi_a>0$. The fallback chooses that row independently at every call and commits at the first productive report, using one fixed $r_0\in C$. Its killed physical map is
\[
 \mathcal T_\theta=\sum_a\pi_a\sum_{y\notin Y_a^+}\mathcal I_{\theta,a,y},
 \quad B_\theta=I-\mathcal T_\theta^*I,
 \quad O_{\theta,n}=I-(\mathcal T_\theta^*)^nI.
\]
A common dark subspace in one model is a nonzero subspace annihilated by every productive effect and invariant under every Kraus operator of every nonproductive report map. This condition is independent of the Kraus representation: it says that every positive state on the subspace remains there on failure and has zero productive probability for every command.

\begin{theorem}[Primitive fallback and its sharp depth test]\label{thm:primitive-fallback}
The following are equivalent: there is one common policy proper at finite uniform expected cost from every model and initial density matrix; no model has a common dark subspace; and $O_{\theta,d_\theta}$ is positive definite for every $\theta$. In that case the displayed one-row fallback is proper. With
\[
 \delta_\theta=\lambda_{\min}(O_{\theta,d_\theta})>0,
 \qquad \kappa=\max_\theta c^*d_\theta/\delta_\theta,
\]
it costs at most $\kappa$ from every current physical state in every model. One may replace this bound by the sharper certified value
\[
 \kappa=\max_\theta\left\|(I-\mathcal T_\theta^*)^{-1}
                  \sum_a\pi_aG_{\theta,a}\right\|.
\]
Only one live label and one output value are needed; these share a label in event stopping and require two labels in state stopping. Depth $d$ cannot be replaced uniformly by $d-1$. There is no dimension-only upper bound on expected cost.
\end{theorem}
\begin{proof}
Work in one model. Put $K_n=\ker O_n$, so $K_0=\mathcal H$ and
$O_{n+1}=B+\mathcal T^*O_n$. For a Kraus representation $\mathcal T(A)=\sum_vL_vAL_v^*$, positivity gives
\[
 K_{n+1}=\ker B\cap\bigcap_vL_v^{-1}K_n.
\]
The spaces decrease. Equality at two consecutive levels makes that space invariant under every $L_v$ and annihilated by $B$, hence makes all later levels equal. There are at most $d$ strict dimension drops, so $K_d$ is stable. Since every $\pi_a$ is positive, $\ker B$ is the intersection of the productive-effect kernels, and the $L_v$ include nonzero multiples of every nonproductive Kraus operator. Thus a nonzero $K_d$ is a common dark subspace. Conversely every such subspace lies in all $K_n$. Starting in it, no history policy can encounter a productive report: induction holds commandwise and reportwise. Such a policy cannot commit legally.

If $K_d=0$, positivity yields $(\mathcal T^*)^dI\preccurlyeq(1-\delta)I$. Iteration bounds survival after $kd$ fallback calls by $(1-\delta)^k$ for every normalized input. Summing blocks gives expected calls at most $d/\delta$ and expected cost at most $c^*d/\delta$. The finite model set makes the maximum bound finite. The same estimate implies convergence of the positive cost series, whose sum is the stated resolvent; its largest eigenvalue is the exact uniform cost for the fixed row. The fallback does not read the model or physical state. Its single row and fixed decoder give the label counts.

For sharp depth, on basis $e_1,\ldots,e_d$ take one nonproductive Kraus operator $L=\sum_{j=2}^d|e_{j-1}\rangle\langle e_j|$ and a productive effect $|e_1\rangle\langle e_1|$. Then $O_n$ is the projection onto the first $\min(n,d)$ basis vectors. The input $e_d$ requires exactly $d$ unit-cost calls. For the lack of a cost bound, already in dimension one a productive probability $\epsilon>0$ with unit costs has expected time $1/\epsilon$. This proves both limitations.
\end{proof}

The qualitative obstruction is inherited mathematical machinery, not a new discovery of invariant subspaces. Its purpose here is to supply an actual common fallback and its paid constant. Absence of an invariant dark subspace on a restricted preparation is a different problem; the all-initial-state hypothesis cannot be silently removed. Finiteness of $\Theta$ likewise cannot be replaced by modelwise termination in an arbitrary family without a uniform bound.

\subsection{Finite problems and restoration of the exact budget}
For a horizon $T\ge1$, use one common realization plan on the action--report tree, as in Lemma~\ref{lem:common-plan}. At every productive report include either commitment to $r\in C$ or continuation. Let $C_{\theta,T}$, $S_{\theta,T}$ and $V_{\theta,T}$ be, respectively, its expected cost through the first $T$ calls, its open effect after the depth-$T$ decision, and its accumulated terminal-loss operator. They are affine functions of the common sequence weights, computed by composing the primitive maps. In particular $0\preccurlyeq S_{\theta,T}\preccurlyeq I$ and $0\preccurlyeq V_{\theta,T}\preccurlyeq L(I-S_{\theta,T})$.

Define the finite semidefinite value
\begin{equation}\label{eq:primitive-finite}
 q_{C,T}(N)=\min\left\{q:\ C_{\theta,T}+c_*S_{\theta,T}\preccurlyeq NI,
 \quad V_{\theta,T}\preccurlyeq qI\quad(\theta\in\Theta)\right\}.
\end{equation}
The minimum is over the common plan and $0\le q\le L$. The reserve $c_*S$ is not a cutoff declaration of success: every open branch owes another paid call. For $N\ge\kappa$, prefixes of the actual fallback show feasibility. Compactness gives attainment. Set
\[
 D_T(N)=\min\{1,(N/c_*-1)/T\},\qquad
 e_T=(\kappa-c_*)D_T(N).
\]

\begin{theorem}[Same-budget constructive risk frontier]\label{thm:primitive-risk}
Under the preceding primitive conditions with no common dark subspace, the finite values are nondecreasing in $T$ and satisfy $q_{C,T}(N)\le R_C(N)$. For every $N>\kappa$,
\begin{equation}\label{eq:primitive-gap}
 0\le R_C(N)-q_{C,T}(N)
 \le L D_T(N)+\frac{L e_T}{N-\kappa+e_T}.
\end{equation}
The upper is realized by a common proper finite-label observer with expected cost at most the original $N$ in every model and from every initial state. Thus $q_{C,T}(N)\uparrow R_C(N)$ with an explicit $O(T^{-1})$ bound at each fixed $N>\kappa$. No reset or scalar posterior dynamics are required.

Writing $a=|\mathcal A|$ and $y$ for the largest report count, a sufficient live-label bound is
\[
 2+\sum_{t=0}^{T-1}(ay)^t.
\]
The event-stopped total is at most the maximum of this bound and $|C|$; the state-stopped total is at most their sum. These are existence bounds, not optimal memory, running-time or program-bit bounds.

At $N=\kappa$ the completed finite plan still gives cost at most $N+e_T$ and loss at most $q_{C,T}(N)+LD_T(N)$. The same-budget claim at that boundary is not inferred from the displayed fraction. If $\kappa=c_*$, no budget excess is present and the same-budget bound $LD_T(N)$ holds for every $N\ge\kappa$.
\end{theorem}
\begin{proof}
A genuine cost-$N$ policy restricted to depth $T$ satisfies the reserved cost inequality, because after every open path it must make at least one further call. Its truncated terminal loss is no greater than its full nonnegative loss. This proves the lower. A feasible plan at $T+1$ restricts to one at $T$: the new calls cost at least $c_*S_T$, new terminations reduce open mass, and their loss contributions are nonnegative. The reserve and risk inequalities therefore survive restriction, proving monotonicity.

For any feasible prefix, normalization and the conditional cost floor give
\[
 C_T\succeq c_*\sum_{t=0}^{T-1}S_t
       \succeq c_*[I+(T-1)S_T],\qquad S_0=I.
\]
Here each difference of successive open effects is positive. With the reserve this gives $S_T\preccurlyeq D_T(N)I$ in every model. Complete an optimizing prefix by executing the actual fallback at every open leaf. The fallback bound holds at every post-report state, including states never reached by a reference execution. Its cost and loss satisfy
\[
 C^{\rm comp}\preccurlyeq C_T+\kappa S_T
       \preccurlyeq(N+e_T)I,\qquad
 V^{\rm comp}\preccurlyeq V_T+LS_T
       \preccurlyeq(q_{C,T}+LD_T)I.
\]
It is proper because the finite prefix is followed by a proper fallback. There is no assertion that the original policy's expected remaining cost is uniformly small.

For $N>\kappa$, choose the completed policy with probability
$p=(N-\kappa)/(N+e_T-\kappa)$ and the fallback with probability $1-p$, storing the branch in the ensuing live role. Its cost is at most $p(N+e_T)+(1-p)\kappa=N$. Its risk is at most
$p(q_{C,T}+LD_T)+(1-p)L\le q_{C,T}+LD_T+Le_T/(N-\kappa+e_T)$.
If $e_T=0$ choose the completed policy directly. This proves the upper and convergence. The boundary statements follow from the unmixed completion.

A prefix stores its current action--report node, at most $\sum_{t<T}(ay)^t$ live roles; its current report is within-call input. One fallback role and one initial mixture role suffice. At that initial call, sample the branch and its command jointly; after the report either commit or enter the selected branch's counted next role. If both branches chose the same command, their conditional branch weights are proportional to their command probabilities, independent of the physical state, since the executed report instrument is identical. Equivalently use the joint within-call coin directly. No branch bit survives outside the retained label. Commitment uses the same fixed dictionary in both branches. Lemma~\ref{lem:roles} gives the two total counts.
\end{proof}

\begin{corollary}[Compact decisions at a fixed output width]\label{cor:primitive-width}
Suppose the decision carrier $D$ is compact metric and the raw audit losses obey the common decision modulus $\omega$ of \eqref{eq:loss-modulus}. Fix width $m$. For a finite $\eta$-net $D_\eta\subset D$, let $q_{m,T,\eta}(N)$ be the least value in \eqref{eq:primitive-finite} among dictionaries in $D_\eta$ of size at most $m$. Then, for $N>\kappa$,
\[
 [q_{m,T,\eta}(N)-\omega(\eta)]_+
 \le R_m(N)\le q_{m,T,\eta}(N)+LD_T(N)
                         +\frac{Le_T}{N-\kappa+e_T}.
\]
The constructed observer uses at most $m$ terminal values. As $T\to\infty$ and $\eta\downarrow0$ this gives a two-parameter constructive characterization of the unrestricted-history finite-output risk, not just a verifier for a supplied balancing program. With algebraic primitive data and polynomial score dependence on algebraic net points, every finite test is a finite real-algebraic problem; exact real samplers and program arithmetic remain separate implementation resources.
\end{corollary}
\begin{proof}
Round each point of any competing dictionary to the same chosen nearest net point at every history. Neither stopping nor cost changes, and at most $m$ values remain. Its expected loss changes by at most $\omega(\eta)$, because terminal events have total mass one. Restrict this rounded program to a finite prefix to obtain the lower. For the upper, choose an attaining net dictionary and apply Theorem~\ref{thm:primitive-risk}; the fallback uses a point of this same dictionary. There are finitely many such dictionaries, so the optimization is finite. With the stated algebraic inputs all coefficients of the finite plan are algebraic, and semidefinite inequalities can be expressed by principal minors. This gives decidability via the already credited real-algebraic method, not an efficient uniform algorithm for arbitrary real input data.
\end{proof}

\subsection{Removing the crude fallback threshold on every interior budget}
The cost of an arbitrary positive command row may greatly exceed the least feasible cost. It is not an intrinsic resource threshold. The following preliminary optimization removes that restriction without assuming the optimal policy is already known.

Let $\kappa_0$ be the uniform cost bound of the one-row fallback and let $H_\theta$ be its exact cost operator, so $0\preccurlyeq H_\theta\preccurlyeq \kappa_0I$. Ignore terminal loss for this preliminary step and commit one fixed dummy decision. Write $N_{\min}$ for the infimum of uniform expected cost among proper common policies. For a finite prefix define
\[
 W_{\theta,T}(H_\theta)
   =\sum_{h:\,|h|=T}\gamma_{\rm open}(h)\,
                       \mathcal I_{\theta,h}^*(H_\theta).
\]
Here $\gamma_{\rm open}$ includes the decision to continue after the depth-$T$ report, and $\mathcal I_{\theta,h}$ is its raw path map. This is the exact cost of executing the fallback at the open leaves. Both $W$ and the earlier $C_T,S_T$ are affine in the one common realization plan. Set
\begin{align*}
 a_T&=\min\{a:C_{\theta,T}+c_*S_{\theta,T}\preccurlyeq aI\quad(\theta\in\Theta)\},\\
 u_T&=\min\{u:C_{\theta,T}+W_{\theta,T}(H_\theta)\preccurlyeq uI\quad(\theta\in\Theta)\}.
\end{align*}

\begin{theorem}[Constructive feasibility threshold and all interior budgets]\label{thm:interior-budget}
Both finite problems attain their values, and
\[
 c_*\le a_T\uparrow N_{\min},\qquad
 \kappa_0\ge u_T\downarrow N_{\min},
\]
with the quantitative enclosure
\begin{equation}\label{eq:cost-enclosure}
 0\le u_T-a_T\le(\kappa_0-c_*)\min\{1,(\kappa_0/c_*-1)/T\}.
\end{equation}
The least common proper cost $N_{\min}$ is attained, though no finite-memory attainment is concluded at that boundary.

For every $N>N_{\min}$ a finite search through these upper problems eventually exhibits $u_{T_0}<N$. Its optimizing prefix completed by the one-row fallback is an actual finite-label fallback of cost $\kappa=u_{T_0}<N$, valid from every initial state and model. Using it in Theorem~\ref{thm:primitive-risk} gives the same-budget risk enclosure for \emph{every strictly interior feasible budget}, not only budgets above the original one-row cost. If this improved fallback has $B$ live labels, the sufficient live bound becomes
\[
 1+B+\sum_{t=0}^{T-1}(ay)^t.
\]
All decisions may still use one prescribed dictionary; the cost search is independent of its values. Algebraic primitive data give finite real-algebraic upper tests, but no uniform efficient bit-complexity claim is made.
\end{theorem}
\begin{proof}
The prefix of a genuine proper policy satisfies the one-call reserve, so $a_T\le N_{\min}$. Restriction to shorter prefixes proves $a_T$ is nondecreasing. Every prefix completed by the fallback is proper and has exactly the cost $C_T+W_T(H)$; hence $u_T\ge N_{\min}$. Prefixes of the fallback itself give $u_T\le \kappa_0$. Extending an optimizing upper prefix by one fallback call reproduces its exact completed cost, by the identity $H_\theta=g_\theta+\mathcal T_\theta^*H_\theta$ and first-productive commitment. Therefore $u_{T+1}\le u_T$. Compactness of the finite plan space and continuity of its cost operators prove both attainments.

For a lower optimizer, the reserve proof gives
$S_{\theta,T}\preccurlyeq\min\{1,(a_T/c_*-1)/T\}I$. Positivity and $H_\theta\preccurlyeq \kappa_0I$ give $W_{\theta,T}(H_\theta)\preccurlyeq \kappa_0S_{\theta,T}$. Complete that very prefix. Its cost is at most
\[
 a_T+(\kappa_0-c_*)\min\{1,(a_T/c_*-1)/T\}.
\]
Since $a_T\le N_{\min}\le \kappa_0$, this proves \eqref{eq:cost-enclosure}. The squeeze proves both limits. For boundary attainment, the compact product of all finite nodewise decision simplexes has the finite-intersection property for the closed prefix cost constraints at budget $N_{\min}$: feasibility at the largest of finitely many horizons follows from $a_T\le N_{\min}$. A common infinite policy therefore obeys all reserved constraints. Their survival bound tends to zero, and monotone convergence gives properness and cost at most $N_{\min}$. This is the same compactness mechanism as the inherited synthesis theorem, not a finite-memory assertion.

If $N>N_{\min}$, convergence of $u_T$ implies that the strict inequality $u_{T_0}<N$ holds at some finite level. An upper optimizer is itself the finite prefix and a physically executable fallback, not a model-dependent analytic certificate without a common policy. Its start-state uniformity permits it at every later open leaf of a time-homogeneous experiment. Substitute its cost and its counted live labels in the completion and initial-mixture proof; every inequality remains unchanged, except that its $B$ live roles replace the one fallback row. A fixed dummy decision may be replaced by any point of the eventual task dictionary without changing precommit calls or cost. Exact algebraic tests decide a strictly feasible upper threshold when the coefficients and proposed budget have the stated algebraic representation. General real data still require an exact oracle or further numerical certification.
\end{proof}

For example, two orthogonal nonzero failure projections have one-row fair fallback cost $\kappa_0=2$, but fair-first alternation has uniform cost $3/2$ and stops by call two. The cost upper problem at depth two therefore supplies a fallback below every budget $N>3/2$, including $3/2<N<2$. Thus the preliminary optimization removes a genuine restriction of the crude fallback, rather than only renaming its constant.

\subsection{Robustness and scope of the construction}
If two fallback killed maps on the same space satisfy
$\|\mathcal T^*-\widetilde{\mathcal T}^*\|_{\infty\to\infty}\le\epsilon$
on Hermitian matrices, telescoping powers gives
$\|O_d-\widetilde O_d\|\le d\epsilon$: each positive subunital adjoint is a contraction in this norm. Hence $\delta-d\epsilon>0$ certifies a perturbed proper fallback and cost at most $c^*d/(\delta-d\epsilon)$, when the perturbed calls obey the same cost upper bound. This is a quantitative stability certificate for the fallback, not a claim that a terminal-law distance controls unbounded cost. The sharp ladder concerns depth; the scalar rare-success example shows why a spectral margin must enter any numerical certificate.

This section treats finite known candidate models, not statistical learning of an unrestricted unknown kernel. It constructs a common policy without a lower witness, but does not prove that its risk is always the geometric profile $\Q$. The geometric transfer and this primitive synthesis are complementary routes: one gives interpretable sharp resource laws under its certificates, the other supplies a convergent constrained optimum and finite implementation on its finite interface. Quantum reachability, sequence form, potential arguments and semidefinite separation are established ingredients; the present assertion is their explicit paid, same-budget, fixed-output synthesis with a derived fallback and all-model/all-initial-state quantifiers.

\section{Compact uncertainty: certified risk at the unchanged budget}\label{sec:compact-risk}
The preceding primitive construction uses finitely many operating models. A finite grid of an unknown continuum is not itself a robust model family: a policy that meets its budget on the grid can overspend between grid points. We now control that discrepancy from the raw instruments, including the unbounded waiting cost after completion. The resulting finite programs enclose the optimal risk for a compact model family and produce a common finite-label observer at the \emph{original} budget. No model index is revealed to the observer. The proof replaces a worst-horizon accumulation by the actual pre-call occupation of the tested policy.

Robustness and approximation of partially observed control have substantial prior theory \cite{KaraYuksel20,Demirci25,Saldi25}. The particular assertion here concerns an undiscounted stopping task, a uniform expected-cost constraint, finite output width, and a physically executable completion which restores that same constraint. It is not priority for model approximation or perturbation bounds in general.

\subsection{A compact family of primitive instruments}
Use the finite command/report interface of Section~\ref{sec:primitive-risk}, but let the unknown, \emph{fixed throughout a run}, model $\theta$ range over a nonempty compact metric space $\Theta$. All models act on the same $d$-dimensional physical space. Productive report sets and legal commands are common. Each normalized instrument is $\{\mathcal I_{\theta,a,y}\}_y$, each nonnegative call cost has observable $G_{\theta,a}$, and the prospective terminal loss at $r\in D$ has observable $F_{\theta,a,y,r}$. Here $D$ is compact metric, and
\[
 c_*I\preccurlyeq G_{\theta,a}\preccurlyeq c^*I,\qquad
 0\preccurlyeq F_{\theta,a,y,r}\preccurlyeq LI,
 \qquad 0<c_*\le c^*<\infty,
\]
with $L>0$. An observed cost or phase used by a policy must be included in the declared report; it is not another free input. Additional preparation, return and terminal-audit charges use the same explicit convention as before.

For two models define symmetric primitive discrepancies
\begin{align}
 \epsilon_I(\theta,\phi)
 &=\frac12\max_a\sup_{\rho\succeq0,\,\tr\rho=1}
       \sum_y\|\mathcal I_{\theta,a,y}(\rho)
                    -\mathcal I_{\phi,a,y}(\rho)\|_1,\label{eq:raw-model-distance}\\
 \epsilon_g(\theta,\phi)&=\max_a\|G_{\theta,a}-G_{\phi,a}\|,\qquad
 \epsilon_F(\theta,\phi)=\sup_{a,y\in Y_a^+,r\in D}
                 \|F_{\theta,a,y,r}-F_{\phi,a,y,r}\|.\nonumber
\end{align}
Assume these discrepancies tend uniformly to zero with the model distance. Assume also a common modulus $\omega(\eta)\downarrow0$ for terminal losses as decision distance tends to zero. Joint continuity on the compact model and decision carriers is one sufficient primitive condition. The report-resolved trace distance in \eqref{eq:raw-model-distance} is for classical feedback on the specified physical system; a coherent external reference or quantum controller is not included by calling it a diamond distance.

A finite set $F\subset\Theta$ is an $(\epsilon_I,\epsilon_g,\epsilon_F)$-cover when each $\theta\in\Theta$ has a $\phi\in F$ with all three discrepancies bounded by these displayed numbers. The bounds, not just the grid points, are the certificate. Compactness and the assumed continuity provide arbitrarily fine finite covers. An \emph{effective} cover requires supplied representations and certified moduli; compactness alone does not provide an algorithm on unspecified real data.

\begin{lemma}[A uniform common fallback from compact primitive data]\label{lem:compact-fallback}
There is a proper common policy of finite uniform expected cost for every model and every initial density matrix if and only if no model has a common dark subspace. In this case any one fixed full-support command row, stopped at its first productive report, is a uniform fallback. With the notation of Theorem~\ref{thm:primitive-fallback},
\[
 \delta=\min_{\theta\in\Theta}\lambda_{\min}
            [I-(\mathcal T_\theta^*)^dI]>0,\qquad
 \kappa_0=c^*d/\delta
\]
is a valid cost bound. If a finite cover has instrument error at most $\epsilon_I$, then
\begin{equation}\label{eq:certified-dark-margin}
 \delta\ge\min_{\phi\in F}\lambda_{\min}
           [I-(\mathcal T_\phi^*)^dI]-d\epsilon_I.
\end{equation}
A positive right side certifies the common fallback on the entire family, not only on the cover.
\end{lemma}
\begin{proof}
A dark subspace in even one model prevents any all-initial-state policy from committing there, by the commandwise argument of Theorem~\ref{thm:primitive-fallback}. If none exists, that theorem makes every displayed depth-$d$ observable positive definite. Its entries depend continuously on $\theta$. Compactness and continuity of the least eigenvalue therefore give a strictly positive minimum. The same block-survival proof bounds every expected cost by $c^*d/\delta$, using the one fixed row and no model-dependent choice.

For the quantitative inequality, if $0\preccurlyeq A_y\preccurlyeq HI$, instrument normalization permits subtraction of $HI/2$ in every report component. Trace duality gives, for every density matrix,
\begin{equation}\label{eq:effect-distance}
 \left|\sum_y\tr\left[A_y
   (\mathcal I_{\theta,a,y}-\mathcal I_{\phi,a,y})(\rho)\right]\right|
       \le H\epsilon_I(\theta,\phi).
\end{equation}
Indeed the sum of the traces of the differences is zero, and the remaining terms are at most $H/2$ times the sum of their trace norms. To compare killed maps on $0\preccurlyeq A\preccurlyeq I$, take $A_y=A$ on nonproductive reports and zero on productive ones, and average the common command row. Telescope the $d$ powers. Positive subunital maps preserve the interval $[0,I]$ and contract the Hermitian operator norm, so each of the $d$ differences has norm at most $\epsilon_I$. The least-eigenvalue inequality gives \eqref{eq:certified-dark-margin}. This argument uses bounded effects, not a normalized posterior or a positive likelihood floor.
\end{proof}

\subsection{Perturbation controlled by actual occupation}
For a common depth-$T$ realization plan, write $C_{\theta,T}$, $S_{\theta,T}$ and $V_{\theta,T}$ as before. After the decision following report $T$, the open branches have not yet executed another call. For a fixed initial density $\rho$, put
\[
 n_{\phi,T}(\rho)=\E_{\phi,\rho}\min(\tau,T).
\]
This is the actual number of begun calls in the finite prefix under $\phi$. The next lemma requires a budget only in this source model, not in a hybrid execution.

\begin{lemma}[Occupation-sensitive prefix comparison]\label{lem:occupation-perturbation}
For every common depth-$T$ plan, every initial density matrix and every $\kappa\ge0$,
\begin{align}
 |\tr\rho(S_{\theta,T}-S_{\phi,T})|
 &\le\epsilon_I n_{\phi,T}(\rho),\label{eq:prefix-survival-error}\\
 \left|\tr\rho\bigl[(V_{\theta,T}+LS_{\theta,T})
                   -(V_{\phi,T}+LS_{\phi,T})\bigr]\right|
 &\le L\epsilon_I n_{\phi,T}(\rho)+\epsilon_F,\label{eq:prefix-risk-error}\\
 \left|\tr\rho\bigl[(C_{\theta,T}+\kappa S_{\theta,T})
                   -(C_{\phi,T}+\kappa S_{\phi,T})\bigr]\right|
 &\le n_{\phi,T}(\rho)
       [\epsilon_g+\epsilon_I(c^*(T-1)+\kappa)].\label{eq:prefix-cost-error}
\end{align}
Here the errors may be replaced by any certified bounds for this pair of models. In particular, a source reserve $C_{\phi,T}+c_*S_{\phi,T}\preccurlyeq NI$ implies $n_{\phi,T}(\rho)\le N/c_*$ for all $\rho$. Thus the loss error need not grow as $T\epsilon_I$ on a fixed-cost prefix, whereas the cost comparison still has an explicit remaining-horizon factor.
\end{lemma}
\begin{proof}
Use backward cost or score observables on the finite history tree. At a pre-call node $h$ with $t<T$ calls already made, average its common command row. For the loss padded by $L$ on open depth-$T$ leaves, the target-model post-report continuation observable is a convex combination of the terminal score and the subsequent padded score; it lies in $[0,LI]$. For the survival indicator the corresponding observable lies in $[0,I]$. For cost padded by $\kappa$ on open leaves it lies in $[0,(c^*(T-t-1)+\kappa)I]$. These are deterministic bounds for the finite prefix, valid before any comparison of its stopping laws.

Subtract the two backward recursions by writing
\[
 \mathcal I_\theta^*A_\theta-\mathcal I_\phi^*A_\phi
 = (\mathcal I_\theta^*-\mathcal I_\phi^*)A_\theta
       +\mathcal I_\phi^*(A_\theta-A_\phi).
\]
At each node, \eqref{eq:effect-distance} bounds the first term by the corresponding continuation bound times $\epsilon_I$. Iterate the second term with the \emph{source} path instruments and source controller weights. Pairing at the initial state expresses the accumulated discrepancy as a sum over actual source pre-call occupation. The total of these masses is $n_{\phi,T}(\rho)$, including all failures and all private random choices.

For cost, the immediate observable difference adds at most $\epsilon_g$ per source call. Replacing $T-t-1$ by $T-1$ gives \eqref{eq:prefix-cost-error}. For terminal loss, the differences of the actual terminal observables add at most $\epsilon_F$ times source commitment mass through the cutoff, which is at most one. They are not charged at every rejected report. The padded terminal value at an open cutoff is the same $L$ in both models. These facts prove \eqref{eq:prefix-risk-error}; survival has no terminal-observable discrepancy. The positive source cost floor gives the last assertion.

One can equivalently obtain the recursion by first fixing all private choices at each visible history and then integrating them. The common behavioral realization is justified by Lemma~\ref{lem:common-plan}. The proof does not reveal $\phi$ to the controller, condition on a future audit, run a competitor on a fictitious reference past, or divide by a possibly zero report mass.
\end{proof}

\subsection{A finite certificate at precisely the required budget}
Fix a nonempty finite dictionary $C\subset D$ and a common fallback which can start from \emph{every} current physical state in every $\theta$, commits with a value in $C$, has uniform cost at most $\kappa$, and uses $B$ live labels. Initially Lemma~\ref{lem:compact-fallback} supplies one with $B=1$; the cost hierarchy below supplies cheaper ones. Assume $N>\kappa$.

Let $q_{F,C,T}(N)$ be the finite semidefinite value \eqref{eq:primitive-finite} with constraints only at the finitely many cover models $\phi\in F$, but with the same common realization plan. Define
\begin{align}
 D_T&=\min\{1,(N/c_*-1)/T\},\nonumber\\
 e_{F,T}&=(\kappa-c_*)D_T+
 \frac{N}{c_*}[\epsilon_g+\epsilon_I(c^*(T-1)+\kappa)],\label{eq:compact-budget-error}\\
 z_{F,T}&=LD_T+L\frac{N}{c_*}\epsilon_I+\epsilon_F,
 \qquad
 \Gamma_{F,T}=z_{F,T}+\frac{Le_{F,T}}{N-\kappa+e_{F,T}}.\label{eq:compact-risk-error}
\end{align}
A proper fallback must make at least one call, so $\kappa\ge c_*$. These errors are nonnegative.

\begin{theorem}[Same-budget risk enclosure for compact uncertainty]\label{thm:compact-risk}
Under the primitive hypotheses just stated, the finite problem is feasible and attains its value, and
\begin{equation}\label{eq:compact-risk-sandwich}
 q_{F,C,T}(N)\le R_{\Theta,C}(N)
       \le\min\{L,q_{F,C,T}(N)+\Gamma_{F,T}\}.
\end{equation}
The upper is realized by one proper finite-label controller of expected cost at most $N$ for \emph{every} $\theta\in\Theta$ and every initial density matrix. It uses one dictionary, not a modelwise family of dictionaries. With $a$ commands and at most $y$ reports per command, sufficient live labels are
\begin{equation}\label{eq:compact-labels}
 1+B+\sum_{t=0}^{T-1}(ay)^t.
\end{equation}
Event-stopped total labels are the maximum of this bound and $|C|$; state-stopped total labels are their sum.

For nested covers and increasing horizons the finite lower values are nondecreasing. If $T\to\infty$, $T\epsilon_I\to0$, and $\epsilon_g,\epsilon_F\to0$, their limit equals $R_{\Theta,C}(N)$. At a fixed slack $N-\kappa>0$, taking $T\asymp\epsilon^{-1/2}$ when all primitive errors are at most $\epsilon$ gives an explicit $O(\sqrt\epsilon)$ enclosure. This is an upper modulus, not a claim that its exponent is optimal.
\end{theorem}
\begin{proof}
The actual fallback has cost below $N$ in all models, so its prefixes prove feasibility on $F$. A finite plan lies in a compact sequence-form polytope, and the constraints are closed with $0\le q\le L$. The minimum is attained. Restriction of any feasible full policy on $\Theta$ to the finite models and horizon proves the lower. Adding cover models or extending the horizon cannot weaken that lower, by the restriction lemma.

Take a finite optimizer. At every $\phi\in F$, the reserve gives $S_{\phi,T}\preccurlyeq D_TI$ and source call occupation at most $N/c_*$. For any $\theta$ choose a covering $\phi$. Complete the same prefix on the actual experiment with the declared fallback at every open leaf. Positivity bounds the completed cost by $C_{\theta,T}+\kappa S_{\theta,T}$ and the completed loss by $V_{\theta,T}+LS_{\theta,T}$. At the covering model these are respectively at most $N+(\kappa-c_*)D_T$ and $q_{F,C,T}(N)+LD_T$. Lemma~\ref{lem:occupation-perturbation} now gives completed cost at most $N+e_{F,T}$ and loss at most $q_{F,C,T}(N)+z_{F,T}$ in \emph{every} model and initial state. The fallback is applied on the actual post-prefix state, not the state at the covering model. Its all-state guarantee is what makes this completion legal and proper.

Before the first physical call, choose the completed prefix with probability
\[
 p=\frac{N-\kappa}{N-\kappa+e_{F,T}}
\]
and the fallback alone otherwise. The choice is independent of the unknown model and input. Its robust expected cost is at most $p(N+e_{F,T})+(1-p)\kappa=N$. Its loss is at most
\[
 p(q_{F,C,T}+z_{F,T})+(1-p)L
       \le q_{F,C,T}+z_{F,T}+(1-p)L.
\]
This is the claimed gap. The actual loss is always at most $L$, giving the truncation of the upper. Both branches commit almost surely, so does their mixture. The sole initial choice is compiled into one initial row whose within-call branch choice routes to prefix or fallback labels; it is not an uncounted persistent seed. The finite tree has at most the displayed sum of live node roles, and the fallback uses $B$. The role-packing lemma proves both total-label counts.

The explicit errors tend to zero under the stated joint refinement. Compact continuity supplies arbitrarily fine model covers, so such refinements exist. For the rate, substitute $D_T=O(T^{-1})$ and bound the mixture denominator below by $N-\kappa$. The gap is at most a constant times $T^{-1}+\epsilon T+\epsilon$, yielding the asserted choice. None of these estimates assumes independence of the physical states between calls or a mean-cost tail converging uniformly over feasible policies.
\end{proof}

\begin{corollary}[Preserving an optimized finite output width]\label{cor:compact-width}
Let $R_{\Theta,m}(N)$ optimize over all dictionaries of at most $m$ points in $D$. For a finite $\eta$-net $D_\eta$, minimize $q_{F,C,T}(N)$ over tuples $C\in D_\eta^m$ with repetition, and call the result $q_{F,m,T,\eta}(N)$. The fallback may use any one value of the chosen dictionary without changing its costs. Then
\[
 [q_{F,m,T,\eta}(N)-\omega(\eta)]_+
 \le R_{\Theta,m}(N)
 \le\min\{L,q_{F,m,T,\eta}(N)+\Gamma_{F,T}\}.
\]
The common observer has at most $m$ outputs. In particular model and decision approximation do not create a hidden continuum-sized decoder or a stored model estimate.
\end{corollary}
\begin{proof}
Round each point of any feasible competitor's single dictionary to one net point, identically at every history. Its costs and stopping law do not change. Bounded terminal-observable error changes expected loss by at most $\omega(\eta)$ because the terminal events have total mass one. Restricting gives the lower. A finite net has finitely many tuples, so choose a minimizer and apply Theorem~\ref{thm:compact-risk} for the upper. Repetitions can only reduce the actual width.
\end{proof}

\subsection{The least feasible cost and every interior budget}
The primitive fallback cost $\kappa_0$ can be far from optimal. It must not become an artificial admission threshold. For each finite cover and horizon, let
\[
 a_{F,T}=\min\{a:C_{\phi,T}+c_*S_{\phi,T}\preccurlyeq aI\ (\phi\in F)\},
\]
ignoring task loss and committing a fixed dummy value. This is an attained finite semidefinite problem with $c_*\le a_{F,T}\le\kappa_0$. Put $D_T^0=\min\{1,(\kappa_0/c_*-1)/T\}$ and
\begin{equation}\label{eq:compact-cost-certificate}
 E_{F,T}=(\kappa_0-c_*)D_T^0+
 \frac{\kappa_0}{c_*}[\epsilon_g+
                  \epsilon_I(c^*(T-1)+\kappa_0)],\qquad
 U_{F,T}=a_{F,T}+E_{F,T}.
\end{equation}
Here a cover needs only instrument and cost errors. Let $N_{\min}(\Theta)$ be the least uniform expected cost of a common proper policy.

\begin{theorem}[Constructive compact-family feasibility threshold]\label{thm:compact-cost}
One has
\[
 a_{F,T}\le N_{\min}(\Theta)\le U_{F,T}.
\]
The upper is the certified cost bound of an actual common finite-label policy. Along nested covers and horizons with $T\epsilon_I\to0$ and $\epsilon_g\to0$, the lower bounds increase to $N_{\min}(\Theta)$ and the running minima of the upper bounds decrease to the same number. The least cost is attained by a full-history common policy.

For every $N>N_{\min}(\Theta)$, a sufficiently fine finite cover and sufficiently large horizon exhibit $U_{F,T}<N$. Their completed policy is a finite-label fallback with certified $\kappa=U_{F,T}<N$. Consequently Theorem~\ref{thm:compact-risk} and Corollary~\ref{cor:compact-width} apply at \emph{every strictly interior feasible budget}, with no supplied balancing program or reference witness. For effective data and certified covers this is a terminating search for a strictly feasible fallback, not a dimension-only efficient synthesis bound. Finite-memory attainment at $N=N_{\min}(\Theta)$ is not asserted.
\end{theorem}
\begin{proof}
Every actual proper policy supplies a feasible reserved prefix on the cover, proving the lower. Complete a lower optimizer with the original uniform fallback. At a covering model, its open mass is at most $D_T(a_{F,T})\le D_T^0$, its expected prefix calls at most $a_{F,T}/c_*\le\kappa_0/c_*$, and its padded cost at most $a_{F,T}+(\kappa_0-c_*)D_T^0$. Apply \eqref{eq:prefix-cost-error} with padding $\kappa_0$ to every covered model. This proves the upper and its operational realization. It bounds the cost on arbitrary current states, since the prefix problem and perturbation bound were uniform in their initial density matrices.

The error in \eqref{eq:compact-cost-certificate} tends to zero along the stated refinement. Restriction and addition of models prove lower monotonicity; the running minimum construction gives upper monotonicity without incorrectly claiming that the raw error bounds themselves decrease. The squeeze proves both limits.

For attainment, choose proper common policies with uniform costs decreasing to $N_{\min}$. The countable product of finite nodewise probability simplexes is compact. Every finite-prefix reserved cost constraint at each fixed model and state is closed in these coordinates. A finite collection of such constraints at budget $N_{\min}$ is feasible: take a convergent subnet of the selected policies' coordinates and use continuity of the finite-prefix operators and convergence of their budgets. Compactness then gives one policy satisfying all these constraints simultaneously. Its survival is at most $D_T(N_{\min})$ after $T$ calls, and the positive floor and monotone convergence prove properness and total cost at most $N_{\min}$ in every model. No uniform expected-cost tail estimate was used.

At a strictly larger requested budget, the upper sequence eventually falls below it. The completed policy attaining that certificate supplies exactly the fallback required for the risk theorem. Relabelling its dummy output to a chosen task decision changes no call, so it works for any eventual dictionary. Its finite live-state count is retained in \eqref{eq:compact-labels}. Effective cover and primitive arithmetic hypotheses are additional inputs for performing these tests; the proof does not infer them from topological compactness.
\end{proof}

\begin{corollary}[Effective synthesis from compact algebraic primitives]\label{cor:compact-algebraic}
Suppose $\Theta\subset\R^p$ and $D\subset\R^k$ are nonempty compact semialgebraic sets defined over the real algebraic numbers. Suppose the real and imaginary entries of the primitive instrument maps and costs are polynomial in $\theta$, and the terminal-loss entries are polynomial in $(\theta,r)$. Assume the common finite interface, positive cost floor and absence of dark subspaces as above. Then $N_{\min}(\Theta)$ is a computable real. For every algebraic $N>N_{\min}(\Theta)$, integer $m\ge1$ and positive algebraic accuracy $\xi$, a terminating finite procedure produces a common finite-label observer of cost at most $N$ whose risk differs from $R_{\Theta,m}(N)$ by at most $\xi$, with algebraic controller probabilities and decoder values.

The procedure also produces rigorous lower and upper value bounds. It need not be computationally efficient; the statement concerns a supplied finite polynomial description, not learning the experiment from data. Exact bounded-bit sampling, minimal program size and finite-memory optimality at the boundary are separate questions.
\end{corollary}
\begin{proof}
All assertions about positivity of the finite-dimensional Hermitian primitive matrices, normalization, and $O_{\theta,d}\succeq\delta I$ are first-order polynomial assertions over real algebraic coefficients. Represent complex entries by their real and imaginary parts, and use principal minors for positive semidefiniteness. Quantifier elimination over real closed fields \cite{Seidenberg54} therefore verifies these assertions and exhibits an algebraic $\delta>0$ valid for all $\theta$ when it exists. Compactness and the no-dark lemma guarantee existence. The positive lower cost bound and finite upper cost/loss bounds may also be chosen algebraic: compactness gives such bounds, and the same polynomial tests verify them. This gives an algebraic bound $\kappa_0$ for the common one-row fallback.

At a fixed horizon, use common sequence weights and a tuple of $m$ decoder values as variables. Composing a finite number of polynomial primitive maps gives polynomial coefficients in $\theta$. The constraints $C_{\theta,T}+c_*S_{\theta,T}\preccurlyeq NI$ and $V_{\theta,T}\preccurlyeq qI$ for all $\theta\in\Theta$ are thus a finite first-order real-algebraic formula. Compactness of the finite plan and decoder space gives an attained minimum. Its value and some minimizing variables are algebraic, by quantifier elimination and the existence of algebraic points in nonempty semialgebraic sets defined over the real algebraic numbers. This is a robust finite-horizon optimization, not necessarily an ordinary finite-list semidefinite program or a convex optimization in the decoder variables.

The same method computes the uniform reserved cost minimum $a_{\Theta,T}$. Complete its optimizer with the known fallback. With no model-grid error, its cost is at most
\[
 a_{\Theta,T}+(\kappa_0-c_*)
           \min\{1,(\kappa_0/c_*-1)/T\}.
\]
This explicit interval converges to $N_{\min}$ and can be made narrower than any rational tolerance, proving computability. If $N>N_{\min}$, enumeration of these finite uniform problems eventually certifies a completed fallback cost $\kappa<N$; it has finite labels and algebraic probabilities. The uniform risk problem just described and the completion-mixture proof then have gap at most
\[
 LD_T(N)+\frac{L(\kappa-c_*)D_T(N)}{N-\kappa+(\kappa-c_*)D_T(N)}.
\]
Choose a finite $T$ making this less than $\xi$, compute an algebraic optimizer, and mix its completion with the certified fallback. The mixing probability is algebraic and restores the original budget. Its one common decoder tuple has at most $m$ values. This proves every claimed termination and approximation assertion without supposing a finite grid represents a continuum exactly. Real-algebraic elimination may be extremely expensive, and equality at the limiting feasibility boundary is not decided by an unbounded strict-threshold search.
\end{proof}

\subsection{Two-sided robustness and boundaries of the theorem}
\begin{corollary}[Robustness of the optimal constrained risk]\label{cor:optimal-robustness}
Let $\Theta_1,\Theta_2$ be two compact families on the same interface. Suppose each member of either family is within $(\epsilon_I,\epsilon_g,\epsilon_F)$ of some member of the other. Suppose one common finite-label fallback costs at most $\kappa<N$ on their union, with the same cost and loss bounds. Then, for fixed $C$ or optimized width $m$,
\[
 |R_{\Theta_1}(N)-R_{\Theta_2}(N)|\le\Gamma_{F,T}
       \quad(T\ge1),
\]
where the right side is the explicit expression \eqref{eq:compact-budget-error}--\eqref{eq:compact-risk-error} evaluated at the pairwise error bounds; no finite cover is needed in this corollary. Hence it is at most the infimum of these bounds over $T$. In particular it tends to zero with the primitive errors at a fixed positive budget slack. The same-budget $O(\sqrt\epsilon)$ upper modulus follows when all three errors are at most $\epsilon$.
\end{corollary}
\begin{proof}
Take any feasible policy on $\Theta_1$ within $\xi$ of its optimal risk and truncate it at depth $T$. Its reserve and nonnegative loss bounds hold at every member of $\Theta_1$. For each target in $\Theta_2$, use its close source member in the occupation comparison. Complete and mix with the common fallback exactly as in Theorem~\ref{thm:compact-risk}. This gives a single policy feasible at $N$ on $\Theta_2$ with risk at most $R_{\Theta_1}(N)+\xi+\Gamma_{F,T}$. For optimized width, its one source dictionary is used on every target and by the fallback. Let $\xi\downarrow0$, then reverse the roles. This compares the optimal risks through feasible strategies and does not exchange a minimization and a maximization.
\end{proof}

\begin{example}[Risk need not be continuous at the feasibility boundary]\label{ex:boundary-discontinuity}
In physical dimension one, both commands $A,B$ always report a productive result. There is one retained decoder value. Let its prospective losses be zero after $A$ and one after $B$. At model $\theta=0$, both call costs are one; at $\theta=\epsilon>0$, they are respectively $1+\epsilon$ and one. There is no instrument or loss perturbation, only cost error $\epsilon_g=\epsilon$, and there is no dark state. At budget $N=1$, any positive probability of choosing $A$ or continuing after the first report costs strictly more than one in the perturbed model. Thus $R_{\{0\}}(1)=0$ while $R_{\{\epsilon\}}(1)=1$. Both instances remain feasible. Therefore the strict interior/slack condition cannot simply be deleted from a general continuity theorem. The example can be realized by command-dependent deterministic future audits and squared loss at the same decoder.
\end{example}

\begin{example}[Why the remaining-horizon cost term is real]\label{ex:cost-horizon}
In dimension one, command $A$ costs one and is immediately productive at the source; at the target it fails with probability $\epsilon\in(0,1/2]$ and is productive otherwise. Command $W$ costs one and is never productive. Use a prefix that first calls $A$, commits on success, and after failure makes $T-1$ calls to $W$ before leaving the cutoff open. Source prefix occupation is exactly one. Instrument discrepancy is $\epsilon$, while costs and task observables have not changed. The target's cost padded by $\kappa$ on the open leaf is exactly
\[
 1+\epsilon(T-1+\kappa),
\]
compared with source value one. Repeated $A$ is a genuine fallback of cost at most two. This attains the instrument term of \eqref{eq:prefix-cost-error} when $c^*=1$. It proves the need for a remaining-horizon factor for arbitrary tested prefixes; it does \emph{not} prove sharpness of the optimized-risk square-root modulus.
\end{example}

\begin{example}[Even affine primitive families cannot be certified at vertices]\label{ex:vertex-risk}
Let $\theta\in[0,1]$, physical dimension one, and one unit-cost command whose productive report is a Bernoulli($\theta$) bit. At commitment a fresh independent Bernoulli($\theta$) bit is the future audit. Fix $C=\{0,1\}$ with squared loss. Both the report instruments and the post-report loss observables are affine in $\theta$. Nevertheless, for every budget $N\ge1$,
\[
 R_{\{0,1\},C}(N)=0,\qquad R_{[0,1],C}(N)=\tfrac12.
\]
For the two endpoint models, committing the first report as prediction has zero loss. For the full interval this rule has risk $2\theta(1-\theta)\le1/2$ at cost one. Under $\theta=1/2$, the prospective audit is fair conditional on any finite stopping history and any selected decision in $C$; every proper policy therefore has loss exactly $1/2$. This gives the matching lower against arbitrary history and randomization. It is a continuum-uncertainty obstruction even with a uniform one-call fallback and no hidden dynamical state. A convex parameter set and affine primitive entries do not justify testing only its vertices: finite-path coefficients and adaptive terminal scores need not be affine in the fixed unknown parameter.
\end{example}

\begin{example}[Pointwise termination and model switching are different]\label{ex:compact-limits}
The noncompact scalar family with productive probability $p\in(0,1]$ terminates in every model but has no finite uniform cost. Its compact closure adds the dark model $p=0$. Compactness and a positive certified margin are therefore substantive.

Nor is the model allowed to change adversarially during a run. With one command on a qubit, let model 1 have failure Kraus operator $|0\rangle\langle1|$ and model 2 have $|1\rangle\langle0|$, with the complementary productive effects. Each fixed model terminates within two calls from every initial state. Alternating models 1 and 2 from state $|1\rangle$ causes failure forever. The compact finite union satisfies the present fixed-model theorem but not a switching-adversary theorem. Such uncertainty requires a different joint survival certificate.
\end{example}

This compact-family result is an intrinsic constrained-risk approximation on its stated finite physical interface, not an intrinsic characterization of when the geometric profile $\Q$ is exact. It supplies no uncharged model-identification phase and requires no normalized-filter contraction. Its cost threshold, risk certificate, finite output and finite observer implementation come from the same primitive data. The separate exact geometric, spectral, reset-communication and predictive-realization theorems retain their own hypotheses and research questions.

\section{The causal task-order transfer theorem}\label{sec:central}
Suppose that every terminating state-stopped competitor reaches at least $B_0\ge1$ live labels under $\theta_*$. Such a certificate must concern actual physical command compatibility, not printed action names. Define
\begin{equation}\label{eq:widths}
\begin{array}{c|cc}
 i&m_-(M)&m_+(M)\\\hline
 \cp&M&\lfloor M/D_0\rfloor\\
 \ev&M&\lfloor M/D_0\rfloor\\
 \st&M-B_0&\lfloor(M-B)/D_0\rfloor .
\end{array}
\end{equation}
The event upper also requires $M\ge B$. The trivial state-stopped lower $B_0=1$ is valid for a procedure initially waiting to make a call.

\begin{theorem}[Affine task-order transfer with a paid reference and common construction]\label{thm:transfer}
Assume all the prepared-kernel, full-report, task-order, effort, occupation, resource and common-construction hypotheses of Section~\ref{sec:data}. Fix $N,M$, put $b_-=N/\kappa_-$, $b_+=N/\kappa_+$, and suppose $b_-\ge1$ and $m_-(M)\ge1$. Set
\[
 t_N=(1-A_N)_+/u_-.
\]
If $t_N>0$, every one of the three risks satisfies
\begin{equation}\label{eq:central-lower}
 R_i^\Theta(N,M)\ge
 \bigl[l_-t_N\Q_{m_-(M)}(b_-/t_N)-E_N\bigr]_+.
\end{equation}
If $t_N=0$, the asserted lower is zero, with no evaluation of $\Q$ at infinity. An empty legal class still has risk $+\infty$.

Suppose also $m_+(M)\ge1$, $b_+\ge1$, and, for event stopping, $M\ge B$. If
\begin{equation}\label{eq:beta}
 \beta_N=\frac{l_+b_+}{1+b_+a_+}\ge1,
\end{equation}
one common Borel program with actual expected cost at most $N$ in every model attains
\begin{equation}\label{eq:central-upper}
 R_i^\Theta(N,M)\le
 \frac{u_+(1+b_+a_+)}{l_+}\Q_{m_+(M)}(\beta_N)+b_+e_+.
\end{equation}
It commits almost surely. If $\beta_N<1$, the always-accept first-attempt construction instead gives $R_i^\Theta(N,M)\le u_+\Q_{m_+(M)}(1)+e_+$ at cost at most $\kappa_+\le N$.

The lower is against every legal competing strategy and uses one actual witness. The upper is one parameter-independent strategy for all models. The two effort constants need not agree. Neither direction assumes that the comparison experiment belongs to $\Theta$ or that a normalized predictive update is uniformly contracting.
\end{theorem}

\begin{corollary}[Scale-uniform relative transfer]\label{cor:relative}
Under Theorem~\ref{thm:transfer}, suppose $A_N=E_N=a_+=e_+=0$. If $l_+b_+\ge1$, then
\begin{equation}\label{eq:relative}
 \frac{l_-}{u_-}\Q_{m_-(M)}(u_-b_-)
 \le R_i^\Theta(N,M)
 \le\frac{u_+}{l_+}\Q_{m_+(M)}(l_+b_+).
\end{equation}
The factors $l_\pm,u_\pm$ may be fixed nonunit constants, independent of budget and resolution. In particular no term proportional to $b_\pm$ times a fixed absolute calibration error is required in this relative-order class.
\end{corollary}
\begin{proof}
Substitute $t_N=1/u_-$ and $\beta_N=l_+b_+$ in the theorem.
\end{proof}

\begin{corollary}[Separate exact checkpoint and autonomous laws]\label{cor:exact}
Suppose the certificates are exact: $l_\pm=u_\pm=1$, all defects vanish, $\kappa_-=\kappa_+=\kappa$, and exports are the identity. Then, for $N\ge\kappa$,
\[
 R_\cp^\Theta(N,m)=\Q_m(N/\kappa).
\]
The constructed autonomous upper uses $B+m$ state-stopped labels, or $\max(B,m)$ event-stopped labels. If the mandatory live count is sharp, $B_0=B$, then
\begin{align}
 R_\st^\Theta(N,B+m)&=\Q_m(N/\kappa)\quad(m\ge1),\label{eq:typed-exact}\\
 R_\ev^\Theta(N,M)&=\Q_M(N/\kappa)\quad(M\ge B).\label{eq:event-exact}
\end{align}
The subtraction by $B$ belongs only to the state-stopped model. Nonidentity exports require the expanded dictionary in \eqref{eq:widths}, unless an optimal dictionary is invariant under all exports. These exact stopping distinctions are inherited from the complete companion and are reproved here to fix the resource semantics.
\end{corollary}
\begin{proof}
The same lower and upper profile coincide in each asserted interface. Identity exports give $D_0=1$. The packing and converse widths are proved in Lemma~\ref{lem:roles}. No lower on live states is used to forbid event-stopped reuse.
\end{proof}

The familiar absolute-defect theorem is also a specialization: take $l_\pm=u_\pm=1$ while retaining $A_N,E_N,a_+,e_+$. Thus relative and absolute calibration are composed in one statement. The chosen metric is the common scored task throughout; it is not changed between the converse and the construction.

\section{Proof of the central theorem}
\begin{lemma}[Profile attainment]\label{lem:attainment}
For every finite $m\ge1$ and $b\ge1$, the minimum in \eqref{eq:profile} is attained. The profile is nonincreasing in $m,b$ and left-continuous in $b>1$.
\end{lemma}
\begin{proof}
Write a dictionary as an ordered tuple in the compact set $H^m$, permitting repetitions. The feasible densities are a weak-star closed subset of the radius-$b$ ball of $L^\infty(\nu)$: positivity, the upper bound and the mass equality are each closed under pairing with $L^1(\nu)$. They are therefore weak-star compact. As the tuple varies, its squared-distance function varies uniformly on $H$, since all points belong to a bounded set. Along a convergent product net, split the objective difference into this uniform variation and the weak-star pairing against the fixed bounded measurable loss at the limit tuple. This proves continuity of the objective and attainment. Nets avoid an unwarranted sequential compactness assumption. Monotonicity follows from inclusion of feasible sets. For left continuity take an optimizer $h$ at $b>1$ and let $h_\lambda=(1-\lambda)h+\lambda$. It has unit mass and cap $b_\lambda=(1-\lambda)b+\lambda<b$. Its objective tends to the optimum as $\lambda\downarrow0$, while monotonicity bounds the optimum at $b_\lambda$ between the optimum at $b$ and this objective.
\end{proof}

\begin{lemma}[Actual-history comparison measure]\label{lem:comparison}
Fix any robustly admissible competitor of cost at most $N$, with its fixed terminal dictionary $C$. There are finite nonnegative measures $\xi_c$ on $X$, $c\in C$, such that $\xi=\sum_c\xi_c$ obeys
\begin{align}
 \xi&\le (\E_*A_\tau)\nu\le b_-\nu,\label{eq:measure-cap}\\
 1&\le u_-\xi(X)+A_N,\label{eq:measure-mass}\\
 \E_*\|Z-\widehat r\|^2
 &\ge l_-\sum_{c\in C}\int d_c\,d\xi_c-E_N.
 \label{eq:measure-risk}
\end{align}
The measure $\xi$ is not asserted to be the terminal law of any comparison execution.
\end{lemma}
\begin{proof}
Adjoin an independent uniform random variable to the current report on both sides. This is current-call randomness, not a seed retained between calls; it is sufficient to represent the competitor's randomized output and stopping decision. A finite-output kernel has such a representation by cumulative probabilities. At a fixed reached productive history $h$ and chosen command $a$, the post-report decision then gives an acceptance indicator $q_h(y)$ and an output $r_h(y)\in C$ on acceptance. Give $r_h$ an arbitrary fixed value on rejection. The notation $y$ now includes the current independent random variable. Integrating the original certificates over that variable preserves their constants. Dependence on the entire earlier history is permitted in these functions.

For each Borel $D\subset X$, set
\begin{equation}\label{eq:xi-definition}
 \xi_c(D)=\E_*\sum_{j\le\tau:\ j\ {
m productive}}
 \int\one_{\{\pi_{H_{j-1},a_j}(y)\in D\}}
 q_{H_{j-1}}(y)\one_{\{r_{H_{j-1}}(y)=c\}}
 P^0_{H_{j-1},a_j}(dy,H).
\end{equation}
The outer expectation is the actual pre-call occupation at $\theta_*$. All terms are jointly measurable and nonnegative. Tonelli defines the measures, whose total is finite because $\E_*A_\tau\le b_-$. Summing over $c$ and dropping the acceptance indicator gives \eqref{eq:measure-cap}: each comparison row has mark law $\nu$.

The analogous sum with the actual rows instead of $P^0$ has mass exactly one, because the actual procedure commits exactly once almost surely. Apply \eqref{eq:lower-mass} at each actual reached history, then sum. The occupation certificate yields \eqref{eq:measure-mass}. Apply \eqref{eq:lower-task} in the same way to the disjoint actual commitment events. Their expected losses sum to the actual risk. The comparison contribution is $d_{r_h(y)}(\pi_{h,a}y)$ by \eqref{eq:full-moments}; this identity conditions on the full current report and therefore does not overlook extra information seen by the competitor. The accumulated loss defect is at most $E_N$, giving \eqref{eq:measure-risk}. One may first truncate the number of calls. Monotone convergence for the nonnegative contributions and the finite defect bounds justify the limit.

No actual past is replaced by a comparison past. In particular the competitor need not terminate, or obey a cost bound, if hypothetically run under the comparison experiment. This distinction is why the argument applies when that experiment is not an operating model.
\end{proof}

\begin{remark}[Why the full report is required]
A constant mark can conceal a visible fair bit which also equals the future audit. A comparison using only that mark has $f=0,w=1$, but a two-value decoder reading the bit predicts it with zero loss. The hypotheses therefore condition on the full report seen by the competitor; matching a compressed mark marginal alone cannot establish the lower.
\end{remark}

\begin{lemma}[Live roles, output roles and physical compatibility]\label{lem:roles}
A program with $\ell$ live roles and $t$ output roles has a state-stopped implementation with $\ell+t$ labels and an event-stopped implementation with $\max(\ell,t)$ labels, preserving every path cost and complete stopped scored law. Conversely an event-stopped $M$-label program has at most $M$ decoder values. A state-stopped $M$-label program with at least $B_0$ reached live labels has at most $M-B_0$ terminal decoder values.

For a physical live-state certificate, consider finitely many mandatory pre-call contexts $I$ with no free context input. Suppose the commands legal at a history of context $i$ are exactly $\mathcal U_i$, and every terminating competitor visits every context with positive probability under the same witness. This positive probability may depend on the competitor; no uniform lower bound over competitors is required. Let $b_{\rm phys}$ be the least size of a cover of $I$ by subsets $J$ with $\bigcap_{i\in J}\mathcal U_i\ne\varnothing$. Then at least $b_{\rm phys}$ live labels are required. This certificate is invariant under duplicate names and bijective stateless recodings of the same physical commands.
\end{lemma}
\begin{proof}
For state stopping take disjoint copies of the live and terminal role sets and type them zero and one. For event stopping embed the live and terminal role sets separately into one alphabet of size $\max(\ell,t)$. On continuing transitions use the live embedding; on commitment emit the final event and use the terminal embedding. A command row is read only on a continuing branch, while the decoder is read only after commitment. Sharing their label therefore causes no conflict and uses no continuing mode bit. Couple the original and recoded coins. Every finite stopped prefix and every path cost agree, and almost-sure termination gives equality of the stopped scored laws. A fixed decoder has one value per retained label. In state stopping the global halt map excludes each reached live label from the output set, proving its subtraction. Event stopping has no such exclusion.

A live label reached in contexts $J$ must have one command distribution supported on all the corresponding $\mathcal U_i$. Since $J$ is finite, the union of the exceptional null command sets is null; hence their intersection supports a probability and is nonempty. The sets of contexts served by the labels cover $I$, proving the bound. Conversely a nonempty intersection contains a command and hence a Dirac command distribution. Replacing names without changing executed kernels, legal domains, costs and reports does not change these intersections. A translator which itself remembers a hidden phase is not such a recoding; its memory must be charged.
\end{proof}

\begin{proof}[Proof of Theorem~\ref{thm:transfer}]
Fix a competing program. Lemma~\ref{lem:roles} bounds its decoder width by $m_-(M)$. Let $s=\xi(X)$ in Lemma~\ref{lem:comparison}. If $t_N>0$, then \eqref{eq:measure-mass} gives $s\ge t_N$. Scale the measure to $\xi'=(t_N/s)\xi$. It has mass $t_N$ and remains bounded above by $b_-\nu$. Thus $h=(d\xi'/d\nu)/t_N$ has unit mass and cap $b_-/t_N$. Nearest-dictionary projection and nonnegativity imply
\[
 \sum_{c\in C}\int d_c\,d\xi_c
 \ge\int [w+\dist(f,C)^2]\,d\xi
 \ge t_N\Q_{m_-(M)}(b_-/t_N).
\]
Together with \eqref{eq:measure-risk} this proves the asserted lower for this competitor. Its worst-model risk dominates its risk at $\theta_*$. Taking the infimum over competitors gives \eqref{eq:central-lower}, without exchanging an infimum and a supremum. Nonnegative loss gives the zero truncation and the stated case $t_N=0$.

For the upper assume \eqref{eq:beta}. Choose an attained optimizer $(C,h)$ at width $m_+(M)$ and cap $\beta_N$. Let $r_C(x)$ be the first nearest point to $f(x)$ in a fixed ordering of $C$. After each attempt of $\mathcal G$, accept $x$ with probability $h(x)/\beta_N$, export $T_Jr_C(x)$ and commit; otherwise return to the start routing label. All these operations are Borel. The choice $J$ is used within the current attempt; any part retained thereafter is charged through the live role or exported decoder, not a free seed. The exports are assumed jointly measurable and take at most $D_0|C|$ distinct values, so their finite decoder and transitions are Borel as well.

Conditional on every surviving attempt history, the acceptance probability is at least
\begin{equation}\label{eq:hazard}
 p=\frac{l_+}{\beta_N}-a_+=\frac1{b_+}>0.
\end{equation}
Consequently the probability of reaching attempt $j+1$ is at most $(1-p)^j$, by induction on conditional survival probabilities. In particular commitment is almost sure and the expected number of attempts is at most $1/p=b_+$. Each begun attempt has conditional expected cost at most $\kappa_+$, not merely a cost conditional on a favorable outcome. Summing the predictable attempt costs gives actual expected cost at most $\kappa_+b_+=N$.

The conditional unnormalized loss contributed by an attempt is at most $u_+\Q_{m_+(M)}(\beta_N)/\beta_N+e_+$. Sum its contributions on the disjoint terminal events, using the same survival bound. The result is at most
\[
 b_+\left[\frac{u_+}{\beta_N}\Q_{m_+(M)}(\beta_N)+e_+\right],
\]
which is \eqref{eq:central-upper}. The program has $B$ live roles and at most $D_0m_+(M)$ terminal roles. Lemma~\ref{lem:roles} gives exactly the sufficient widths in \eqref{eq:widths}. A checkpoint counts only the terminal width. No physical renewal law was used: the conditional hazard controls survival even if the hidden state after rejection differs arbitrarily from its former state.

If $\beta_N<1$, commit after the first attempt, and use an optimizer at cap one. The loss certificate with $q=1$ gives $u_+\Q_{m_+(M)}(1)+e_+$, and the whole-attempt cost is at most $\kappa_+$. This proves all remaining assertions.
\end{proof}

\section{The exact task defect at a prescribed output width}\label{sec:width-defect}
The raw matrix order of Section~\ref{sec:order} is exact for localized affine-square tests. Ordinary squared prediction tests form a smaller family. Requiring the stronger order can therefore discard valid task comparisons, while taking a supremum over arbitrarily many readouts can overcharge a finite decoder. We identify the least additive defect for the actual task and its output width.

\subsection{Finite dictionaries and an attained optimal defect}
Let $Y,Z$ be standard Borel spaces, $D$ a nonempty compact metric decision space, and $\ell:Z\times D\to[0,L]$ Borel. Assume a uniform decision modulus
\begin{equation}\label{eq:loss-modulus}
 |\ell(z,r)-\ell(z,s)|\le\omega(d(r,s)),\qquad
 \lim_{\delta\downarrow0}\omega(\delta)=0,
\end{equation}
with $\omega$ nondecreasing. For finite positive measures $P,Q$ on $Y\times Z$, put $\lambda=P_Y+Q_Y$. The unnormalized conditional loss densities $g_P(y,r),g_Q(y,r)$ relative to $\lambda$ have jointly Borel versions continuous in $r$. Indeed disintegrate $P+Q$ over $\lambda$, multiply its conditional law by the Radon--Nikodym densities of $P,Q$ relative to $P+Q$, and integrate $\ell$. Off one completed null set,
\[
 |g_P(y,r)-g_P(y,s)|\le p(y)\omega(d(r,s)),\quad
 |g_Q(y,r)-g_Q(y,s)|\le q(y)\omega(d(r,s)),
\]
where $p=dP_Y/d\lambda$, $q=dQ_Y/d\lambda$ and $p+q=1$. Zero total mass causes no difficulty. Put $F(y,r)=g_P(y,r)-g_Q(y,r)$ and
\[
 \Delta_C^\ell(P\mid Q)=\int\max\{0,\max_{r\in C}F(y,r)\}\,d\lambda
\]
for a nonempty finite dictionary $C\subset D$.

\begin{theorem}[Width-resolved optimal task defect]\label{thm:width-defect}
For every integer $m\ge1$, the following equal quantities are attained:
\begin{align}
 \Delta_m^\ell(P\mid Q)&=\max_{(r_1,\ldots,r_m)\in D^m}
 \int\max\{0,F(y,r_1),\ldots,F(y,r_m)\}\,d\lambda,
 \label{eq:defect-envelope}\\
 &=\sup_{\substack{a:Y\to[0,1]\ {
m Borel}\\r:Y\to D\ {
m Borel},\ |r(Y)|\le m}}
       \int a(y)\ell(z,r(y))\,(P-Q)(dy,dz).
 \label{eq:defect-tests}
\end{align}
Thus this is the smallest additive error valid for every selected width-$m$ task. An optimum uses a deterministic report event and readout from one fixed dictionary; independent current randomization cannot improve it. The value is independent of the dominating measure and satisfies
\begin{equation}\label{eq:width-growth}
 \Delta_1\le\Delta_m\le m\Delta_1,\quad
 \Delta_{m+n}\le\Delta_m+\Delta_n,\quad
 \frac{\Delta_n}{n}\le\frac{\Delta_m}{m}\quad(1\le m\le n).
\end{equation}
Moreover
\begin{equation}\label{eq:infinite-defect}
 \Delta_\infty^\ell(P\mid Q):=\int[\sup_{r\in D}F(y,r)]_+\,d\lambda
       =\lim_{m\to\infty}\Delta_m^\ell(P\mid Q).
\end{equation}
If a $\delta$-net of $D$ has $n$ points, then
\begin{equation}\label{eq:defect-net}
 0\le\Delta_\infty-\Delta_n\le
 \omega(\delta)[P(Y\times Z)+Q(Y\times Z)].
\end{equation}
In particular $\Delta_1=0$ if and only if $\Delta_\infty=0$. The width hierarchy distinguishes positive defect magnitudes, not different zero-defect orders.
\end{theorem}
\begin{proof}
Fix an ordered dictionary. At each report choose the first index maximizing its positive loss-density difference, and accept exactly when that maximum is positive. These decisions are Borel and give the displayed envelope. They pointwise dominate every other selection and every randomized readout in that dictionary. Every finite-image rule uses some such dictionary, proving equality of the suprema.

Moving every tuple coordinate by at most $\delta$ changes the envelope by at most $(p+q)\omega(\delta)$ almost everywhere. Integration makes the objective continuous on compact $D^m$, with repetitions allowed. It attains its maximum, and the selected deterministic test attains the second supremum. Averaging over an independent current seed cannot increase that value. A seed choosing dictionaries whose union exceeds $m$ does not describe a width-$m$ decoder.

The test expression gives independence of the dominating measure. Inclusion gives monotonicity. A nonnegative maximum is bounded by the sum of the single-center positive parts, proving $\Delta_m\le m\Delta_1$. Partition a dictionary into two groups to obtain subadditivity. For $m\le n$, choose an optimal $n$-tuple and average over its $m$-element index subsets. A selected maximizing positive index is retained with probability $m/n$. Thus the average subdictionary envelope is at least $m/n$ times the original one. Integrate and take one subdictionary at least as good as its average. This proves the last inequality, including repeated entries.

Continuity makes a countable dense decision set determine the pointwise supremum. The positive envelopes of its finite initial segments increase to that supremum; monotone convergence proves the limit and measurability. A finite net gives the quantitative estimate by the same modulus calculation. Finally, $\Delta_1=0$ makes $F(y,r)\le0$ almost everywhere for each point in the countable dense set. Intersect these full-measure sets and use continuity to obtain the inequality for all decisions on one common set. Hence $\Delta_\infty=0$; the converse is immediate.
\end{proof}

For fixed $P_0,P_1$, the map $t\mapsto\Delta_m^\ell(tP_0\mid P_1)$ is convex, nondecreasing and $LP_0(Y\times Z)$-Lipschitz on $[0,\infty)$. It is a supremum of affine functions of $t$ with slopes between zero and that bound. The positive-variation and measurable-selection ingredients are classical. The decision-cardinality restriction should not be confused with Blackwell $k$-decision orders \cite{BR14}: the present localized fixed-loss order at zero defect already follows from width one.

\subsection{The exact squared-task order}
Take $D=H\subset\R^k$ compact convex and $\ell(z,r)=\|z-r\|^2$, with bounded audits. Relative to a common report measure, write the signed moment densities of $P_1-tP_0$ as
\[
 a=\frac{d(P_1-tP_0)_Y}{d\lambda},\quad
 b=\frac{d\int z(P_1-tP_0)(dy,dz)}{d\lambda},\quad
 c=\frac{d\int\|z\|^2(P_1-tP_0)(dy,dz)}{d\lambda}.
\]
The corresponding quadratic is $\psi_y(r)=a(y)\|r\|^2-2b(y)\cdot r+c(y)$.

\begin{theorem}[Necessary and sufficient squared-task certificate]\label{thm:quadratic-defect}
The least unrestricted-width defect is
\begin{equation}\label{eq:quadratic-defect}
 \Delta_\infty^{\rm sq}(tP_0\mid P_1)
     =\int[-\min_{r\in H}\psi_y(r)]_+\,d\lambda.
\end{equation}
Localized squared-task order with zero defect holds for all finite dictionaries if and only if $\min_H\psi_y\ge0$ on one common full-measure set. No inverse covariance, equal matrix support or positive eigenvalue floor is required. For $a>0$ the minimum is
\begin{equation}\label{eq:quadratic-projection}
 c-\|b\|^2/a+a\,\dist(b/a,H)^2;
\end{equation}
for $a=0$ it is $c-2\sup_H b\cdot r$. On the radius-$R$ closed Euclidean ball,
\begin{equation}\label{eq:quadratic-ball}
 \min_H\psi=\begin{cases}
 c-\|b\|^2/a,&a>0,\ \|b\|\le aR,\\
 c+aR^2-2R\|b\|,&\text{otherwise}.
 \end{cases}
\end{equation}
The second line includes the nonpositive-mass and zero-radius boundary cases whenever the first does not apply.
\end{theorem}
\begin{proof}
Expanding the square makes the directed loss-density difference for $(tP_0,P_1)$ equal to $-\psi$. The preceding theorem proves the integral and common-null-set statements. Completing the square proves the first formula; $a=0$ gives a linear optimization. On a ball with $a>0$, either the unconstrained minimizer $b/a$ lies inside or its radial projection minimizes. If $a\le0$, the least value at radius $s$ is $c+as^2-2\|b\|s$, nonincreasing for $0\le s\le R$. Its minimum is at $R$. These arguments include ties, $b=0$ and $R=0$.
\end{proof}

Report mass is a separate resource test: its least selected-event defect in $P_Y\le uQ_Y+e$ is $\int[p-uq]_+d\lambda$. It is not implied by one arbitrary task inequality. Equality of the entire task family has a sharper conclusion.

\begin{proposition}[Task equality and its exact moment content]\label{prop:task-rigidity}
Suppose the affine hull of $H$ has dimension at least one. Localized squared-task equality for $P,Q$ at every constant readout in $H$ is equivalent to equality of their report masses, their first-moment measures projected onto the direction space of that affine hull, and their full centered squared-norm measures $\int\|z-h_0\|^2P(dy,dz)$ for an origin $h_0$ in the affine hull. Full second-moment matrix measures need not agree. For a singleton carrier the mass must instead be tested separately.
\end{proposition}
\begin{proof}
Choose a countable dense set of readouts and intersect their full-measure equality sets. The common modulus extends equality to all of $H$. Translate by $h_0$ and restrict the quadratic difference to the direction space $V$ of the affine hull. A compact convex set with this affine hull has nonempty relative interior. A polynomial vanishing there vanishes on $V$: restrict successively to line segments in a relative open ball, where a one-variable polynomial with an interval of zeros is identically zero, and identify its coefficients. The quadratic coefficient is the scalar mass difference times the identity on $V$, so that difference is zero because $\dim V\ge1$. The linear coefficient is the projected centered first-moment difference, and the constant coefficient is the full centered squared-norm difference. They vanish. Conversely these three equalities make every expanded squared-task difference zero. When $H$ is a singleton, only one constant loss measure is tested, so a zero-loss audit does not identify mass.
\end{proof}

\begin{example}[Identical tasks with incompatible matrix supports]\label{ex:strict-task-order}
There is one report and the decision carrier is the unit ball in $\R^2$. Let $P_0$ audit the two signs of $e_1$ fairly and $P_1$ the two signs of $e_2$ fairly. Every ordinary squared decision has loss $1+\|r\|^2$ under both laws. Their affine-square moment matrices are respectively $\operatorname{diag}(1,1,0)$ and $\operatorname{diag}(1,0,1)$ in the coordinates $(1,z_1,z_2)$. No two-sided positive relative matrix order can compare their different kernels. Thus exact squared-task comparison is strictly broader than the sufficient positive-matrix criterion, without contradicting that criterion's exact statement for its larger test family.
\end{example}

\begin{theorem}[Sharp width amplification in every finite dimension]\label{thm:sharp-width}
For every integer $k\ge1$ there are probability report--audit laws with identical report masses and squared decisions in the unit ball of $\R^k$ such that
\begin{equation}\label{eq:sharp-width-law}
 \Delta_m^{\rm sq}(P_0\mid P_1)=\frac{\min(m,k)}{5k}
 \quad(m\ge1).
\end{equation}
Both the linear width growth and the discrepancy between fixed-width and unrestricted defects are sharp in this family.
\end{theorem}
\begin{proof}
Let the report $i$ be uniform on $\{1,\ldots,k\}$. Under $P_0$ the audit is zero. Under $P_1$, conditional on $i$, it is $e_i$ with probability $4/5$ and $-e_i$ with probability $1/5$. The conditional loss difference is $(6/5)r_i-1$. It is positive only when $r_i>5/6$. A unit-ball vector cannot meet this condition for two different coordinates, since $2(5/6)^2>1$. Its positive contribution in its one possible coordinate is at most $1/5$. A dictionary of $m$ vectors therefore contributes at most $\min(m,k)/(5k)$. Choosing that many distinct basis vectors and accepting their corresponding reports attains the bound. The construction has bounded audits, exact report-mass equality, and no asymptotic dimension assumption.
\end{proof}

\subsection{The same dictionary across actual histories}
Let $u$ range over a standard Borel context space with finite positive occupation measure $\mu$. Let $P_u,Q_u$ be finite kernels with integrable total masses; assume the same bounded decision-continuous loss. Form the finite measures $\mathbf P=\mu(du)P_u(dy,dz)$ and $\mathbf Q=\mu(du)Q_u(dy,dz)$, where the full report is $(u,y)$.

\begin{proposition}[Exact fixed-dictionary occupation defect]\label{prop:occupation-defect}
One has
\begin{equation}\label{eq:occupation-defect}
 \Delta_m^\ell(\mathbf P\mid\mathbf Q)
   =\max_{C\in D^m}\int\Delta_C^\ell(P_u\mid Q_u)\,\mu(du)
   \le\int\Delta_m^\ell(P_u\mid Q_u)\,\mu(du).
\end{equation}
The maximum uses a single dictionary for every context. The inequality can be strict by an arbitrarily large factor as the number of contexts grows. Neither side asserts that a prescribed occupation measure is causally realizable.
\end{proposition}
\begin{proof}
For a fixed dictionary, take the positive envelope of its unnormalized loss densities on the product context/report space. Disintegration and integration give exactly the integral of the rowwise fixed-dictionary envelopes. This can also be checked without selecting contextwise densities: a finite partition of reports and rational acceptance/readout choices approximate every bounded selected task in the finite measure $P_u+Q_u$. A countable algebra generating the standard Borel report space, together with a countable dense decision set, suffices for these approximations. Hence the rowwise suprema are measurable. The uniform decision modulus and integrable total masses make the integrated fixed-dictionary envelope continuous on compact $D^m$, so it attains its maximum. Applying the width-defect theorem to the product measures yields the equality. Moving the maximum inside the integral gives the inequality.

For strictness use the preceding sharp example with context $u=i$ and no additional report. Every row has one-center defect $1/5$, whereas a single common width-$m$ dictionary over uniform contexts has defect $\min(m,k)/(5k)$. Thus even individually exact rowwise errors may overcharge a common finite decoder. In a stopped experiment $\mu$ must be the actual pre-call occupation of the competitor; this algebraic identity does not optimize or fabricate that occupation.
\end{proof}

\section{Width-sensitive transfer for a general acquired loss}\label{sec:loss-transfer}
The exact defects just constructed can be inserted into the paid causal proof without enlarging the observer's decoder. We state the resulting theorem for bounded nonnegative decision-continuous losses. Squared prediction is its principal example, not an implicit restriction on its quantifiers.

Let $D$ be a nonempty compact metric decision space and let $\ell:Z\times D\to[0,L]$ satisfy \eqref{eq:loss-modulus}. A genuinely prepared comparison law $K_0(dx,dz)$ has report law $\nu$ and conditional task functions $g_r(x)=\E_0[\ell(Z,r)\mid x]$. Disintegration supplies jointly Borel versions with the same decision modulus. Define
\begin{equation}\label{eq:general-profile}
 \Q_m^\ell(b)=\min_{\substack{(r_1,\ldots,r_m)\in D^m\\0\le h\le b,\ \int h\,d\nu=1}}
       \int h(x)\min_j g_{r_j}(x)\,\nu(dx),\qquad b\ge1.
\end{equation}
Repetition in the tuple permits fewer than $m$ outputs. Values at infeasible width or cap are $+\infty$.

Fix one actual lower model, the same actual-history occupation and effort convention as in Section~\ref{sec:data}, and the same lower report-mass comparison with constants $u_-,a_-$. The full-report comparison rows must satisfy
\begin{equation}\label{eq:full-loss-factor}
 \E^0[\ell(Z,r)\mid Y=y]=g_r(\pi y)\qquad(r\in D)
\end{equation}
on one common full-measure set. It is enough to verify this on a countable dense decision set; the uniform modulus extends it to all decisions. This excludes an unaccounted informative component of the report.

At a reached lower context $u=(h,a)$ let
\begin{equation}\label{eq:row-width-error}
 e_{-,m}(u)=\Delta_m^\ell(l_-P^0_u\mid P_u),\qquad 0<l_-\le1.
\end{equation}
These optimal defects may be replaced by measurable certified upper bounds. Suppose every robustly feasible budget-$N$ competitor has actual occupation $\mu$ for which
\begin{equation}\label{eq:global-width-error}
 \Delta_m^\ell(l_-\mu(du)P^0_u\mid\mu(du)P_u)\le E_{N,m}.
\end{equation}
The common terminal dictionary is fixed across these contexts. Proposition~\ref{prop:occupation-defect} shows that it suffices to bound $\int e_{-,m}\,d\mu$; the displayed global certificate can be sharper. A uniform row bound $e_{-,m}\le\bar e_m$ gives $E_{N,m}\le(N/\kappa_-)\bar e_m$. The killed-potential lemma also applies to nonuniform certified row errors. The occupation is that of each actual competing program. It is not a freely chosen measure or the occupation of a fictitious comparison execution.

For the upper, retain one common acquisition program $\mathcal G$, its conditional whole-attempt cost $\kappa_+$, live count $B$, actual rooted-continuation validity, and export multiplicity $D_0$. Its report acceptance obeys the same lower mass certificate with $l_+,a_+$. For every virtual dictionary of width at most $m$, assume
\begin{equation}\label{eq:upper-general-task}
 \E[q(X)\ell(Z,T_Jr(X))]
 \le u_+\int q(x)g_{r(x)}(x)\nu(dx)+e_{+,m},
 \qquad u_+\ge1,
\end{equation}
conditionally on every surviving history in every model. This is one program and one virtual dictionary, not modelwise optimized decisions. The exact defect theorem computes the least possible $e_{+,m}$ whenever the compared effective loss family has the stated common decision modulus. For example, uniformly Lipschitz exports on a compact carrier preserve such a modulus. The two measures can be placed on a disjoint union of their audit carriers with the corresponding loss definitions; equality of hidden audit coordinates is not needed. Otherwise \eqref{eq:upper-general-task} itself is the required task certificate.

\begin{theorem}[Paid transfer at the actual decoder width]\label{thm:width-transfer}
Assume the prepared-kernel, commitment, cost, mass, effort, occupation and resource conditions just stated. Use exactly the lower and upper widths $m_-(M),m_+(M)$ of \eqref{eq:widths}, with the event upper requiring $M\ge B$. Put $b_\pm=N/\kappa_\pm$, $t=(1-A_N)_+/u_-$ and suppose $b_-\ge1$ and $m_-(M)\ge1$. If $t>0$, then
\begin{equation}\label{eq:width-transfer-lower}
 R_i^{\Theta,\ell}(N,M)\ge
 \left[l_-t\Q_{m_-(M)}^\ell(b_-/t)-E_{N,m_-(M)}\right]_+.
\end{equation}
For $t=0$ the lower is zero, without evaluating a profile at infinity. Empty legal classes retain value $+\infty$.

If $b_+\ge1$, $m_+(M)\ge1$ and
\[
 \beta=\frac{l_+b_+}{1+b_+a_+}\ge1,
\]
there is one common Borel finite-role construction with actual expected cost at most $N$ and
\begin{equation}\label{eq:width-transfer-upper}
 R_i^{\Theta,\ell}(N,M)\le
 \frac{u_+(1+b_+a_+)}{l_+}\Q_{m_+(M)}^\ell(\beta)
             +b_+e_{+,m_+(M)}.
\end{equation}
It commits almost surely in every model. If $\beta<1$, the first-attempt upper is $u_+\Q_{m_+(M)}^\ell(1)+e_{+,m_+(M)}$, at cost at most $\kappa_+$. Exact matched certificates with identity exports give the checkpoint and event/state-stopped equalities of Corollary~\ref{cor:exact}, with $\Q^\ell$ replacing $\Q$.

Optimal row defects are necessary and sufficient for their declared selected-task inequalities. Sharpness of those defects does not, by itself, assert sharpness of the entire stopped risk bound when other mass, cost or implementation certificates are not matched.
\end{theorem}
\begin{proof}
First the profile minimum exists. The density set is weak-star compact in $L^\infty(\nu)$, and $D^m$ is compact. Moving its decision tuple by at most $\delta$ changes its pointwise minimum loss by at most $\omega(\delta)$. Split an objective difference into this uniform change and weak-star pairing with the bounded fixed limiting loss. The objective is therefore continuous on the product, so the minimum is attained. This argument uses product nets if necessary, not unjustified sequential compactness. A first nearest conditional-loss index is Borel.

Fix any competing observer. Its terminal values form one dictionary $C$ with at most $m_-(M)$ values, by the exact stopping-role lemma. Adjoin independent current coins to the full report, and write the current acceptance and readout as functions of the actual context and report. In the lower comparison, summing their contributions against actual pre-call occupation gives exactly the two product measures in \eqref{eq:global-width-error}. The defining defect inequality yields
\[
 \E_*\ell(Z,\widehat r)
 \ge l_-\sum_{c\in C}\int g_c\,d\xi_c-E_{N,m_-(M)},
\]
where the accepted comparison measures $\xi_c$ are constructed as in \eqref{eq:xi-definition}. Condition \eqref{eq:full-loss-factor} is what identifies their conditional tasks after selection. The actual model commits exactly once; the same report-mass proof gives $\xi=\sum_c\xi_c\le b_-\nu$ and $\xi(X)\ge t$. All these measures are finite by the actual effort inequality.

For $s=\xi(X)\ge t>0$, scale to $\xi'=(t/s)\xi$. Nonnegative loss and pointwise minimization over $C$ give
\[
 \sum_c\int g_c\,d\xi_c
 \ge\int\min_{c\in C}g_c\,d\xi'
 \ge t\Q_{m_-(M)}^\ell(b_-/t).
\]
The normalized density $t^{-1}d\xi'/d\nu$ has unit mass and the asserted cap. Taking the worst-model risk and then the infimum proves the lower without exchanging infimum and supremum. Nonnegativity and the empty-class convention give its boundary cases.

For the upper take an attained optimal tuple and density at width $m_+(M)$ and cap $\beta$. In each attempt accept with probability $h(X)/\beta$, choose the first minimizing conditional-loss dictionary value, and apply its declared export. The conditional acceptance probability is at least $l_+/\beta-a_+=1/b_+$. Induction bounds survival through $n$ attempts by $(1-1/b_+)^n$; no physical renewal or independent-attempt law is assumed. Thus commitment is almost sure and the expected number of begun attempts is at most $b_+$. Summing the conditional whole-attempt cost gives actual expected cost at most $N$.

The conditional unnormalized terminal loss of an attempt is at most $u_+\Q_{m_+(M)}^\ell(\beta)/\beta+e_{+,m_+(M)}$. Summing on the disjoint actual terminal events and the same survival bound proves the displayed upper. Its actual decoder has at most $D_0m_+(M)$ values and its live roles are those of $\mathcal G$. The role-packing lemma gives the claimed event or state counts, with no extra retained mark or free seed. When $\beta<1$, use $q=1$ in the first attempt instead. Finally, in the exact matched case the lower and upper profiles agree at the stated respective widths. All assertions follow.
\end{proof}

\begin{corollary}[Acquired Bregman geometry]\label{cor:bregman-profile}
Let $H$ be compact convex in $\R^k$, and let $\Phi$ be continuously differentiable and convex on a neighborhood of $H$. For the Bregman loss
\[
 D_\Phi(z,r)=\Phi(z)-\Phi(r)-\nabla\Phi(r)\cdot(z-r),
\]
the general theorem applies, with
\begin{equation}\label{eq:Bregman-profile}
 g_r(x)=w_\Phi(x)+D_\Phi(f(x),r),\quad
 f(x)=\E_0[Z\mid x],\quad
 w_\Phi(x)=\E_0[\Phi(Z)\mid x]-\Phi(f(x))\ge0.
\end{equation}
If additionally $\Phi$ is twice continuously differentiable and $aI\preccurlyeq\nabla^2\Phi\preccurlyeq AI$ on $H$, with $0<a\le A$, then
\[
 \frac a2\Q_m^{\rm sq}(b)\le\Q_m^{D_\Phi}(b)
                  \le\frac A2\Q_m^{\rm sq}(b).
\]
Without the curvature lower bound the squared-loss exponents are not asserted for every Bregman loss.
\end{corollary}
\begin{proof}
Expand the conditional expected loss and insert $f$. The affine term cancels, leaving the displayed identity; Jensen's inequality gives $w_\Phi\ge0$. This conditional-expectation property has a classical Bregman-prediction antecedent \cite{Banerjee05}. Compactness and continuity give a bounded nonnegative loss and a common decision modulus, so the general theorem applies. Under the Hessian bounds, integration of the second derivative along the segment between $z$ and $r$ gives $(a/2)\|z-r\|^2\le D_\Phi(z,r)\le(A/2)\|z-r\|^2$. Integrate for every feasible dictionary and selected density and then take the infima. The profile comparison follows with the same acquisition and width constraints.
\end{proof}

\section{The exact zero-risk resource spectrum}\label{sec:zero-spectrum}
The mass of one favorable atom gives a sufficient resource threshold, not in general a necessary one. Other zero-loss atoms may be available to the same decoder. We determine the complete finite-width threshold, for the general acquired loss first and then for squared prediction. This statement is about a genuinely acquired probability law; converting it to a physical budget still requires the exact transfer hypotheses.

Use the compact decision space and nonnegative conditional losses $g_r(x)$ of Section~\ref{sec:loss-transfer}. For a dictionary $C$, put $d_C(x)=\min_{r\in C}g_r(x)$ and define
\[
 z_m=\max_{C\subset D,\,1\le |C|\le m}\nu\{d_C=0\}.
\]
\begin{theorem}[Exact finite-width zero-loss criterion]\label{thm:zero-general}
The maximum defining $z_m$ is attained. For every finite $m\ge1$ and finite $b\ge1$,
\begin{equation}\label{eq:zero-general}
 \Q_m^\ell(b)=0\quad\Longleftrightarrow\quad bz_m\ge1.
\end{equation}
When the inequality holds, an optimal procedure accepts only reports in a zero-loss set of an attaining dictionary. When it fails, the attained profile is strictly positive. Neither implication requires positive residual variance, a dimension assumption, a density, or an absence of atoms.
\end{theorem}
\begin{proof}
Represent a dictionary by a tuple in compact $D^m$, allowing repetitions. Conditional losses are continuous in each decision on one common full-measure set. For a convergent sequence of tuples, the nonnegative minima converge pointwise. The indicator of the closed condition $d_C(x)=0$ is upper semicontinuous in the tuple. Reverse Fatou for bounded functions therefore makes its integral upper semicontinuous; compactness gives a maximizing tuple.

If $v=\nu\{d_C=0\}\ge1/b$, the density $h=v^{-1}\one_{\{d_C=0\}}$ has unit mass and cap at most $b$, with objective zero. Conversely, take an attained optimal tuple and density for a zero profile, whose existence was proved in Theorem~\ref{thm:width-transfer}. Nonnegativity implies $h(x)d_C(x)=0$ almost everywhere. Thus its unit mass is supported on $\{d_C=0\}$ and is at most $b\nu\{d_C=0\}\le bz_m$. If $bz_m<1$, no optimizer has zero value; nonnegativity and attainment then make the value strictly positive.
\end{proof}

For squared prediction, form the finite measure on $H$
\[
 \mu_0(A)=\nu\{x:w(x)=0,\ f(x)\in A\}.
\]
List its positive atomic masses in nonincreasing order as $a_1\ge a_2\ge\cdots$, padding with zeros if there are finitely many, and set $a_j=0$ for every $j$ if there are none. The total mass is at most one. Repeated values in this list belong to distinct atom locations.

\begin{theorem}[Atomic zero-risk spectrum and exact resource thresholds]\label{thm:zero-spectrum}
For every $m\ge1$,
\begin{equation}\label{eq:atomic-zero-spectrum}
 z_m=\sum_{j=1}^m a_j,\qquad
 \Q_m(b)=0\quad\Longleftrightarrow\quad
 b\sum_{j=1}^m a_j\ge1.
\end{equation}
The nonatomic part of $\mu_0$ contributes nothing at a fixed finite width, even when $w=0$ everywhere. All finite-width zero-risk acquisition thresholds are therefore determined by the ordered atomic masses of the zero-residual predictive law, not by its matrix rank or its ambient dimension.

Under the exact matched certificates with identity exports, the checkpoint threshold at width $m$ is $N=\kappa/z_m$ when $z_m>0$, and no finite budget yields zero risk when $z_m=0$. For the event-stopped equality, replace $m$ by the total width $M\ge B$. For the state-stopped equality with a sharp mandatory live count, use total labels $B+m$. These conclusions are necessary and sufficient within those declared exact-transfer classes, not for arbitrary experimental interfaces.
\end{theorem}
\begin{proof}
For squared loss, $d_C=w+\dist(f,C)^2$ is zero exactly when $w=0$ and $f\in C$. Hence $\nu\{d_C=0\}=\mu_0(C)$. A finite set receives mass only from atoms. There are at most countably many positive atoms: for each positive integer $n$, at most $n$ atoms can have mass exceeding $1/n$. A nonempty collection of remaining positive masses has a largest member. Indeed if its supremum is $a>0$, only finitely many members can exceed $a/2$, and one of this finite set attains the supremum. Successive removal therefore defines the displayed ordering. Choosing the locations of its first $m$ positive masses maximizes $\mu_0(C)$ among all dictionaries of size at most $m$. Missing locations may be filled by arbitrary repeated points in $H$. This proves the first identity; Theorem~\ref{thm:zero-general} proves the equivalence. Finally substitute $b=N/\kappa$ into the respective exact transfer equalities, keeping their different memory counts.
\end{proof}

\begin{example}[Why one atom does not determine the necessary threshold]\label{ex:two-zero-atoms}
Let the acquired marked law put mass $1/4$ at $(0,0)$ and mass $3/4$ at $(1,0)$, in a carrier containing $[0,1]$. Treating only the first atom would suggest budget cap four. In fact $z_1=3/4$ and $z_2=1$. Thus $\Q_1(b)=0$ exactly for $b\ge4/3$, while $\Q_2(b)=0$ already at $b=1$. The remainder has zero residual variance, which is precisely why the single-atom necessary threshold does not apply. Proposition~\ref{prop:atom} now explicitly states the additional hypothesis needed for its one-atom necessary threshold; the remainder of that proposition is unchanged.
\end{example}

\subsection{An exact variational form, including ties}
The following fixed-dictionary formula records the classical selection mechanism explicitly. It is a population trimming or bounded-density variational principle, not a new invention of a trimming functional \cite{CGM97,GG99}. Its role here is to keep boundary mass and the physical acquisition cap in the same proof.
\begin{proposition}[Selected-mass dual and optimal selection]\label{prop:selection-dual}
For $0\le d\le L$ measurable on a probability space and $b\ge1$,
\begin{equation}\label{eq:selection-dual}
 \min_{0\le h\le b,\ \int h\,d\nu=1}\int hd\,d\nu
 =\max_{0\le t\le L}\left(t-b\int(t-d)_+\,d\nu\right).
\end{equation}
There is a threshold $t$ satisfying $\nu(d<t)\le1/b\le\nu(d\le t)$. An optimizer has density $b$ on $\{d<t\}$, zero on $\{d>t\}$ and a constant in $[0,b]$ on the tie set, chosen to give unit mass. More generally these saturation rules characterize every optimum for an attaining threshold, up to null sets. Applying this formula with $d=d_C$ and then minimizing over $C\in D^m$ gives the exact acquired profile; no interchange of the dictionary minimum and threshold maximum is asserted.
\end{proposition}
\begin{proof}
For every feasible $h$ and every $t$, pointwise minimization of $h(d-t)$ over $[0,b]$ gives $h(d-t)\ge-b(t-d)_+$. Integrating and using unit mass proves weak duality. Take a lower $1/b$ quantile of $d$ in $[0,L]$; elementary continuity of probability measures from below and above gives the two quantile inequalities, including $b=1$ and atoms at the endpoints. The prescribed density exists because the missing mass lies between zero and $b\nu(d=t)$. It attains equality in the pointwise inequality, proving equality and both attainments. The nonnegative difference in that inequality is $h(d-t)$ on $d>t$ and $(b-h)(t-d)$ on $d<t$. Its integral vanishes exactly under the stated saturation rules. On $d=t$ it is zero for every allowed density. This also proves the claim about all optimizers and ties.
\end{proof}

\section{An intrinsic hierarchy for paid causal synthesis}\label{sec:synthesis}
A comparison inequality does not create the acquisition program used by its upper bound. We give an intrinsic criterion over all classical feedback programs on a finite-branching causal interface. It applies to unknown-model classical experiments and finite-dimensional quantum instruments, with one program shared by every model and initial state. Exact unrestricted-history synthesis and exact fixed-memory synthesis remain different conclusions.

\subsection{Primitive data and the controller class}
A visible history records executed commands and reports. At each open decision node $h$, the legal next-command set $\mathcal A(h)$ is finite and nonempty. After a productive report there may also be a finite set $\mathcal O(h)$ of legal terminal routes. There is no terminal route at the root or after a nonproductive report. Each command has finitely many reports, so the tree is locally finite and countable. Legal domains are common to the admitted models and are part of the interface. A physical phase determined by visible history is not an additional free input to a finite-memory observer.

The unknown model $\theta$ belongs to an arbitrary nonempty set $\Theta$, not necessarily finite, compact or equipped with a prior. At each history the hidden quantum state lies in a specified finite-dimensional matrix algebra. Command $a$ has completely positive maps $\mathcal I_{\theta,h,a,y}$ with
\[
 \sum_y\mathcal I_{\theta,h,a,y}^*(I)=I.
\]
Dimensions may depend on the model and finite history. The initial identity is $I_\theta$. Nonnegative classical transition matrices are the diagonal special case. There is no positive-likelihood or eigenvalue floor.

Costs are nonnegative. Their conditional expected-cost observables satisfy the uniform floor
\begin{equation}\label{eq:syn-cost-floor}
 G_{\theta,h,a}\succeq\underline c I,\qquad \underline c>0.
\end{equation}
A pathwise floor implies this assumption, but only the displayed conditional floor is needed here. The cost $C_\tau$ is the core cost through commitment. A separately executed preparation or audit is charged in addition or represented as a compulsory call; it is not silently included without its physical operation.

Fix finitely many bounded terminal scores $j=1,\ldots,J$. Route $o$ at $h$ has a Hermitian score observable $F_{\theta,h,o,j}$ on the current state, with
\begin{equation}\label{eq:syn-score-bounds}
 \alpha_jI\preccurlyeq F_{\theta,h,o,j}\preccurlyeq\beta_jI,
 \qquad -\infty<\alpha_j\le\beta_j<\infty.
\end{equation}
The route is chosen before observing its future audit. For example, $F$ is the adjoint of the declared audit instrument applied to a bounded score of its outcome and the chosen route. These are primitive unnormalized observables, not assumed posteriors. Prescribed Hermitian targets on the initial state are denoted by $A_{\theta,j}$; constant targets are scalar multiples of $I_\theta$.

A controller uses full visible history and independent private randomness, but neither the unknown model nor an unknown initial density matrix. It must commit exactly once almost surely and satisfy
\begin{equation}\label{eq:syn-goal}
 \sup_{\theta,\rho}\E_{\theta,\rho}C_\tau\le N,\qquad
 \E_{\theta,\rho}\Phi_j=\tr(\rho A_{\theta,j})
 \quad(\theta\in\Theta,\ \rho\succeq0,\ \tr\rho=1).
\end{equation}
This is classical feedback over the specified quantum experiment, not an uncharged coherent controller memory. Uniformity over initial states is distinct from uniformity over rooted continuations of a history-dependent experiment; both are required before a program may be repeated after rejection. Restricted preparation faces and deadlocks are treated explicitly in the current supplement.

\subsection{One realization plan, rather than modelwise policies}
At a horizon $T\ge1$, include decisions following the first $T$ reports. Continuing at depth $T$ leaves an open branch; its next physical call has not occurred. The weight $\gamma(h)$ is the product of controller choice probabilities along $h$, excluding physical report probabilities. Write $\gamma(h,d)$ after a command or terminal decision. The common sequence-form polytope $\mathcal P_T$ is given by
\begin{equation}\label{eq:realization-plan}
 \gamma(\varnothing)=1,\quad \gamma(h,d)\ge0,\quad
 \sum_d\gamma(h,d)=\gamma(h),\quad
 \gamma(h,a,y)=\gamma(h,a)
\end{equation}
for every present report child. All variables lie in $[0,1]$, and the same variables are used for every model. Sequence form itself is classical \cite{Stengel96}; its role is to preserve the common-policy quantifier.

Let $\mathcal I_{\theta,h}$ be the composition of the primitive maps along $h$, without controller weights. Define affine functions on this polytope by
\begin{align}
 \mathsf B_{\theta,T}&=\sum_{\substack{|h|\le T\\o\in\mathcal O(h)}}
       \gamma(h,o)\mathcal I_{\theta,h}^*(I),&
 \mathsf S_{\theta,T}&=I_\theta-\mathsf B_{\theta,T},\label{eq:syn-mass}\\
 \mathsf V_{\theta,j,T}&=\sum_{\substack{|h|\le T\\o\in\mathcal O(h)}}
       \gamma(h,o)\mathcal I_{\theta,h}^*(F_{\theta,h,o,j}),\label{eq:syn-value}\\
 \mathsf C_{\theta,T}&=\sum_{\substack{|h|<T\\a\in\mathcal A(h)}}
       \gamma(h,a)\mathcal I_{\theta,h}^*(G_{\theta,h,a}).\label{eq:syn-cost}
\end{align}
Every coefficient is a raw finite-path matrix. No report probability is divided out.

\begin{lemma}[One plan for all models]\label{lem:common-plan}
The polytope is nonempty and compact. Its elements are exactly the finite-prefix laws of common randomized feedback controllers, with expectations given by traces against the displayed operators, simultaneously for every model and initial state. Arbitrary private random memory does not enlarge these laws relative to visible-history policies using fresh randomness. Moreover $\mathsf B_{\theta,T},\mathsf S_{\theta,T}\succeq0$ and their sum is $I_\theta$.
\end{lemma}
\begin{proof}
A behavioral policy gives the product weights. Conversely, at a node of positive weight choose $d$ with probability $\gamma(h,d)/\gamma(h)$; at a zero-weight node choose any legal row. Induction recovers every specified weight. The finite linear constraints form a closed bounded set, and any choice of legal rows proves nonemptiness.

To include arbitrary private memory, fix a visible action--report path and integrate over all private random tapes producing those controller choices when supplied that report path. This probability is independent of the physical model and initial state. The raw path instrument is the same for every compatible tape, because the executed commands and reports have been fixed. Thus it factors from the tape integral, and the integrated choice weights obey the same equations. The ratio construction reproduces the laws in all models at once. It uses full visible-history storage; it does not make a tape into free finite memory.

Multiply the raw path effects by these weights. Summing report effects uses normalization of each instrument, and summing controller choices uses the plan equation. Induction shows that terminated and open effects sum to the incoming effect, starting from $I_\theta$. Every such effect is positive. The cost and score formulas follow by summing their actual finite-prefix contributions.
\end{proof}

\subsection{The necessary cost reserve and terminal-score intervals}
For a finite $F\subset\Theta$, let $\mathcal H_T(F,N)$ be the feasibility problem in $\gamma\in\mathcal P_T$ defined by
\begin{align}
 \mathsf C_{\theta,T}+\underline c\mathsf S_{\theta,T}
       &\preccurlyeq NI_\theta,\label{eq:hier-cost}\\
 \alpha_j\mathsf S_{\theta,T}
       &\preccurlyeq A_{\theta,j}-\mathsf V_{\theta,j,T}
        \preccurlyeq\beta_j\mathsf S_{\theta,T}
       \quad(\theta\in F,\ 1\le j\le J).\label{eq:hier-tail}
\end{align}
These are linear inequalities in the classical case and linear matrix inequalities in the quantum case. They are not implemented by pretending that open paths terminate at the cutoff. Those paths owe at least one further paid call, and their remaining bounded scores lie in the displayed interval.

\begin{lemma}[Restriction and the sharp budget tail]\label{lem:synthesis-tail}
A feasible plan at level $T+1$ restricts to a feasible plan at level $T$. Every feasible level with $N\ge\underline c$ satisfies
\begin{equation}\label{eq:sharp-tail}
 0\preccurlyeq\mathsf S_{\theta,T}\preccurlyeq
 \delta_T I_\theta,\qquad
 \delta_T=\min\{1,(N/\underline c-1)/T\}.
\end{equation}
If $N<\underline c$, there is no feasible proper program.
\end{lemma}
\begin{proof}
Write $D_T=\mathsf B_{T+1}-\mathsf B_T\succeq0$ for the new terminal effect. Normalization gives $\mathsf S_T=D_T+\mathsf S_{T+1}$. New score increments lie between $\alpha_jD_T$ and $\beta_jD_T$. Adding them to the next-level remainder intervals proves the current-level intervals. Every path still open after depth $T$ executes another call before the next decision. Therefore $\mathsf C_{T+1}-\mathsf C_T\succeq\underline c\mathsf S_T$. Together with the next-level cost constraint this implies the current one.

Let $S_t$ denote the open effect after $t$ calls, so $S_0=I$ and $S_t\succeq S_T$ for $t<T$. The conditional cost floor gives $\mathsf C_T\succeq\underline c\sum_{t=0}^{T-1}S_t\succeq\underline c[I+(T-1)S_T]$. Adding the reserve gives $\underline c[I+TS_T]\preccurlyeq NI$. Positivity and $S_T\preccurlyeq I$ prove the claimed bound. Every proper program must execute its first paid call, proving the infeasible boundary below the floor.
\end{proof}

\begin{theorem}[Finite-obstruction characterization of paid causal synthesis]\label{thm:causal-hierarchy}
Under the stated finite-branching, positive-cost and bounded-score hypotheses, the following are equivalent:

\textup{(i)} One common proper visible-history program satisfies \eqref{eq:syn-goal} for every model and every initial state.

\textup{(ii)} $\mathcal H_T(F,N)$ is feasible for every finite model set $F\subset\Theta$ and every integer $T\ge1$.

Consequently nonexistence is witnessed by a finite horizon and a finite set of operating models. Feasibility at arbitrarily many tested levels alone is not a finite certificate for the infinite condition. No fixed bound on controller memory is asserted for the limiting policy.
\end{theorem}
\begin{proof}
For necessity take the prefix of the actual program. Every open path must make at least one additional call before commitment, so its conditional future expected cost is at least $\underline c$. Nonnegative conditional summation proves the reserved cost inequality. Scores on the disjoint future terminal events lie between $\alpha_j$ and $\beta_j$ times their total open mass. Almost-sure commitment makes that mass exactly $\mathsf S_T$, proving the score intervals. These same prefix weights work in every model.

For sufficiency consider the compact product of all nodewise probability simplexes on the countable history tree. A finite-prefix plan can be extended to an element of this product by arbitrary legal rows at later nodes. Its realization weights are finite products of coordinates, hence continuous. Each finite collection of the cost and score constraints is closed. For such a collection take the union of its finitely many model indices and its greatest horizon. Assumption (ii) and the restriction lemma give one feasible plan for every constraint in the collection. Extend it to all nodes. The closed constraints therefore have the finite-intersection property in a compact space, and one common infinite policy satisfies them all.

For this policy \eqref{eq:sharp-tail} gives uniformly vanishing survival probabilities, including the exact boundary $N=\underline c$. Thus it commits almost surely in every model and initial state. Its partial costs are nonnegative and at most $N$; monotone convergence yields the actual total cost bound. The norm of the score remainder is at most $\max(|\alpha_j|,|\beta_j|)\delta_T$. It tends to zero, so the bounded terminal scores have exactly the prescribed operator expectations. The program is common because the product coordinates were common throughout. Finally, negating the equivalence proves existence of a finite obstruction when no such program exists.
\end{proof}

\begin{corollary}[Attainment of the least finite cost]\label{cor:synthesis-attainment}
If any common proper program satisfies the prescribed targets at finite uniform cost, then the infimum of such cost budgets is attained.
\end{corollary}
\begin{proof}
Let feasible budgets decrease to their finite infimum $N_*$. For each fixed finite model set and horizon, the corresponding nonempty compact feasible sets decrease with the budget. Their intersection is nonempty and is feasible at $N_*$. Every finite problem at $N_*$ is therefore feasible; the theorem gives a proper program at that budget. The first-call floor ensures $N_*\ge\underline c$.
\end{proof}

\subsection{A strict dual certificate at an obstructing level}
List all left-minus-right matrix violations of \eqref{eq:hier-cost}--\eqref{eq:hier-tail} as $K_l(\gamma)$, $1\le l\le L_0$. Thus feasibility is $K_l\preccurlyeq0$ for every $l$. Define
\[
 v=\min_{\gamma\in\mathcal P_T}
       \max\{0,\max_l\lambda_{\max}(K_l(\gamma))\}.
\]
Compactness makes this minimum attained and $v>0$ exactly when the finite problem is infeasible.

\begin{theorem}[Strict finite dual obstruction]\label{thm:synthesis-dual}
The exact violation value is
\begin{equation}\label{eq:synthesis-dual}
 v=\max_{\substack{W_l\succeq0\\\sum_l\tr W_l\le1}}
       \min_{\gamma\in\mathcal P_T}\sum_l\tr[W_lK_l(\gamma)].
\end{equation}
The maximum is attained. An infeasible level has a positive-matrix separating family with strictly positive value. The inner minimization is a finite linear program in the common realization weights. No Slater hypothesis or inverse covariance is required.
\end{theorem}
\begin{proof}
For a fixed collection of violations, maximizing the displayed pairing over positive matrices with total trace at most one gives their largest positive eigenvalue, or zero by taking all matrices zero. Taking a minimizer defining $v$ gives the weak upper bound for every family $W$.

For the converse when $v>0$, fix $0\le\eta<v$. In the finite direct sum of Hermitian matrix spaces, the point $(\eta I_l)_l$ does not belong to the closed convex set
\[
 \{(K_l(\gamma))_l:\gamma\in\mathcal P_T\}
       +\prod_l\{D_l\succeq0\}.
\]
This set is closed because its first summand is compact. Finite-dimensional strict separation gives a nonzero linear functional whose infimum on that set exceeds its value at the excluded point. The functional must be nonnegative on the positive cone, or its infimum along a cone ray would be minus infinity. Hence it has the form of pairing with positive $W_l$. Their total trace is positive; normalize it to one. Separation yields an inner minimum greater than $\eta$. Let $\eta\uparrow v$ and take a convergent subsequence in the compact normalized positive-matrix set. The inner minimum is continuous in $W$, being the minimum of a continuous bilinear function over compact $\mathcal P_T$. Its limiting value is at least $v$. If $v=0$, a feasible plan gives the upper bound zero and $W=0$ attains zero. This proves equality and attainment in all cases.
\end{proof}

The finite dual is a genuine obstruction, not an assertion that the first obstructing horizon is uniformly bounded or effectively discoverable for arbitrary real primitive data. Existing quantum decision-process and reachability theories already contain substantial control, duality and computability results \cite{Barry14,Ying18,Saldi25}. Generic quantum MDPs, sequence form and convex separation are not claimed as new ingredients.

\subsection{Paid finite-prefix completion}
Suppose every open leaf reachable by a selected finite prefix admits one common legal fallback that commits with conditional expected cost at most $\kappa_0\ge\underline c$, uses at most $B_0$ live roles, and has terminal scores in the same $[\alpha_j,\beta_j]$ intervals. Its command and physical phase compatibility must be certified at these leaves.

\begin{corollary}[A paid finite-prefix completion]\label{cor:synthesis-fallback}
A feasible level-$T$ plan completed by this fallback has, for every model and initial state,
\begin{align}
 \E C&\le N+(\kappa_0-\underline c)\delta_T,\label{eq:fallback-budget}\\
 - (\beta_j-\alpha_j)\delta_T I
 &\preccurlyeq A_{\theta,j}-\E[\Phi_j]\preccurlyeq
       (\beta_j-\alpha_j)\delta_T I.\label{eq:fallback-error}
\end{align}
Here the score expectation is its operator on the initial state. If the finite prefix uses $L_T$ live node roles and the completed terminal decoder has $t$ actual values, a sufficient state-stopped label count is $L_T+B_0+t$ and an event-stopped count is $\max(L_T+B_0,t)$. Program constants, arithmetic and the possible exponential growth of $L_T$ are separate costs.
\end{corollary}
\begin{proof}
The completed cost is at most $\mathsf C_T+\kappa_0\mathsf S_T$. Subtract the reserved cost inequality and use the survival bound. Both the prescribed remaining score and the fallback's remaining score lie between $\alpha_j\mathsf S_T$ and $\beta_j\mathsf S_T$. Their difference lies between the stated opposite multiples of $\mathsf S_T$, proving the error bound. A finite prefix is implemented by storing its current node; the fallback uses its declared live roles. The exact role-packing lemma gives the two sufficient counts. These are existence counts, not a claim that a real-valued finite plan has an exact finite-bit sampler.
\end{proof}

\begin{proposition}[Sharpness of the uniform tail rate]\label{prop:tail-sharp}
For unit cost per call, every $b\ge1$ and integer $T\ge1$ with $b-1\le T$ admit a proper finite-branching program of mean cost $b$ and survival probability $(b-1)/T$ after $T$ calls.
\end{proposition}
\begin{proof}
With probability $p=(b-1)/T$ commit on call $T+1$, and otherwise on call one. Then $\E\tau=1+Tp=b$ and $\Pr(\tau>T)=p$. Its prefix has cost $1+(T-1)p$ and reserve $p$. Thus the reserve and the boundary $b=1$ are exact. If $b-1>T$, the universal tail bound one is attained by a deterministic sufficiently long path. This concerns the survival bound, not a uniform expected cost-tail estimate.
\end{proof}

\subsection{Synthesis of a finite marked squared-task reference}
For a finite virtual mark set $X$, fix masses $\nu_x$, target means $f(x)\in\R^k$ and nonnegative residuals $w(x)$ of a genuinely prepared reference law. Assume bounded audit coordinates. Use as terminal scores
\[
 \mathbf1_{\{X=x\}},\quad
 \mathbf1_{\{X=x\}}Z_i\ (1\le i\le k),\quad
 \mathbf1_{\{X=x\}}\|Z\|^2,
\]
with respective scalar-identity targets $\nu_x$, $\nu_x f_i(x)$ and $\nu_x[w(x)+\|f(x)\|^2]$. They are computed from the declared terminal audit instrument and route.

\begin{corollary}[Intrinsic synthesis of finite-mark squared-task balance]\label{cor:balance-hierarchy}
There is one common proper finite-branching acquisition program with these unnormalized marked moments and cost at most $N$ for all models and initial states if and only if every finite problem of the hierarchy for this finite score list is feasible. Such a program satisfies
\[
 \E[q(X)\|Z-r(X)\|^2]
   =\sum_x\nu_x q(x)[w(x)+\|f(x)-r(x)\|^2]
\]
for every selected readout. If, in addition, the same return rule is legal and the same certificate holds on every actual rooted continuation after rejection, it gives exact task balance for repeated acquisition. The hierarchy alone does not supply a finite live-memory bound.
\end{corollary}
\begin{proof}
The first equivalence is the hierarchy theorem applied to the explicit bounded score list. Expand the squared loss in each mark and substitute the three unnormalized moment identities. This proves the task equation without division by $\nu_x$, so zero-mass marks are harmless. Conversely the displayed marked moment requirement is exactly that finite target list. Applying this program repeatedly requires its hypotheses on the actual later roots; no equality of initial-state matrices can replace missing history-dependent continuation data. A finite prefix with a legal fallback, or the separate finite-memory theorem, supplies an explicitly charged implementation when its additional hypotheses hold.
\end{proof}

\section{Exact finite-memory synthesis on reachable quantum supports}\label{sec:finite-memory}
The preceding hierarchy characterizes full-history strategies. We now fix the observer alphabet and give a finite certificate that respects its persistent memory. The physical models in this section are time-homogeneous and finite-dimensional, with finite command and report sets and globally legal command rows. Restricted preparation faces and physical command domains have an explicit extension in the supplement. A finite model family is required for the finite algebraic synthesis conclusion. The analytic certificate itself can be imposed modelwise on any declared family, always with the same controller variables.

\subsection{The actual continuing instrument}
Fix $\ell$ live classical labels, a prescribed initial label $s_0$, and $t$ terminal decoder values $r_u$ in a compact carrier $D$. In a model $\theta$, the hidden physical algebra is $M_{d_\theta}(\mathbb C)$. A command $a$ has completely positive report maps $\mathcal I_{\theta,a,y}$ with normalized report sum, and expected-cost observable $G_{\theta,a}\succeq\underline c I$, $\underline c>0$. A terminal audit following a productive report has bounded score observable $F_{\theta,a,y,j}(r_u)$ on the post-report state. Every terminal choice is made before that audit.

The observer uses common command probabilities $p_{s,a}$. After command $a$ and report $y$, it uses common probabilities $q^0_{s,a,y,v}$ to continue at live label $v$, or $q^1_{s,a,y,u}$ to commit to decoder value $r_u$. These nonnegative probabilities sum to one. Nonproductive reports have all $q^1=0$. Only the current command, report and fresh randomness are available within that call; there is no external clock, port, persistent seed or hidden controller state.

In the event-stopped interface, live and terminal roles may share labels, and $\max(\ell,t)$ total labels suffice. In the state-stopped interface their types are disjoint and $\ell+t$ labels are required. Any fixed overlap or halt partition can instead be declared explicitly. The physical quantum state is part of the experiment, not a free coherent controller memory.

Use the block-diagonal live algebra $\mathcal A_\theta=\bigoplus_{s=1}^\ell M_{d_\theta}(\mathbb C)$. The actual continuing map and one-call observables are
\begin{align}
 (\mathcal T_\theta X)_v
 &=\sum_{s,a,y}p_{s,a}q^0_{s,a,y,v}\mathcal I_{\theta,a,y}(X_s),\label{eq:memory-T}\\
 (g_\theta)_s&=\sum_a p_{s,a}G_{\theta,a},\label{eq:memory-g}\\
 (f_{\theta,j})_s&=\sum_{a,y,u}p_{s,a}q^1_{s,a,y,u}
       \mathcal I_{\theta,a,y}^*(F_{\theta,a,y,j}(r_u)).\label{eq:memory-f}
\end{align}
The continuing map is completely positive and trace-nonincreasing. Its complementary one-call terminal effect is $b=I-\mathcal T^*(I)$. Normalization and the score bounds give
\begin{equation}\label{eq:memory-bound}
 g\succeq\underline c I,\qquad
 \alpha_jb\preccurlyeq f_j\preccurlyeq\beta_jb.
\end{equation}
We suppress $\theta$ in modelwise arguments. Let $\iota(\rho)$ insert an initial physical matrix in block $s_0$. The cost through $n$ calls is the pairing with $\sum_{k<n}(\mathcal T^*)^kg$, and the terminal score through $n$ calls is the pairing with $\sum_{k<n}(\mathcal T^*)^kf_j$. These formulas count the actual continuing paths, including failures.

\begin{lemma}[Finite determination of the reachable support]\label{lem:reachable-support}
Put $D_\theta=\ell d_\theta$ and
\begin{equation}\label{eq:reachable-Gram}
 R_\theta=\sum_{n=0}^{D_\theta-1}\mathcal T_\theta^n\iota(I),\qquad
 P_\theta=\operatorname{supp}R_\theta.
\end{equation}
This is the least block projection whose initial block is $I$ and whose support is invariant under the continuing instrument:
\begin{equation}\label{eq:support-invariant}
 \mathcal T(P)=P\mathcal T(P)P.
\end{equation}
Every actual live state is supported on $P$. Invariance is a support statement, not the equality $\mathcal T(P)=P$.
\end{lemma}
\begin{proof}
Let $R_n=\sum_{k=0}^n\mathcal T^k\iota(I)$. Its supports increase. If the supports of $R_n$ and $R_{n+1}$ agree, then the support of $\mathcal T(R_n)=\sum_{k=1}^{n+1}\mathcal T^k\iota(I)$ lies there. Since $R_n$ is strictly positive on its support, it dominates a positive multiple of that support projection. Positivity shows that the projection's image also lies there. The support is therefore invariant and all later terms remain inside it. Otherwise its rank strictly increases. Starting with nonzero rank in dimension $D_\theta$ proves stabilization by $D_\theta-1$.

Any invariant block projection containing the initial block contains every summand by positivity and hence their support. This proves leastness. Every raw positive reachable path operator is bounded by an appropriate positive sum in some iterate from $\iota(I)$, so its support lies in the least invariant projection. Equivalently apply the Kraus representation along the path. Thus all actual live states are contained in $P$.
\end{proof}

\subsection{An exact cost and task certificate}
For a block projection $P$ satisfying invariance, let $\mathcal T_P$ be the restriction to its supported algebra, with adjoint $\mathcal T_P^*X=P\mathcal T^*(X)P$. All inverses below act on this finite-dimensional supported algebra, not on unreachable blocks of the full system.

\begin{theorem}[Reachable-support potential certificate]\label{thm:memory-certificate}
For one fixed common observer, the following are equivalent:

\textup{(a)} In every model and from every initial density matrix it commits almost surely, has expected cost at most $N$, and has the prescribed terminal-score operators $A_{\theta,j}$.

\textup{(b)} In every model there are a block projection $P$, a positive $H=PHP$, and Hermitian $K_j=PK_jP$, with $P_{s_0}=I$, \eqref{eq:support-invariant}, and
\begin{align}
 P(H-\mathcal T^*H-g)P&\succeq0,& H_{s_0}&\preccurlyeq NI,\label{eq:memory-cost-cert}\\
 P(K_j-\mathcal T^*K_j-f_j)P&=0,& (K_j)_{s_0}&=A_{\theta,j},\label{eq:memory-score-cert}\\
 \alpha_jP&\preccurlyeq K_j\preccurlyeq\beta_jP.\label{eq:memory-K-bounds}
\end{align}
One may choose the explicitly computed least reachable projection in \eqref{eq:reachable-Gram}. On it, proper finite-cost operation is equivalent to
\begin{equation}\label{eq:memory-spectral}
 \operatorname{spr}(\mathcal T_P)<1.
\end{equation}
Under this condition the unique exact cost and score solutions are
\begin{equation}\label{eq:memory-resolvents}
 H^{\min}=(I-\mathcal T_P^*)^{-1}(PgP),\qquad
 K_j=(I-\mathcal T_P^*)^{-1}(Pf_jP).
\end{equation}
Every positive cost superpotential dominates $H^{\min}$ on this support. Testing its initial block and those of the $K_j$ is therefore necessary and sufficient, not merely a Lyapunov sufficient condition. Model-dependent certificates are analytic objects and are not stored by the common observer.
\end{theorem}
\begin{proof}
Assume (b) and work on the invariant support. Iterating the cost inequality gives
\[
 \sum_{k=0}^{n-1}(\mathcal T_P^*)^k(PgP)
       +(\mathcal T_P^*)^nH\preccurlyeq H.
\]
Trace against an initial density matrix, discard the nonnegative remainder and use the initial budget bound. The left partial sum is the actual cost through $n$ calls. Monotone convergence gives total expected cost at most $N$. The floor $PgP\succeq\underline cP$ bounds $\underline c\sum_{k<n}\Pr(\tau>k)$ by the same constant. Hence commitment is almost sure. The same argument gives finite cost from every state on $P$, with its bound supplied by $H$.

Iteration of the score equation yields
\[
 K_j=\sum_{k=0}^{n-1}(\mathcal T_P^*)^k(Pf_jP)
                    +(\mathcal T_P^*)^nK_j.
\]
The remaining trace is bounded in magnitude by $\max(|\alpha_j|,|\beta_j|)\Pr(\tau>n)$ and tends to zero. The finite sum records the scores on disjoint terminal events. Bounded scores and proper operation identify its limit with the actual terminal expectation. The initial block equation proves (a).

For necessity fix a model, let $X_0=\iota(I)$, take the least reachable projection and write $D=\ell d$. Assumption (a) gives
\[
 \sum_{k\ge0}\langle g,\mathcal T^kX_0\rangle\le dN.
\]
The corresponding series starting at $\mathcal T^nX_0$ is a nonnegative tail of this series. Sum these tails for $0\le n<D$ to obtain
\begin{equation}\label{eq:reachable-finite-cost}
 \sum_{k\ge0}\langle g,\mathcal T^kR\rangle\le DdN.
\end{equation}
Since $R$ is strictly positive on its support, every positive supported matrix is bounded by a finite multiple of $R$. Positivity proves finite expected cost from every such matrix. This includes states which need not occur as individual pure conditional states on any one path.

The positive partial cost potentials increase, and their pairing with $R$ is bounded by \eqref{eq:reachable-finite-cost}. Strict positivity of $R$ on its finite-dimensional support bounds their norms. They converge to a positive $H^{\min}$ solving $H^{\min}-\mathcal T_P^*H^{\min}=PgP$. Its initial block is precisely the cost operator and is at most $NI$.

The cost floor also makes $\sum_k(\mathcal T_P^*)^kP$ converge in norm. In particular $(\mathcal T_P^*)^nP\to0$. Every supported Hermitian matrix lies between $-\|X\|P$ and $\|X\|P$; positivity makes its image under these powers tend to zero. Decompose any complex matrix into Hermitian parts. In finite dimension, convergence of all these powers forces every eigenvalue to have modulus less than one. Otherwise an eigenvector of modulus at least one would not tend to zero. This proves \eqref{eq:memory-spectral} and its norm-convergent Neumann inverse.

The score series converges absolutely in norm by \eqref{eq:memory-bound} and the summable survival series. Normalization gives $PbP=P-\mathcal T_P^*P$, whose series sums to $P$. Consequently the score series sums to a $K_j$ between $\alpha_jP$ and $\beta_jP$, solves its equation and has the prescribed initial block. This proves (b) and both resolvent formulas.

Conversely, spectral radius less than one gives finite cost through the same Neumann series and proper operation through normalization. Budget and targets are then exactly the initial-block conditions. Finally any positive cost superpotential dominates every partial cost potential by the first telescoping inequality, and thus dominates their limit. More precisely, its nonnegative drift slack $D_H=P(H-\mathcal T^*H-g)P$ satisfies
\[
 H-H^{\min}=(I-\mathcal T_P^*)^{-1}D_H\succeq0.
\]
This proves both minimality and the exact excess-certificate formula.
\end{proof}

\subsection{An attained finite-memory algebraic frontier}
For synthesis the common probabilities, decoder values and modelwise certificate matrices become variables. The state-stopped interface is a finite union over permitted halt partitions. The event-stopped interface permits declared live/output reuse. Both keep a fixed initial live label.

\begin{corollary}[Finite algebraic synthesis at an exact memory budget]\label{cor:memory-algebraic}
Assume finitely many physical models with real algebraic primitive coefficients, a compact semialgebraic decoder carrier, and audit-score observables polynomial in decoder coordinates. Fix a total label budget $M$ and algebraic cost budget $N$. Existence of a common observer with the stated exact targets and cost at most $N$ is equivalent to feasibility of a finite system of polynomial equalities and positive-semidefinite inequalities. It is decidable by real-algebraic quantifier elimination. A feasible instance has a common controller and certificates with real algebraic entries.

For bounded continuous nonnegative terminal loss, the worst-model, worst-initial-state loss at budget $N$ attains its minimum whenever the feasible fixed-memory class is nonempty. Under the algebraic hypotheses its optimum is real algebraic. These are exact existence statements, not polynomial-time bounds, rational-optimum claims or finite-bit exact implementation of all algebraic samplers.
\end{corollary}
\begin{proof}
Use polynomial normalization and positivity for the probability rows. Encode Hermitian matrices by their real diagonal and real and imaginary off-diagonal coordinates. Positive semidefiniteness is equivalent to nonnegativity of every principal minor. Encode $P=P^*=P^2$, $P_{s_0}=I$, invariance, support of $H,K_j$, and the cost and score conditions of the theorem. All expressions in \eqref{eq:memory-T}--\eqref{eq:memory-f} and the certificate are polynomial in these variables. There are finitely many choices of live and terminal role counts and halt partitions under $M$. Every model retains its own certificate matrices but the same controller variables. For a loss upper bound $R$, replace its exact initial score target by $(K_{\rm loss})_{s_0}\preccurlyeq RI$. The theorem identifies this exactly with the required worst-state loss bound. Classical real-algebraic elimination proves decidability and algebraic witnesses \cite{Seidenberg54}; no new elimination algorithm is claimed.

Attainment needs a separate argument because the auxiliary potentials may be unbounded away from their initial block. The probability rows and decoder tuples form a compact space. Every finite-prefix cost and score is continuous in these parameters. A limit of observers with cost at most $N$ has every finite-prefix cost at most $N$. Monotone convergence and the positive floor give the same total budget and proper operation for the limit. The uniform bound $\Pr(\tau>T)\le N/[\underline c(T+1)]$ makes bounded terminal-score tails uniformly vanish on this budget class. Therefore full terminal scores are continuous there, as are the finitely many largest-eigenvalue worst-model losses. Exact score constraints are closed. The loss minimum on a nonempty compact feasible set is attained.

The achievable loss-bound set is semialgebraic over the real algebraic numbers by the first part. An attained finite endpoint of such a one-variable set is algebraic. Its nonempty fiber also has real algebraic coordinates. This proves the last statement without imposing a false compactness bound on the certificate matrices or claiming rational optimal probabilities.
\end{proof}

Finite-state POMDP synthesis, potential equations, spectral transience, quantum decision processes and real-algebraic methods have established antecedents \cite{Chatterjee15,Andriushchenko23,Barry14,Ying18,Saldi25}. The present claim is the precise common-model, all-initial-state, paid certificate with exact event/state semantics, its necessary reachable-support construction and its connection to acquired task balance. Those classical ingredients separately are not presented as new.

\begin{example}[An unreachable recurrent block is not an obstruction]\label{ex:unreachable-loop}
There are two live labels and one physical state. From the prescribed start $1$, a productive cost-one call commits immediately. Label $2$ is never reached and has a cost-one self-loop. The full continuing matrix is $\operatorname{diag}(0,1)$, with spectral radius one. Nevertheless the actual observer is proper and has cost one. Its reachable projection is $\operatorname{diag}(1,0)$; on that support the continuing matrix is zero and $H^{\min}=P$. Thus an unrestricted inverse of $I-\mathcal T$ or a global strict-transience requirement would incorrectly reject it. Scalar blocks embed this example in any quantum dimension.
\end{example}

\begin{corollary}[Same-task balance with a certified finite live count]\label{cor:memory-balance}
Apply the finite-memory certificate to the localized mass, first-moment and squared-norm scores in Corollary~\ref{cor:balance-hierarchy}. A feasible certificate gives one actual finite-label common balancing program at its certified cost. Conversely every such observer on the declared time-homogeneous interface yields the certificate. If the return rule and all actual rooted continuations are also certified, the program's actual live roles and exported decoder values may be used in the upper widths of the transfer theorem. Failure at a prescribed $M$ is an obstruction for that observer class, not merely failure of a searched single-call mixture.
\end{corollary}
\begin{proof}
The score list is finite and its quadratic expansion is exactly the unnormalized squared-task identity. The potential theorem is necessary and sufficient for the finite observer producing it. The common $p,q,r$ variables are the observer itself, so no uncounted history tape or model index is used. Its live labels and costs are those of its certified transitions. The rooted-continuation condition supplies the required uniformity when attempts repeat. The supplement gives the post-report compilation: prospective audits are not observed on rejected branches, and a logical subroutine return is not mistaken for an irrevocable outer halt. The transfer theorem then uses these same actual resource counts.
\end{proof}

\section{A finite test for reusable causal task balance}\label{sec:reusable}
The synthesis theorems produce an acquisition routine, but its first invocation is not the whole stopped experiment. A rejected report changes the physical state. Checking only the average returned state loses the report history on which rejection depends. Requiring task equality on every density matrix on the reachable Hilbert support, on the other hand, can be unnecessarily strong. The correct object for exact linear task identities is the report-resolved real linear span of the unnormalized returned states.

Fix a time-homogeneous routine with a finite visible return alphabet $X$. At a routine start, the physical state is a density matrix on a finite-dimensional carrier $K$. The routine is legal and proper on that carrier, as certified for example by Theorem~\ref{thm:memory-certificate}; all its physical operations and live roles are counted. On return $x$, a specified legal rejection operation produces the unnormalized next-start state
\[
 \mathcal R_x(\rho).
\]
The maps $\mathcal R_x$ are completely positive and $\sum_x\mathcal R_x$ is trace preserving. The rejection operation is trace preserving conditional on that return, so it does not change the return probability. It is performed only when the caller rejects and is charged separately when it has a physical cost. Its phase must be in the counted controller labels. There is no implicit physical reset.

For finitely many bounded prospective audit scores, let $F_{x,j}$ be the Hermitian start-state observable of the unnormalized score on return $x$, before the caller's commitment decision. The future audit is not observed on rejection. Put $F_{x,0}=\mathcal R_x^*(I)$ for report mass. Fix constants $m_{x,j}$, including $m_{x,0}=\nu_x$, with $\sum_x\nu_x=1$. The initial preparation family $\mathcal S$ is a nonempty set of density matrices on $K$; it need not contain every such matrix. All data may be indexed by an unknown model $\theta$, but the actual routine and caller rules are the same in every model. The following statements are imposed modelwise without allowing the controller to read $\theta$.

Let $d=\dim K$, let $\operatorname{Herm}(K)$ be the real vector space of Hermitian matrices, and define
\begin{align}
 W_0&=\operatorname{span}_{\R}\mathcal S, & r_0&=\dim W_0,\\
 W_{n+1}&=W_n+\sum_{x\in X}\mathcal R_x(W_n),& n_0&=d^2-r_0.\label{eq:return-span}
\end{align}
The additions are sums of real linear subspaces, not sums of channels or support projections.

\begin{theorem}[Finite report-resolved continuation criterion]\label{thm:return-span}
One has $W_{n_0}=W_{n_0+1}=:W_\infty$. This is exactly the span of all positive, unnormalized routine-start states
\[
 \mathcal R_{x_n}\cdots\mathcal R_{x_1}(\rho),
 \qquad \rho\in\mathcal S,\quad n\ge0.
\]
The following conditions are equivalent.

\textup{(i)} After every finite positive-probability history of rejected returns, from every preparation in $\mathcal S$, the next routine has conditional marked audit moments $m_{x,j}$, including its report masses $\nu_x$.

\textup{(ii)} Every residual $F_{x,j}-m_{x,j}I$ annihilates $W_\infty$ under the trace pairing.

\textup{(iii)} Those residuals annihilate $W_{n_0}$.

These conditions imply the same identities for every caller using the past returned marks and private randomness to choose acceptance, provided rejection uses exactly the declared return maps. They are also necessary when all finite sequences of rejections are legal caller choices. No unseen future audit, fresh model-dependent program or stored physical posterior is used by the caller.
\end{theorem}
\begin{proof}
An equality $W_n=W_{n+1}$ makes $W_n$ invariant under every $\mathcal R_x$, so all subsequent subspaces agree. Before equality each dimension increases by at least one, starting at $r_0$ and never exceeding $d^2$. Stabilization therefore occurs by $n_0$. Induction expands $W_n$ as the span of word images of length at most $n$ of $W_0$; taking the invariant limit proves the first assertion.

For a word $w=(x_1,\ldots,x_n)$ and $\rho\in\mathcal S$, write $A_w=\mathcal R_{x_n}\cdots\mathcal R_{x_1}(\rho)$. Complete positivity gives $A_w\succeq0$. If its trace is positive, the conditional next-start state is $A_w/\tr A_w$. Its marked expectation is $\tr(F_{x,j}A_w)/\tr A_w$, so equality to $m_{x,j}$ is precisely $\tr[(F_{x,j}-m_{x,j}I)A_w]=0$. If its trace is zero, positivity gives $A_w=0$, and the same unnormalized equality is automatic. Thus (i) is equivalent to annihilation of all word images, hence to (ii). Stabilization gives (iii).

A caller's randomized rejection decisions contribute scalar nonnegative weights to each fixed visible word, because its choices use only that word and private randomness, not the hidden matrix. Conditioning also on its private tape changes these weights but not the physical word map. Integrating over compatible tapes preserves every displayed linear identity. Therefore the identities hold at every actual surviving caller history. Conversely deterministic rejection of the first $n$ returns makes each positive-probability word of length $n$ a possible start history, proving necessity. These arguments use finite histories; conditional-hazard summation, rather than a claim of renewal, supplies the infinite stopping conclusion below.
\end{proof}

\begin{corollary}[A finite-dimensional certificate for repeated finite-mark balance]\label{cor:reusable-balance}
Take the mass, first-moment and squared-norm score targets of Corollary~\ref{cor:balance-hierarchy}. Suppose their residuals annihilate the report-resolved span in every model, and one invocation including its executed rejection return costs conditionally at most $\kappa$ on all actual start states. For any finite dictionary and any feasible $0\le h\le b$ with $\sum_x h(x)\nu_x=1$, accept return $x$ with probability $h(x)/b$. The resulting routine commits almost surely, costs at most $\kappa b$, and has exact terminal task value
\[
 \sum_x h(x)\nu_x\,[w(x)+\|f(x)-r(x)\|^2].
\]
The same finite observer implements every attempt, with no added intercall history register. A rejection return performed within the productive call adds no live role; an extra physical return subroutine must have its calls and live roles counted. This yields the upper construction needed by the transfer theorem from a finite report-resolved continuation test, rather than leaving its history-uniform task identity as an unverified premise.
\end{corollary}
\begin{proof}
The three marked moment identities and expansion of the square give the same unnormalized task identity conditional on every actual start history by the theorem. Thus the conditional acceptance probability is $1/b$ and the unnormalized loss is $b^{-1}$ times the displayed value. Induction gives survival $(1-1/b)^n$ through $n$ invocations; summing begun invocations gives expected number $b$. Summing their conditional costs and disjoint terminal losses proves both assertions. For a final accepted invocation the rejection operation is not performed, and nonnegative cost only improves its budget. The decision is made at the final productive report before the prospective audit. The existing routine's live labels are reused on rejection, while its finite terminal decoder is read only on commitment. This proves the memory statement with the explicit return-operation qualification.
\end{proof}

The cost assumption in this corollary is not silently converted into a linear-span equality. A positive cost superpotential on the reachable support is a sufficient, checkable bound. If the exact whole-invocation cost observable $H$ is available and constant cost $\kappa$ is required at every reachable start, the condition is exactly $(H-\kappa I)\perp W_\infty$ by the same proof. A one-sided inequality on only the reachable positive states need not be equivalent to a Loewner inequality on their whole support; no such equivalence is asserted.

\begin{example}[Averaging is too weak; support equality is too strong]\label{ex:return-span}
On a qubit let $\rho_0=I/2$ and
\[
 \mathcal R_0(A)=\tfrac12\tr(A)|0\rangle\langle0|,
 \qquad
 \mathcal R_1(A)=\tfrac12\tr(A)|1\rangle\langle1|.
\]
The average return sends $\rho_0$ to itself, so iterating only that average generates the one-dimensional span of $I$. But the report-resolved span is the entire diagonal Hermitian subspace. The observable $\sigma_z$ has zero expectation at the initial and averaged returned state, and expectations $+1,-1$ at the two actual conditional starts. An initial or pooled calibration of that task therefore fails the required repeated identity.

Conversely $\sigma_x$ annihilates the complete report-resolved diagonal span, though it is a nonzero operator on its full two-dimensional Hilbert support. A test insisting on the operator identity $\sigma_x=0$ on that support would reject an exact reachable-task identity. Both examples can be realized by a fair visible coin, a prospective Pauli audit at commitment, and a charged conditional replacement preparation on rejection. They distinguish a span of density matrices from a Hilbert-space support and do not make replacement preparations free.
\end{example}

\begin{proposition}[Finite exact arithmetic verification]\label{prop:return-arithmetic}
For a fixed finite-dimensional routine whose return maps, task observables and a finite spanning preparation list have algebraic real and imaginary parts, the continuation criterion is decidable by finitely many exact linear-algebra operations. At most $d^2-r_0$ span enlargements are needed; each enlarges a space of dimension at most $d^2$. For a finite family of physical models the same test is performed in every model on the one fixed controller. This is a verification result for the supplied routine, not a polynomial-time synthesis theorem for its controller parameters or a complexity bound for evaluating arbitrary algebraic numbers.
\end{proposition}
\begin{proof}
Represent Hermitian matrices by $d^2$ real coordinates and every return map by a real matrix on those coordinates. Starting with a basis of the supplied preparation span, append all its images under all return maps and retain a basis by exact elimination. The dimension argument in Theorem~\ref{thm:return-span} bounds the number of enlargements. At stabilization check the finitely many trace pairings with every residual. All operations and zero tests are exact in the finite algebraic field generated by the input entries. This implements the equivalence of the theorem.
\end{proof}

Finite reachable linear spans and word-equivalence tests have classical probabilistic-automata antecedents \cite{Tzeng92}; reachable Hilbert supports and quantum termination also have established theory \cite{LiYuYing12,Fu24}. Those algorithms and the dimension argument are not claimed as new. The assertion used here is the exact report-resolved test for repeated marked tasks, its distinction from pooled calibration and support-wide identities, and its paid post-report compilation into the same causal resource theorem.

\subsection{Continuous reports and a sharp finite-depth obstruction}
The finite-return restriction can be removed from the exact continuation test without asserting general effective synthesis. Let $X$ now be standard Borel. A rejected return is described by a completely positive instrument $D\mapsto\mathcal R(D)$, countably additive in the finite-dimensional space of linear maps, with $\mathcal R(X)$ trace preserving. Let $F_j(D)$ be finite Hermitian operator-valued prospective task measures and $m_j(D)$ their prescribed finite scalar measures; include report mass as $F_0(D)=\mathcal R(D)^*I$ and $m_0=\nu$. Each bounded score gives such a finite measure. Define
\[
 W_{n+1}=W_n+\operatorname{span}_{\R}\{\mathcal R(D)A:
        D\in\mathcal B(X),\ A\in W_n\},\qquad
 E_j(D)=F_j(D)-m_j(D)I.
\]
The operational identity below is an almost-sure identity at actual reached histories, not a condition imposed at individually null continuous reports.

\begin{theorem}[Finite-depth criterion for standard Borel return instruments]\label{thm:borel-return}
With $W_0=\operatorname{span}_{\R}\mathcal S$ of dimension $r_0$ and physical dimension $d$, the span stabilizes by depth $d^2-r_0$. Exact marked task identities after every finite report-dependent rejection history hold if and only if
\begin{equation}\label{eq:borel-return}
 \tr[E_j(D)A]=0\quad
 (A\in W_{d^2-r_0},\ D\in\mathcal B(X),\ \text{every }j).
\end{equation}
If an identity fails, there is a witness using an allowed initial preparation, at most $d^2-r_0$ rejected returns, a rectangular cylinder of report events of positive probability, and one subsequent task event.

There exists a finite Borel partition $\Pi$ of the return space with at most $2^{d^2-r_0}$ cells such that the maps $\{\mathcal R(C):C\in\Pi\}$ generate exactly the same reachable span as all Borel events. This is an existence statement about exact instrument measures. An arbitrary coarse partition or numerical binning is not such a certificate, and testing task-measure equality still requires its exact measure data.
\end{theorem}
\begin{proof}
Dimension growth and invariance give stabilization exactly as before. Induction identifies $W_n$ with the span of the matrices
$\mathcal R(D_n)\cdots\mathcal R(D_1)\rho$ for words of length at most $n$ and $\rho\in\mathcal S$. These are positive unnormalized states for rectangular cylinders. If a residual does not annihilate their span, it does not annihilate one such generator. That generator is positive and nonzero, so its trace is positive. This proves the bounded-depth witness assertion and necessity of the displayed condition.

For sufficiency, fix an initial preparation and a finite depth. Iterating the instrument gives a finite positive operator-valued measure on $X^n$. Pairing it with a fixed residual $E_j(D)$ gives a finite signed scalar measure. Its value on every rectangular cylinder is zero by the invariant-span condition. Rectangles form a generating pi-system, so uniqueness of finite signed measures makes it zero on every measurable history event. A bounded nonnegative report-dependent rejection weight can be integrated against this zero measure; approximation by simple functions proves that the conditional task discrepancy is zero at the actual surviving histories. To obtain one common almost-sure set for all next task events, first use a countable generating algebra of the standard Borel report space, intersect its full-measure sets for the finitely many scores, and then use uniqueness of finite measures at each remaining history. Private caller randomness is integrated in the same way. This gives the asserted operational identities, including measurable rejection depending on the full report history rather than only on rectangles.

For the finite partition, start with $V=W_0$. If $V$ is not invariant under all event maps, choose an event $D$ and an $A\in V$ such that $\mathcal R(D)A\notin V$. Add this event to a list and replace $V$ by the invariant span generated by all listed event maps and $W_0$. Its dimension increases by at least one. After at most $d^2-r_0$ such choices, the resulting space is invariant under every event map and hence equals $W_\infty$. Take the common Boolean refinement of the listed events. Every selected map is the sum of the maps of its partition cells. The invariant span under the partition maps therefore contains $W_\infty$, while each cell is a Borel event and so cannot generate a vector outside $W_\infty$. Equality follows. At most $2^{d^2-r_0}$ cells are needed. Nothing in this argument supplies an effective representation of those selected Borel events for arbitrary input measures.
\end{proof}

\begin{theorem}[Sharp and physically realizable continuation depth]\label{thm:sharp-return-depth}
For every physical dimension $d\ge2$, there is a completely positive trace-preserving return channel, one initial full-rank preparation, one visible return mark and a bounded two-outcome prospective audit such that the prescribed squared-task balance holds after $n$ rejections for every $0\le n<d^2-1$, but fails after $d^2-1$ rejections. All real and imaginary parts of primitive matrix entries can be chosen real algebraic. Thus the witness depth $d^2-r_0$ is sharp already for $r_0=1$, even with one return mark and uniformly finite paid invocation cost.
\end{theorem}
\begin{proof}
Put $q=d^2-1$. Choose a Hilbert--Schmidt orthonormal real basis $F_1,\ldots,F_q$ of traceless Hermitian matrices with algebraic entries. A generalized matrix-unit basis gives one, and $\|F_i\|_{\rm op}\le1$. Choose distinct numbers
\[
 \lambda_i=\frac{i}{2dq(q+1)},\qquad
 \varepsilon=\frac1{2dq},\qquad
 \rho_0=I/d+\varepsilon\sum_{i=1}^q F_i.
\]
The norm bound makes $\rho_0\succeq I/(2d)$ and its trace is one. Define
\[
 \Phi(A)=\tr(A)I/d+\sum_{i=1}^q\lambda_i\tr(F_iA)F_i.
\]
Its Choi matrix in the unnormalized convention is
\[
 J_\Phi=I\otimes I/d+\sum_i\lambda_i F_i^{\mathsf T}\otimes F_i.
\]
Since $\sum_i\lambda_i=1/(4d)$, the perturbation has norm at most $1/(4d)$. Thus $J_\Phi\succeq3I/(4d)$, which proves complete positivity. Taking traces proves trace preservation, and $\Phi(I)=I$, $\Phi(F_i)=\lambda_iF_i$. Consequently
\[
 \rho_n=\Phi^n(\rho_0)=I/d+\varepsilon\sum_i\lambda_i^nF_i.
\]
The $q+1=d^2$ matrices $\rho_0,\ldots,\rho_q$ are linearly independent by the Vandermonde determinant on the distinct nodes $1,\lambda_1,\ldots,\lambda_q$. There is therefore a Hermitian $L$ annihilating $\rho_0,\ldots,\rho_{q-1}$ but not $\rho_q$. This already proves the claimed depth. We give exact coefficients to make the physical realization explicit.

Set $x_0=1$, $x_i=\lambda_i$ and $w_i=\prod_{j\ne i}(x_i-x_j)^{-1}$. The coefficient of degree $q$ in Lagrange interpolation of $x^n$ at these $q+1$ nodes gives
\[
 \sum_{i=0}^q w_i x_i^n=0\ (0\le n<q),\qquad
 \sum_{i=0}^q w_i x_i^q=1.
\]
Let $L_0=w_0I+\varepsilon^{-1}\sum_{i=1}^q w_iF_i$ and
$K=|w_0|+\varepsilon^{-1}\sum_{i=1}^q|w_i|>0$.
Then $L=L_0/K$ is Hermitian with operator norm at most one, and
$\tr(L\rho_n)=0$ for $n<q$, while $\tr(L\rho_q)=1/K>0$.
The effects $(I+L)/2,(I-L)/2$ give a valid two-outcome audit $Z\in\{-1,1\}$ with mean $\tr(L\rho)$ and second moment one. All specified matrices and constants are algebraic; the square roots needed for a Kraus implementation of their positive Choi matrix may also be chosen algebraic by finite-dimensional spectral calculus.

One routine emits a deterministic productive mark without changing the state. On acceptance, commit a readout and perform this prospective audit. On rejection, execute the physical channel $\Phi$ and return to the start. Charge the productive operation and rejected return separately; each invocation has cost at most two unit calls and a fixed finite implementation. After $n$ rejected invocations the actual start state is $\rho_n$. For the genuinely prepared initial reference, $f=0,w=1$, and the squared loss at readout $r\in[-1,1]$ is $1+r^2$ at every $n<q$. At $n=q$ it is $1+r^2-2r/K$, which differs for any nonzero $r$. Thus all marked mass, mean and squared-norm identities hold up to the asserted depth and the task identity then fails. No report of zero probability, infinitely expensive routine, free physical reset or numerical approximation is used.
\end{proof}

The Hermitian reachable-space method itself, including its distinction from Hilbert-space reachability, is expressly present in Fu et al.\ \cite{Fu24}; it is not introduced here. The extensions proved above are the standard Borel report-event formulation for reusable marked tasks, its finite-coarsening statement, and the matching physically realizable $d^2-1$ witness depth connected to paid rejection. Neither the finite witness bound nor its sharp example constitutes a general decision procedure for arbitrary continuous instrument measures.

\section{Primitive optimal control for finite reset instruments}\label{sec:reset-control}
An exact application should not require that a balancing group has already been built into the experiment. We give a class in which the best continuation and a least-favorable preparation can instead be derived from the primitive measurement instrument. The class includes any finite family of rank-one failure projections in any finite Hilbert dimension, as well as mixed post-failure states. It does not subsume an arbitrary-rank failure projection, whose posterior still depends on its input. The complete arbitrary-rank two-detector theorem is retained in the next section.

Let $\mathcal H$ be finite-dimensional. There is a nonempty finite command set $\mathcal A$, nonterminal reports $i=1,\ldots,s$, and a terminal success. Assume $s\ge1$. A nonterminal branch has the physically specified map
\begin{equation}\label{eq:reset-instrument}
 \mathcal I_{a,i}(\rho)=\tr(E_{a,i}\rho)\sigma_i,
 \qquad E_{a,i}\succeq0,\quad \sum_i E_{a,i}\preccurlyeq I,
\end{equation}
where $\sigma_i$ is a known density matrix and the report $i$ is observed. The aggregate success effect is $B_a=I-\sum_iE_{a,i}$. The instrument itself prepares the state $\sigma_i$ on failure; no additional freely available reset is assumed. One call has a nonnegative cost with conditional mean $\tr(G_a\rho)$, where
\[
 0<c_*I\preccurlyeq G_a\preccurlyeq c^*I<\infty.
\]
Success requires immediate commitment and failure forbids it. The decoder has one fixed terminal value. All commands are legal. The initial state is arbitrary and unknown to the observer; the known post-report state is part of the physical model, not an unknown-state oracle. Commands and reports may be relabelled without changing these data.

Set
\[
 Q_{iaj}=\tr(E_{a,j}\sigma_i),\qquad c_{ia}=\tr(G_a\sigma_i),
 \qquad q_{ia}^{\rm suc}=1-\sum_jQ_{iaj}.
\]
A nonempty set $S\subset\{1,\ldots,s\}$ is a \emph{closed failure set} when $\sum_{j\in S}Q_{iaj}=1$ for every $i\in S$ and every command $a$.

\begin{theorem}[Finite reset classification, optimum and preparation saddle]\label{thm:reset-control}
There exists one common visible-history policy which is proper at finite uniform expected cost for every initial density matrix if and only if there is no closed failure set. If there is none, the finite Bellman system
\begin{equation}\label{eq:reset-bellman}
 v_i=\min_{a\in\mathcal A}\left(c_{ia}+\sum_jQ_{iaj}v_j\right),\qquad 1\le i\le s,
\end{equation}
has a unique finite nonnegative solution. It is the optimal remaining cost from each reported state $\sigma_i$. A deterministic minimizing command at every reported state attains it and is proper.

Define Hermitian start-cost operators
\[
 L_a=G_a+\sum_jv_jE_{a,j}.
\]
The unknown-initial-state value is exactly
\begin{align}
 V_*&=\min_{p\in\Delta(\mathcal A)}\lambda_{\max}\left(\sum_a p_a L_a\right)\label{eq:reset-sdp}\\
 &=\max_{\rho\succeq0,\ \tr\rho=1}\min_a\tr(\rho L_a).\label{eq:reset-dual}
\end{align}
Both extrema are attained. A minimizing initial command distribution followed by the deterministic Bellman continuation is one robustly minimax policy whose uniform bound holds for every initial state; its initial row need not minimize cost separately at each known state. A maximizing density matrix is a primitive least-favorable preparation against every visible-history policy. The latter statement is for unrestricted history; it is not automatically a saddle at every smaller memory budget.

If $g$ distinct command values occur in a chosen deterministic Bellman selector, then $g+1\le\min(s,|\mathcal A|)+1$ live labels suffice with a prescribed initial label. If a minimizing initial distribution is a pure command already used by that selector, $g$ suffice. Event stopping permits reuse of a live label for the constant terminal value; state stopping adds one disjoint terminal label.
\end{theorem}
\begin{proof}
If a closed failure set exists, start in any $\sigma_i$ with $i\in S$. Every command fails with probability one and reports another state in $S$. Induction makes success impossible under every policy; the positive cost floor makes its expected cost infinite.

Suppose there is no such set. Choose every command independently with strictly positive probability. From each failure state there is a positive-probability directed path to success: otherwise the set of failure states reachable from that state would be nonempty and closed under every command, contradicting the assumption. A shortest such path visits no failure state twice and has at most $s$ edges. Finiteness gives a positive lower bound $\delta$ for the probability of following one such path in the next $s$ calls, uniformly in its starting failure state. Consequently survival after $ns$ calls is at most $(1-\delta)^n$. From an arbitrary initial state the first call either succeeds or enters one of the reported states. The expected number of calls is therefore uniformly at most $1+s/\delta$, and the expected cost at most $c^*(1+s/\delta)$. This proves existence using even a one-live-label randomized command row. Its role is feasibility, not optimality.

For explicit primitive constants, if the row chooses each command with probability at least $\eta>0$ and $\epsilon>0$ is the smallest positive number among the finite collection $Q_{iaj},q_{ia}^{\rm suc}$, then one may take $\delta=(\eta\epsilon)^s$. Factors for a shorter path can be bounded below by this same number. Zero probabilities are not divided out or assigned a false positive floor.

Start value iteration at $v^{(0)}=0$. The finite-horizon dynamic programming recursion gives $v^{(n+1)}_i=\min_a(c_{ia}+\sum_jQ_{iaj}v^{(n)}_j)$. These values increase and are bounded by the proper randomized policy's finite cost. Their limit $v$ is finite and solves \eqref{eq:reset-bellman}, since a finite minimum commutes with this coordinatewise limit. Select a minimizing command $\alpha(i)$ in each row. Iterating its equality bounds every partial expected cost by $v_i$. The positive floor then bounds the sum of survival probabilities, so the selected policy is proper and has finite cost. The bounded terminal remainder $\E[v_{I_n}\one_{\{\tau>n\}}]$ vanishes. Passing to the limit in the iterated equality identifies its cost as $v_i$.

For any other proper policy, Bellman's inequality for each chosen command telescopes to $v_i$ no greater than its expected cost plus the same bounded survival remainder. Letting the horizon grow proves its cost at least $v_i$. An infinite-cost policy cannot improve this bound. The same argument for any finite nonnegative Bellman solution $w$ shows that its greedy policy realizes $w$ and that $w$ is below the cost of every proper policy. Comparing to $v$ in both directions proves uniqueness.

Before the first call the observer has no report, so its first command has one model-independent distribution $p$. After a failure $j$, the physical state is exactly $\sigma_j$, irrespective of the original state or previous private coins. The already proved continuation lower therefore makes its initial cost at least
$\tr[\rho\sum_a p_aL_a]$. The Bellman continuation attains this bound for every $\rho$. Taking the worst density matrix gives the largest eigenvalue, proving \eqref{eq:reset-sdp}. Compactness of the probability simplex gives its minimum.

For completeness the dual equality follows directly from finite-dimensional separation. Let $\mathcal K=\operatorname{conv}\{L_a\}$ and let $v_* =\min_{L\in\mathcal K}\lambda_{\max}(L)$. The compact convex set $\mathcal K$ is disjoint from the open convex set $\{H:H\prec v_*I\}$. A separating real linear functional on Hermitian matrices is pairing with a nonzero Hermitian $R$, with $\inf_{L\in\mathcal K}\tr(RL)\ge\sup_{H\prec v_*I}\tr(RH)$. Finiteness of the right side forces $R\succeq0$, since otherwise moving $H$ arbitrarily far in a negative semidefinite direction makes that pairing unbounded above. Normalize $\tr R=1$; the right supremum is then $v_*$. This gives a density matrix with $\min_a\tr(RL_a)\ge v_*$. Conversely, pairing any density matrix with a minimizing $L$ bounds that minimum by $v_*$. This proves equality and dual attainment. The continuation lower makes the maximizing state a witness against every policy, not merely against first-call mixtures.

To implement the policy, use one label for its initial distribution and a pure command label for each distinct value of $\alpha$. On failure report $j$, retain the label for $\alpha(j)$. The next row depends only on that retained label. Success commits immediately. If the initial row is already one of these pure rows, use that existing label as the start. The event/state packing lemma gives the terminal counts. No running clock, earlier report or persistent seed is available in addition to the labels.
\end{proof}

\begin{proposition}[Equality, rigidity and a memory lower certificate]\label{prop:reset-rigidity}
Let $\rho_*$ be an attaining preparation in \eqref{eq:reset-dual}, and put
\[
 d_{ia}=c_{ia}+\sum_jQ_{iaj}v_j-v_i\ge0,\qquad
 \mathcal A_i^*=\{a:d_{ia}=0\}.
\]
For every finite-cost policy with first-command distribution $p$, its cost satisfies the exact residual identity
\begin{equation}\label{eq:reset-residual}
 \E_{\rho_*}C
 =\tr\!\left(\rho_*\sum_a p_aL_a\right)
    +\E_{\rho_*}\sum_{n\ge1:\,\tau>n}d_{I_n,A_{n+1}}.
\end{equation}
Consequently an unrestricted optimum uses only commands in $\mathcal A_i^*$ at every positive-probability continuation from reported state $i$, under this preparation.

Suppose a set $I$ of reported states is visited with positive probability by every robust optimum evaluated at $\rho_*$. Let $\chi(I)$ be the least number of subsets covering $I$ whose respective intersections of optimal command sets are nonempty. Every fixed-row autonomous optimum needs at least $\chi(I)$ live labels. If every minimizing initial row has support not contained in $\mathcal A_i^*$ for every $i\in I$, it needs at least $\chi(I)+1$. These are checkable lower certificates once the mandatory set and the initial-row condition have been proved, not assertions that every experiment has the same minimal memory.
\end{proposition}
\begin{proof}
After each failed call, subtract $v_{I_n}$ from the cost of the next call plus the conditional next potential. The conditional difference is exactly $d_{I_n,A_{n+1}}$. Sum to a deterministic horizon, including the first-call cost and potential. The remaining potential has absolute value at most $\max_i v_i$ times the survival indicator. Finite expected cost and the floor make survival tend to zero. Monotone convergence for the nonnegative residual sum yields \eqref{eq:reset-residual}. Its first term is at least $V_*$ by the dual witness. A robust optimal policy's cost at $\rho_*$ is at most $V_*$, so every visited positive-probability residual vanishes.

A retained label serving states in a subset $J$ has one command row. Its support must lie in every $\mathcal A_i^*$, $i\in J$, because each such visit has positive probability and any action of positive row weight with positive residual would contradict equality. Thus $\bigcap_{i\in J}\mathcal A_i^*$ is nonempty. The labels used by mandatory states cover $I$, proving the first lower. For an optimum, its first row is a minimizing row in \eqref{eq:reset-sdp}: its worst-state cost dominates the first-call lower, whose minimum is $V_*$. Under the additional support condition this initial row can serve none of the mandatory states, so its label is additional to their covering labels. This proves the sharpened count.
\end{proof}

\begin{corollary}[Arbitrary finite rank-one detector families]\label{cor:rank-one-family}
Let the failure projection of detector $a$ be $P_a=|u_a\rangle\langle u_a|$ on any finite-dimensional Hilbert space, with $\|u_a\|=1$, and let its deterministic call cost be $k_a>0$. There is no intervening evolution or additional control. If all rays $\mathbb C u_a$ coincide, that ray is dark and the all-state value is infinite. Otherwise the value is finite and is given by Theorem~\ref{thm:reset-control}, with
\[
 v_i=\min_a\{k_a+|\langle u_i,u_a\rangle|^2v_a\},
 \qquad L_a=k_aI+v_aP_a.
\]
Thus the continuation is a finite scalar system, while the unknown-initial-state optimization retains the full spectral geometry of the physical projections. There is no fixed prior or restriction to two detectors or two physical dimensions.
\end{corollary}
\begin{proof}
A rank-one failure satisfies $P_a\rho P_a=\tr(P_a\rho)P_a$, so it is a reset branch with report $a$ and state $P_a$. If all rays agree, their common ray never succeeds. Otherwise choose two distinct rays and alternate their detectors. Their squared overlap is less than one. After the first failure every switch has this same strictly positive success probability, giving finite uniform cost. The classification theorem applies. Substitution of the reset transition probabilities gives the displayed formulas. The general equality and memory certificates retain the actual observed failure label and therefore apply to this family as well.
\end{proof}

The reduction to a finite absorbing decision problem, Bellman verification, convex separation and spectral duality are classical methods; quantum decision-process frameworks already provide much broader control languages \cite{Barry14,Ying18,Saldi25}. The contribution claimed here is the explicit primitive reset-instrument criterion, the common all-initial-state spectral value, its realization with charged rows, and the residual-based memory certificate within this experiment. It is distinct from optimizing the probing interval for a fixed detector with unitary evolution \cite{Kessler21,DelVecchio26}. It also does not replace the preserved two-stage reset communication and rank-constrained discrimination results by a first-detection theorem.

\section{A spectral memory law for two quantum detectors}\label{sec:two-detectors}
We apply the finite-memory framework without a balancing group action or reset to a prescribed reference state. The measurements disturb the state: this is first detection, not estimation of an untouched pre-command parameter. Quantum first-detection laws, dark states and probing-time optimization have substantial antecedents \cite{Kessler21,DelVecchio26}. Our question is the exact paid-call versus classical-label frontier when the observer chooses between two specified detectors.

Let $P,Q$ be orthogonal projections on a finite-dimensional Hilbert space. Detector $A$ has failure map $\rho\mapsto P\rho P$ and success map $\rho\mapsto(I-P)\rho(I-P)$; detector $B$ uses $Q$ in the same way. Every call costs one. There is no intervening evolution, additional measurement, reset or coherent controller memory. A failure does not permit commitment. The first success requires immediate productive commitment, with one constant terminal decoder value. The initial density matrix is arbitrary and unknown. An autonomous controller has a prescribed initial live label, fixed command rows and post-failure transitions, with no free clock or persistent external seed. Write $V_j$ for the least worst-initial-state expected call count using at most $j$ live labels, and $V_*$ for the infimum over all visible-history policies.

\begin{theorem}[Exact two-detector resource law]\label{thm:detector-memory}
Suppose $P,Q\ne0$ and $c=\|PQ\|<1$. In every finite dimension,
\begin{equation}\label{eq:detector-global}
 V_*=1+\frac1{2(1-c)}.
\end{equation}
It is attained by choosing the first detector fairly and then switching after every failure. The one-live-label value is
\begin{equation}\label{eq:detector-one}
 V_1=1+\frac1{1-c}=2V_*-1,
\end{equation}
and a fair independent detector choice at every call attains it. With a fixed initial label, three live labels are necessary and sufficient to attain $V_*$. The two-live-label optimum is attained and satisfies
\begin{equation}\label{eq:detector-two}
 V_*<V_2\le1+\frac1{1-c^2};
\end{equation}
its exact value is determined in Theorem~\ref{thm:two-label-exact} below.

If an explicitly randomized initial distribution over counted labels is instead part of the interface, two live labels are necessary and sufficient to attain $V_*$. Under event stopping the stated live counts are also total-label counts, since the constant terminal value may reuse a live label. State stopping requires one additional disjoint terminal label.

If either failure projection is zero, one call to that detector gives value one. If $c=1$, their ranges have a common nonzero dark subspace and no policy is proper for every initial state. The robust cost is then infinite.
\end{theorem}
\begin{proof}
We give the required spectral reduction explicitly. Diagonalize $PQP$ on $\operatorname{ran}P$. For an eigenvector $u$ with eigenvalue $q\in(0,1)$, put $v=Qu/\sqrt q$. Then $\|v\|=1$ and $\langle u,v\rangle=\sqrt q$. Their span is invariant under both projections, which restrict to the rank-one projections onto $u,v$. Distinct orthonormal eigenvectors give orthogonal such two-dimensional blocks, including after resolving multiplicities. The zero eigenspace gives $\operatorname{ran}P\cap\ker Q$; the remaining part of $\operatorname{ran}Q$ is orthogonal to $P$. The orthogonal complement of the two ranges lies in both kernels. Eigenvalue one is exactly a common failure range. This is the classical two-subspace geometry \cite{Halmos69}, proved here in the form used below. The largest nontrivial squared overlap is $q=c^2$ when $c>0$.

After a failure of $A$, a second immediate $A$ fails with probability one and leaves the state unchanged. The analogous statement holds for $B$. Such a repetition gives no information. In a full-history policy it can be deleted, with its private coins and deterministic failure simulated internally, leaving all subsequent physical states and detector choices coupled and weakly reducing calls. Infinite repetitions give infinite cost and cannot improve the optimum. Thus, after its first choice, a pointwise optimal policy switches after every failure.

On a block of squared overlap $q$, the post-failure pure state switches between $u$ and $v$, and every further failure has probability $q$. Its expected remaining cost is $1/(1-q)$. On the whole $P$ failure range this continuation cost is the positive operator $H_P=(I_P-PQP)^{-1}$, and similarly $H_Q=(I_Q-QPQ)^{-1}$. These inverses exist because $c<1$. A first choice of $A$ with probability $r$ therefore has cost operator bounded below by
\begin{equation}\label{eq:detector-lower-operator}
 I+rPH_PP+(1-r)QH_QQ,
\end{equation}
with equality for subsequent alternation. On a principal block its largest eigenvalue is
\begin{equation}\label{eq:detector-block}
 1+\frac{1+\sqrt{1-4r(1-r)(1-q)}}{2(1-q)}.
\end{equation}
It increases with $q$ and is minimized uniquely at $r=1/2$. Unpaired parts have the corresponding zero-overlap bound, and the common kernel of $P,Q$ succeeds on the first call. The largest block thus gives \eqref{eq:detector-global}. If $c=0$, use a unit vector from each of the two nonzero orthogonal failure ranges; the same formula with $q=0$ applies.

With one live label, the same probability $r$ of choosing $A$ is used after every failure. A persistent shared seed would be extra memory and is not admitted. On a block, let $x,y$ be the remaining costs after $A,B$ fail. For $0<r<1$,
\[
 x=1+rx+(1-r)qy,\qquad y=1+rqx+(1-r)y.
\]
Solving and returning to the initial state gives $I+aP+bQ$ on that block, where
\[
 a=\frac{r/(1-r)+q}{1-q^2},\qquad
 b=\frac{(1-r)/r+q}{1-q^2}.
\]
The larger eigenvalue is $1+[a+b+\sqrt{q(a+b)^2+(1-q)(a-b)^2}]/2$. Since $a+b\ge2/(1-q)$, it is at least $1+(1+\sqrt q)/(1-q)=1+1/(1-\sqrt q)$, with equality at $r=1/2$. The deterministic endpoints have an undetectable failure state. The largest principal block proves \eqref{eq:detector-one}; the fair row realizes that operator bound in every block.

For the fixed-start unrestricted optimum, its first command row must be fair. On a largest principal block, both possible first failures have positive probability from the top vector of the fair cost operator. Every positive-probability repetition of the last failed detector strictly increases the expected cost there. Consequently an optimum needs a pure $B$ row after $A$ fails and a pure $A$ row after $B$ fails. These are two distinct rows, neither equal to the initial mixed row. Three live labels are necessary; those three rows and their obvious switching transitions suffice. This is a rigidity statement on the extremal block, not a claim about behavior on irrelevant physical subspaces.

The two-label parameter space is compact. For a sequence of proper two-label observers with uniformly bounded worst-state costs, every finite-prefix cost is continuous in the parameters and remains bounded at the limit. Monotone convergence gives a finite full cost there; unit costs imply proper operation. A minimizing sequence therefore has a proper minimizing limit. Such a limit cannot attain $V_*$ by the preceding row argument, proving the strict lower in \eqref{eq:detector-two}. Deterministic alternation from a fixed first detector uses two labels and has worst cost $1+(1-c^2)^{-1}$, proving the upper.

If the interface permits a randomized initial distribution, start fairly in the two pure switching labels. This stores the first choice in a counted label and attains $V_*$. One live label cannot attain it by \eqref{eq:detector-one}. The event/state counts follow from their exact halt semantics. Finally, a zero failure projection always succeeds. In finite dimension $\|PQ\|=1$ implies a nonzero vector in both failure ranges by the equality case for orthogonal projections. It remains there under every permitted command and is never detected. This proves the degenerate cases.
\end{proof}

\begin{corollary}[A primitive least-favorable preparation]\label{cor:detector-saddle}
Choose unit failure vectors $u,v$ in a largest principal-angle block, with $\langle u,v\rangle=c$, and put
\[
 \psi_*=(u+v)/\sqrt{2(1+c)},\qquad \rho_*=|\psi_*\rangle\langle\psi_*|.
\]
For $c=0$ choose any unit vectors in the two nonzero orthogonal failure ranges. This one pure state gives expected cost at least $V_*$ against every visible-history policy. Fair-first alternation has cost at most $V_*$ against every initial state. Thus it is an explicit minimax saddle without a policy-dependent preparation or a minimax exchange. The same preparation gives the one-live-label lower $V_1$ against every one-label row.
\end{corollary}
\begin{proof}
The first failure probability from $\rho_*$ is $(1+c)/2$ for either detector. Its post-failure state is respectively $u$ or $v$, whose least remaining cost is $1/(1-c^2)$ by the switching argument. Every first command mixture therefore costs at least $1+(1+c)/[2(1-c^2)]=V_*$. The uniform upper was proved above. For a one-label row, the cost at $\rho_*$ is $1+(a+b)(1+c)/2$. The bound $a+b\ge2/(1-c^2)$ gives $V_1$, attained by the fair row. This preparation is a witness for the stated detection task, not an automatically recurrent reference row after arbitrary rejected experiments.
\end{proof}

\subsection{The exact two-live-label frontier}
Put $q=c^2$. For $0<r\le1$ define
\begin{align}
 a_q(r)&=r+q,&b_q(r)&=(1-r)(r^{-1}+q),\label{eq:two-label-ab}\\
 F_q(r)&=1+\frac{a_q(r)+b_q(r)+\sqrt{(a_q(r)-b_q(r))^2+4q a_q(r)b_q(r)}}{2(1-q^2)}.
 \label{eq:two-label-F}
\end{align}

\begin{theorem}[Exact two-label optimum and a golden-ratio transition]\label{thm:two-label-exact}
In the nondegenerate class above,
\begin{equation}\label{eq:two-label-opt}
 V_2(c)=\min_{0<r\le1}F_{c^2}(r).
\end{equation}
A minimizing observer has a fixed initial label $S$ choosing $A$ with probability $r$ and $B$ otherwise, and a label $T$ choosing $B$ purely. After an $A$ failure it uses $T$; after a $B$ failure it uses $S$. Success commits immediately. It realizes the optimum for every finite-dimensional pair with the same largest overlap $c$.

Let $q_*=(\sqrt5-1)/2$. The minimizer is unique and is characterized by
\begin{equation}\label{eq:two-label-transition}
 \begin{array}{c|c|c}
 q&r_{\rm opt}&V_2\\\hline
 0&q_*&(1+\sqrt5)/2\\
 0<q<q_*&\text{unique }r\in(0,1)\text{ with }F'_q(r)=0&F_q(r_{\rm opt})\\
 q_*\le q<1&1&1+(1-q)^{-1}.
 \end{array}
\end{equation}
In the interior case, put $a=a_q(r)$, $b=b_q(r)$, $D=\sqrt{(a-b)^2+4qab}$ and $b'=-r^{-2}-q$. The unique equation is explicitly
\begin{equation}\label{eq:two-label-stationary}
 (1+b')D+(a-b)(1-b')+2q(b+ab')=0.
\end{equation}
Squaring does not license extraneous roots. Together with the one-label and unrestricted laws, this determines the exact fixed-start cost at every integer live-memory budget: $V_1$ at one, $V_2$ at two, and $V_*$ at three or more. State-stopped total counts add one terminal label; event-stopped counts do not.
\end{theorem}
\begin{proof}
First prove a converse against every two-label controller, including randomized post-failure transitions. Restrict the initial state to a largest principal block with failure states $u,v$ and squared overlap $q$. For $q=0$ use any unit vectors in the two nonzero orthogonal ranges. Relabel detector names if necessary so that the initial label has the larger probability $r$ of choosing $A$. The other label has probability $t$, with $0\le t\le r\le1$. Proper worst-state operation requires $r>0$ and $t<1$.

Even allowing the controller to choose whichever of its two labels it prefers after a failure, the least remaining costs at $u,v$ are
\begin{equation}\label{eq:two-row-values}
 x=\frac{(1-t)^{-1}+q/r}{1-q^2},\qquad
 y=\frac{r^{-1}+q/(1-t)}{1-q^2}.
\end{equation}
Here is a direct verification for arbitrary subsequent choices, not an assumed stationary optimum. One has $x-qy=(1-t)^{-1}$ and $y-qx=r^{-1}$. For any command probability $z\in[t,r]$,
\[
 1+zx+(1-z)qy-x=\frac{z-t}{1-t}\ge0,
 \qquad
 1+zqx+(1-z)y-y=1-\frac zr\ge0.
\]
These are lower potential inequalities for every allowed label and randomized label choice. Iteration to a finite horizon bounds actual expected calls below by the initial potential minus a remaining term bounded by $\max(x,y)$ times survival. The remainder vanishes for a proper policy. An improper policy has infinite call cost. Thus $x,y$ are universal lower bounds. They are attained by using the $t$-row after an $A$ failure and the $r$-row after a $B$ failure, as their equations verify.

From the prescribed initial label, the first call therefore has cost operator at least $I+rxP+(1-r)yQ$ on this block. For fixed $r$, both $x,y$ are nondecreasing in $t$, so this operator is minimized at $t=0$. Its coefficients then are
\[
 rx=\frac{r+q}{1-q^2},\qquad
 (1-r)y=\frac{(1-r)(r^{-1}+q)}{1-q^2}.
\]
Its largest eigenvalue is $F_q(r)$. This proves the lower without assuming deterministic transitions, a prior, or an uncounted initial seed.

For the matching upper use the displayed $S,T$ observer. Its exact block costs are \eqref{eq:two-row-values} at $t=0$. In a smaller principal block with squared overlap $q'\le q$, both nonnegative coefficients $a_{q'}(r)/(1-q'^2)$ and $b_{q'}(r)/(1-q'^2)$ increase with $q'$. The larger eigenvalue of $aP+bQ$ is nondecreasing in either nonnegative coefficient by positive matrix order, and in overlap by its square-root formula. Thus the worst block is the largest. Unpaired ranges have no greater than the corresponding zero-overlap bound, and the common kernel succeeds immediately. The same observer therefore has cost at most $F_q(r)$ for every initial state and attains it on the largest block. Its finite post-failure costs prove proper operation. Equality in \eqref{eq:two-label-opt} follows.

For $q=0$, $F_0(r)=1+\max\{r,r^{-1}-1\}$. The increasing and decreasing terms meet uniquely at $r^2+r-1=0$, giving the first line of \eqref{eq:two-label-transition}.

For $0<q<1$, use one fixed principal block. Its numerator matrix is $a_q(r)P+b_q(r)Q$, with $a_q$ affine and $b_q''(r)=2/r^3>0$. Its largest eigenvalue is strictly convex. Indeed the convex-combination matrix at two different parameters exceeds the matrix at the combined parameter by a positive multiple of $Q$. A top eigenvector of the latter has positive $Q$-expectation since $q>0$ and $a_q>0$. Evaluation at that vector proves strict convexity. The top eigenvalue is simple, so the function is differentiable, and it diverges as $r\downarrow0$.

At $r=1$ the numerator is $(1+q)P$. Its top vector is $u$, with $Q$-expectation $q$. Differentiation at this simple eigenvector gives
\[
 (1-q^2)F'_q(1)=1-q-q^2.
\]
For $q\ge q_*$ this is nonpositive, so strict convexity makes $r=1$ the unique minimum. For $q<q_*$ it is positive; divergence at zero and strict convexity give a unique interior stationary point. Differentiating the square root gives \eqref{eq:two-label-stationary}. This proves every case. The already proved live-role counts give the full integer frontier.
\end{proof}

The golden-ratio threshold concerns the squared principal-angle cosine and the optimal controller row, not a universal physical constant. At small overlap one continuing row is pure and the initial row is mixed. At and above the threshold the optimal two-label policy is fixed-start deterministic alternation. A third live label restores a fair initial choice without contaminating either continuing row and strictly improves the worst-state cost.

\subsection{A sharp memory-constrained minimax gap}
The exact finite-memory frontier also identifies a limitation of replacing a uniform converse by an optimized single-prior calculation. Let $\mathcal C_2$ be the proper fixed-start observers with at most two live labels in this detection interface.

\begin{corollary}[Minimax gap at two live labels]\label{cor:memory-minimax-gap}
For $P,Q\ne0$ and $c<1$,
\begin{equation}\label{eq:memory-minimax-gap}
 \inf_{\pi\in\mathcal C_2}\sup_\rho\E_\rho^\pi\tau
       =V_2(c)>V_*(c)
       =\sup_\rho\inf_{\pi\in\mathcal C_2}\E_\rho^\pi\tau.
\end{equation}
In particular no single initial preparation supplies the sharp two-label robust lower $V_2$. At zero overlap the gap is $(\sqrt5-2)/2$. Permitting a randomized initial distribution over the two counted pure labels removes this gap and gives robust value $V_*$. No analogous gap occurs at one live label or at three or more for this experiment.
\end{corollary}
\begin{proof}
The left equality and strict inequality are the memory theorems. For any fixed initial state, choose the better of the two deterministic first detectors and then alternate. Both choices are proper two-label observers. Their better cost is no larger than their average, which is the cost of fair-first alternation and is at most $V_*$. Hence the maximin is at most $V_*$. The single state $\rho_*$ in Corollary~\ref{cor:detector-saddle} gives cost at least $V_*$ against every policy, proving the reverse inequality. When $c=0$, $V_2=(1+\sqrt5)/2$ and $V_*=3/2$, giving the displayed gap. A randomized initial label realizes fair-first alternation directly. For one label, the same least-favorable state and fair independent row form a saddle at $V_1$; for three or more the full-history saddle applies.
\end{proof}

The policy in the inner maximin may depend on the fixed preparation; the robust policy on the left may not. The fixed-start two-label class is not closed under the convexification that would mix the two pure initial rows while remembering which continuing row to use. A free random initial distribution changes precisely that resource. Thus a minimax exchange would erase a real memory constraint. This is an exact result for the stated detector model, not a failure of the general sufficient transfer theorem or a claim that every finite-memory problem has such a gap.

\section{Raw task order and changing-rank stability}\label{sec:order}
The central theorem is formulated in measures because selection acts before normalization. This section gives a verification principle from primitive kernels and a singular extension of that principle. It does not assume a covariance or information matrix in the experiment.

\subsection*{A measure determined by executable affine-square tests}
For a report--audit law $P(dy,dz)$ on $Y\times H$ define the finite positive matrix measure
\begin{equation}\label{eq:moment-matrix}
 J_P(D)=\int_{D\times H}v(z)v(z)^\mathsf T P(dy,dz),
 \qquad v(z)=(1,z_1,\ldots,z_k)^\mathsf T.
\end{equation}
For finite positive matrix measures write $A\preccurlyeq B$ when $B(D)-A(D)$ is positive semidefinite for every Borel $D$. This definition is independent of the choice of dominating scalar measure.

\begin{proposition}[An exact order criterion for localized affine-square tests]\label{prop:matrix-order}
For two joint laws $P,P^0$ on the same full report space and $0<l\le1\le u$, the following assertions are equivalent:
\begin{equation}\label{eq:matrix-order}
 lJ_{P^0}\preccurlyeq J_P\preccurlyeq uJ_{P^0};
\end{equation}
for every Borel report event $D$ and every $a\in\R^{k+1}$,
\begin{equation}\label{eq:affine-tests}
 l\int_D|a^\mathsf T v(z)|^2dP^0
 \le\int_D|a^\mathsf T v(z)|^2dP
 \le u\int_D|a^\mathsf T v(z)|^2dP^0.
\end{equation}
It suffices to test rational vectors $a$. These equivalent conditions imply, for all Borel $q\in[0,1]$ and finite-image Borel $r:Y\to H$, both the report-mass and the squared-task inequalities with zero additive defect and these same factors. A violation of \eqref{eq:matrix-order} is witnessed by a single Borel report event and a rational affine-square test.
\end{proposition}
\begin{proof}
The integral in \eqref{eq:affine-tests} is $a^\mathsf T J_P(D)a$. The definition of positive semidefiniteness gives the equivalence. Continuity of the quadratic form gives sufficiency of rational vectors and a rational witness for a strict violation. Taking $a=(1,0,\ldots,0)$ gives mass order. For a fixed readout $r$, sum the $k$ quadratic tests with vectors $(-r_j,e_j)$, obtaining the order of $\|z-r\|^2$. Partition $Y$ into the finitely many cells where $r$ is constant. For a simple nonnegative $q$, apply order on intersections with these cells and add. Bounded monotone approximation extends to arbitrary $q\in[0,1]$. All measures are finite and all tested losses bounded.
\end{proof}

Thus the criterion is necessary and sufficient for the declared task-rich affine-square family. It is a sufficient condition for the narrower squared-prediction task of the central theorem. It is not claimed necessary for that narrower task, or for the existence of an arbitrary causal balancing strategy. The matrix is built from raw joint audit laws, not from an assumed parameterization. The upper construction may use its virtual report space after averaging its current command; the lower criterion must use the full report seen by a competitor, conditional on that competitor's chosen command.

\begin{theorem}[Order on measurable singular ranges]\label{thm:singular-order}
Let $J_0,J_1$ be finite positive matrix measures of size $k+1$. With
$\lambda=\tr J_0+\tr J_1$, take measurable positive semidefinite densities $A_0,A_1$ with respect to $\lambda$. Then
\begin{equation}\label{eq:singular-equivalence}
 lJ_0\preccurlyeq J_1\preccurlyeq uJ_0
\end{equation}
is equivalent, $\lambda$-almost everywhere, to equality of kernels
$\ker A_0=\ker A_1$ and to
\begin{equation}\label{eq:range-spectrum}
 lP_0\preccurlyeq A_0^{\dagger/2}A_1A_0^{\dagger/2}
            \preccurlyeq uP_0,
\end{equation}
where $P_0$ is the orthogonal projection onto $\operatorname{ran}A_0$ and $\dagger$ denotes the Moore--Penrose inverse. On the zero range both sides are zero. All these maps are measurable, even when rank changes. No positive lower bound on a nonzero eigenvalue of $A_0$ is required. In particular every zero affine-square task on an acquired report event remains exactly zero under the order.
\end{theorem}
\begin{proof}
A finite signed matrix measure is positive if and only if its density is positive semidefinite almost everywhere. To see the nontrivial direction, test the density against all rational vectors; a negative set of positive measure for one such quadratic form contradicts positivity of its integral on that set. The countable intersection of the resulting full-measure sets handles all rational vectors, hence all real vectors by continuity. Thus \eqref{eq:singular-equivalence} is equivalent to $lA_0\preccurlyeq A_1\preccurlyeq uA_0$ almost everywhere.

For two such matrices, a vector in $\ker A_0$ has zero quadratic form under $A_1$, and positivity implies that it is in $\ker A_1$. The lower inequality gives the converse. On their common range $A_0$ is positive definite; congruence by its inverse square root yields \eqref{eq:range-spectrum}. Conversely equality of kernels makes $A_1$ supported on that range. Congruence of \eqref{eq:range-spectrum} by $A_0^{1/2}$ gives the original two matrix inequalities, extended by zero on the common kernel.

Measurability follows from finite-dimensional spectral functional calculus: $P_0$ is the pointwise limit of $A_0(A_0+n^{-1}I)^{-1}$, and the scalar function equal to $x^{-1/2}$ on $x>0$ and zero at zero is Borel. Applying it spectrally gives a Borel matrix function on each rank stratum, hence on their finite measurable union. We use these matrices pointwise; their unbounded inverse norms are never taken as a uniform error constant. Finally a nonnegative task of zero integral under $J_0$ has integral at most $u$ times zero under $J_1$, proving preservation of zero tasks.
\end{proof}

There is an important distinction between this assertion and a rank formula for risk. Rank is allowed to change; the rate still comes from the acquired masses and coverings of Section~\ref{sec:geometry}. For example a nonzero matrix density with arbitrarily small positive eigenvalues can obey the order with fixed $l,u$, while an arbitrarily small absolute perturbation can create a new positive direction in its kernel and fail the order. This is precisely the difference relevant to a zero-loss stratum.

\begin{proposition}[Positive-kernel verification and relative command calibration]\label{prop:positive-order}
Suppose raw positive input states satisfy $ls_0\le s\le us_0$ in an ordered vector space and a fixed positive normalized report--audit instrument sends an input to a joint probability law. Its output laws and task matrices satisfy the same order. More generally let $\mathcal I_j$ be positive instruments and let a nominal probability vector $(\lambda_j^0)$ satisfy a desired joint-law identity after averaging. If
\[
 (1-\eta)\lambda_j^0\le\lambda_j\le(1+\eta)\lambda_j^0
 \quad\hbox{for every }j,\qquad 0\le\eta<1,
\]
then the actual averaged joint law lies between $1-\eta$ and $1+\eta$ times the nominal law, for every input. Successive relative input and instrument bounds multiply their lower factors and upper factors. They do not add an acquisition-amplified defect.
\end{proposition}
\begin{proof}
Positivity sends the two positive differences $s-ls_0$ and $us_0-s$ to positive measures. Integration of $v(z)v(z)^\mathsf T$ against a positive measure preserves positive matrix order. The command assertion follows by multiplying each nonnegative instrument measure by its nonnegative weight difference and summing. Applying the two inequalities successively multiplies the factors. These arguments do not divide by any report probability.
\end{proof}

\subsection*{Actual-mass-weighted predictive stability}
The following is the trace/base-norm contraction principle for positive normalized instruments; in the quantum case compare \cite[Corollary 3.40]{Watrous18}. We include its proof to retain the unnormalized and zero-mass conventions used here.
\begin{proposition}[Continuity through likelihood zeros and pure states]\label{prop:continuity}
A classical stochastic instrument on signed measures, or a finite-dimensional completely positive normalized quantum instrument, maps normalized inputs $s,t$ to unnormalized post-report states $A_y,B_y$. Use any common finite dominating measure for these two outputs and let $p_y=u(A_y)$, $q_y=u(B_y)$ be their masses. Then
\begin{align}
 \int\|A_y-B_y\|_1dy&\le\|s-t\|_1,\label{eq:stable1}\\
 \|p-q\|_{\TV}&\le\tfrac12\|s-t\|_1,\label{eq:stable2}\\
 \int\min(p_y,q_y)\left\|\frac{A_y}{p_y}-\frac{B_y}{q_y}\right\|_1dy
 &\le2\|s-t\|_1.\label{eq:stable3}
\end{align}
Terms of zero mass are zero. The same unnormalized and report-law inequalities hold for a complete stopped instrument under one common legal policy which terminates almost surely under both inputs.
\end{proposition}
\begin{proof}
Write $s-t=D_+-D_-$ with positive parts. Positivity, the triangle inequality and preservation of integrated mass give
$\int\|\mathcal I_y(D_+)-\mathcal I_y(D_-)\|_1dy\le u(D_+)+u(D_-)=\|s-t\|_1$.
Taking masses gives \eqref{eq:stable2}. If $p\le q$, then
\[
 p\|A/p-B/q\|_1\le\|A-B\|_1+(q-p).
\]
Interchanging the roles for $q\le p$, integrating, and using $|p-q|\le\|A-B\|_1$ proves \eqref{eq:stable3}. A finite causal continuation with a fixed policy is again a positive normalized instrument. Truncate a stopped continuation at a deterministic horizon, retaining a distinct surviving outcome. Apply the finite inequality, then let the horizon increase. The surviving masses vanish under both inputs by almost-sure termination, and the sum over terminal times converges monotonically. No expected-time assertion is deduced from these norm inequalities.
\end{proof}

The upper finite-state construction never stores a normalized posterior. It uses the acquired report and a finite decoder. Consequently stability of the raw instrument, rather than a globally contracting posterior recursion, suffices for the singular realizations below. More general closed-loop optimal-control claims would require their own strategy comparison and stability; they are not consequences of comparing the same fixed policy.

\section{Acquired geometry at regular, stratified and atomic scales}\label{sec:geometry}
All masses in this section refer to the actually prepared pushforward
$\sigma=(f,\sqrt w)_\#\nu$. They are not volumes on a parameter space. Conditional expectation of any feasible density given $(f,\sqrt w)$ preserves its mass, density cap and objective for every dictionary; conversely a density on this pushforward lifts by composition. Thus the profile can equivalently be optimized on this actual marked law. The transfer theorem compares operating task measures with this physical acquisition law. In the relative-order realizations it also bounds actual upper-program report measures above and below by fixed multiples of $\nu$. Thus the comparison remains tied to acquired mass, even though perturbed conditional means need not coincide with $f$.

\begin{proposition}[A non-power geometric comparison]\label{prop:sublevel}
For $r\ge0$ let
\[
 \Omega_m(r)=\max_{C\subset H,\ 1\le|C|\le m}
       \nu\{x:w(x)+\dist(f(x),C)^2\le r^2\}.
\]
For all $b,m\ge1$,
\begin{equation}\label{eq:sublevel-bounds}
 \sup_{r\ge0}r^2[1-b\Omega_m(r)]_+
 \le\Q_m(b)
 \le\inf\{r^2:\Omega_m(r)\ge1/b\}.
\end{equation}
In particular a verified estimate $b\Omega_m(r_-)\le1/2$ and a verified covering by $m$ distortion sublevel sets of total acquired mass at least $1/b$ at radius $r_+$ imply $r_-^2/2\le\Q_m(b)\le r_+^2$. The assertion requires no integer local dimension or isotropic geometry.
\end{proposition}
\begin{proof}
For a fixed dictionary and a density cap $b$, at most $b\Omega_m(r)$ of the selected mass lies in its radius-$r$ sublevel set. The remaining mass costs at least $r^2$, proving the lower. To see that the maximum defining $\Omega_m$ is attained, use an ordered tuple in $H^m$. For every report, the indicator of the closed sublevel condition is upper semicontinuous in that tuple. Reverse Fatou for these bounded indicators makes its integral upper semicontinuous. Compactness gives a maximizing tuple. If the corresponding sublevel set has mass $v\ge1/b$, the density $\one_D/v$ is feasible and has loss at most $r^2$. Taking the infimum over such radii gives the upper. This proof keeps the density cap and actual mass; a cover of the formal reachable set alone provides neither.
\end{proof}

The proposition can be used with anisotropic boxes, different local scales, or non-power small-ball laws, provided both the upper mass estimate and the required acquired cover are verified. It is a finite-scale statement, not a claim that an unspecified anisotropic dimension determines risk. We give explicit closed curves for a useful class and its finite stratifications.

Assume $\sigma$ is supported on $K\times[0,1]$, where $K\subset H$ is compact, and for some $d,\alpha>0$, $0<r_0\le1$ and $c,C>0$,
\begin{equation}\label{eq:rectangle}
 c r^d s^\alpha\le\sigma(B(x,r)\times[0,s])\le C r^d s^\alpha
 \quad(x\in K,\ 0<r\le r_0,\ 0<s\le1).
\end{equation}
These are global estimates on the stated component, including small reliability mass. Singular Ahlfors-regular location measures of noninteger $d$ are permitted.

\begin{proposition}[The marked acquired profile]\label{prop:regular}
Under \eqref{eq:rectangle}, uniformly for $b\ge1$ and integer $m\ge1$,
\begin{equation}\label{eq:F}
 \Q_m(b)\asymp F(b,m):=
 \max\{b^{-2/\alpha},(bm)^{-2/(d+\alpha)}\}.
\end{equation}
The constants depend only on the displayed data and the bounded task carrier. This geometric calculation is inherited from the stopped-experiment and predictable-geometry work retained in \cite{Companion}; its proof is included for independence of the central theorem.
\end{proposition}
\begin{proof}
Cover $K$ by finitely many balls of radius $r_0/2$. Their upper bounds yield $\sigma(K\times[0,s])\le C_0s^\alpha$. Under density cap $b$, choose $s=c_0b^{-1/\alpha}$ sufficiently small so that at most half the selected mass has reliability at most $s$. This gives a risk lower $s^2/2$.

For a dictionary with $m$ centers, its distortion sublevel at radius $r$ lies in the union of $m$ location balls of radius $r$ with reliability at most $r$. A center outside $K$ whose ball intersects $K$ can, for this mass bound, be replaced by one intersection point and the radius doubled. For $2r\le r_0$, the union has acquired mass at most $C_1m r^{d+\alpha}$. Set $r=c_1(bm)^{-1/(d+\alpha)}$ with one sufficiently small constant $c_1$ so that the restriction and the half-mass bound both hold for all $b,m\ge1$. Proposition~\ref{prop:sublevel} gives the other lower in \eqref{eq:F}.

For the upper, the location marginal has two-sided $d$-dimensional ball bounds by taking $s=1$. A maximal $4r$-separated set covers $K$ by $4r$-balls. The marginal upper bound then forces at least $c_2r^{-d}$ points for small $r$. Its radius-$r$ balls are disjoint. Choose constants $c_3>0$ small and $C_2$ large such that, for all sufficiently large $b$, the integer
\[
 j=\min\{m,\lfloor c_3b^{d/\alpha}\rfloor\}\ge1,
 \qquad r=C_2(bj)^{-1/(d+\alpha)}
\]
satisfies $r\le\min(r_0/8,1)$, $j\le c_2r^{-d}$, and $cj r^{d+\alpha}\ge1/b$. The last condition fixes a sufficiently large $C_2$; after fixing it, a sufficiently small $c_3$ enforces the packing condition because
$j r^d=C_2^d j^{\alpha/(d+\alpha)}b^{-d/(d+\alpha)}$.
Increasing the lower threshold for $b$ then enforces the radius condition and $j\ge1$. Use $j$ separated centers. The disjoint rectangles $B(x_i,r)\times[0,r]$ have union mass at least $1/b$ and distortion to those centers at most $2r^2$. Selecting a unit mass under cap $b$ from that union proves the upper, and substitution of $j$ gives \eqref{eq:F}. For $b$ in the remaining bounded interval, always accept and use one center. Its bounded loss is controlled by a fixed multiple of $b^{-2/\alpha}$, uniformly over arbitrarily large $m$. This proves all claimed uniformities.
\end{proof}

\begin{corollary}[All-scale resource curves under fixed relative calibration]\label{cor:uniform}
Suppose \eqref{eq:rectangle} and the relative hypotheses of Corollary~\ref{cor:relative} hold. On a set of feasible budgets suppose the positive effort ratio $\kappa_+/\kappa_-$ and positive width ratio $m_-/m_+$ are bounded above and below by fixed constants. Then
\begin{equation}\label{eq:uniform-law}
 R_i^\Theta(N,M)\asymp
 F(N/\kappa_-,m_-(M))
\end{equation}
throughout that set, including arbitrarily fine resolution. The constants may depend on $l_\pm,u_\pm$, but these factors need not approach one as $N,M$ grow.

More generally the same conclusion holds with absolute defects if $A_N\le A_0<1$, $b_+a_+\le A_1<\infty$, $E_N\le c_0F(b_-,m_-)$ and $b_+e_+\le C_0F(b_-,m_-)$, for a sufficiently small $c_0>0$ depending on the fixed geometric and order constants.
\end{corollary}
\begin{proof}
For pure relative order and $l_+b_+\ge1$, both effective budgets and widths in \eqref{eq:relative} differ from $b_-,m_-$ by bounded positive factors. Apply Proposition~\ref{prop:regular}. If $1\le b_+<1/l_+$, then $b_-$ lies in a bounded interval. The first-attempt upper is bounded, while the lower and $F(b_-,m_-)$ are bounded below by a fixed positive constant independently of $m_-$. This treats the boundary regime. With absolute defects, $t_N$ is bounded below by $(1-A_0)/u_-$ and the corrected upper budget differs by fixed factors. The geometric lower therefore exceeds $cF$ before subtracting $E_N$. Take $c_0\le c/2$. The upper defect has the stated order, yielding the result. A fixed positive absolute defect need not satisfy these scale conditions; relative order and absolute proximity are distinct hypotheses.
\end{proof}

\begin{theorem}[A mixed-dimensional acquired law]\label{thm:stratified}
Suppose the acquired marked law decomposes as a finite probability mixture
$\sigma=\sum_{j=1}^J\lambda_j\sigma_j$, with $\lambda_j>0$, $\sum_j\lambda_j=1$. The components are distributions of the actual task coordinates, not the pushforwards of unspecified hidden conditional means. Each $\sigma_j$ satisfies \eqref{eq:rectangle} with its own $d_j,\alpha_j$ and fixed geometric constants. Define the feasible acquired-mass simplex
\[
 \mathcal S_b=\{(s_j):s_j\ge0,\ \sum_js_j=1,\ s_j\le b\lambda_j\}.
\]
For $m\ge2J$ and all $b\ge1$,
\begin{equation}\label{eq:stratified-law}
 \Q_m(b)\asymp\min_{s\in\mathcal S_b}
 \sum_{j:s_j>0}s_j
 \max\left\{
 \left(\frac{b\lambda_j}{s_j}\right)^{-2/\alpha_j},
 \left(\frac{b\lambda_jm}{s_j}\right)^{-2/(d_j+\alpha_j)}\right\}.
\end{equation}
A zero $s_j$ contributes zero. The constants depend on the finite component data and $J$, not on $b,m$. Applying Theorem~\ref{thm:transfer} to this profile gives the corresponding two-sided causal resource law; no stratum is renormalized without retaining its selected mass $s_j$.
\end{theorem}
\begin{proof}
For any feasible density and dictionary, set $s_j=\lambda_j\int h\,d\sigma_j$. The masses belong to $\mathcal S_b$. When $s_j>0$, the normalized density $\lambda_jh/s_j$ on the $j$th component is bounded by $b\lambda_j/s_j$. Its dictionary still has at most $m$ elements. Thus the full objective is at least
\[
 \sum_{j:s_j>0}s_j\Q_m^j(b\lambda_j/s_j).
\]
This proves the lower after Proposition~\ref{prop:regular}.

Conversely fix feasible $s$ and allocate $m_j=\lfloor m/J\rfloor$ centers to every component with positive mass. Their total is at most $m$ and $m_j\ge m/(2J)$. Choose an optimal component density $h_j$ at cap $b\lambda_j/s_j$ and dictionary $C_j$. The measure $\zeta=\sum_{s_j>0}s_jh_j\sigma_j$ has mass one and is bounded by $b\sigma$. Use its Radon--Nikodym density with respect to $\sigma$ and the union of the $C_j$. Nearest-center distortion to that union is no greater than distortion to the component's own dictionary. This gives the upper $\sum_js_j\Q_{m_j}^j(b\lambda_j/s_j)$, within a fixed factor of the expression in \eqref{eq:stratified-law}. This is a measurable existence construction, not automatically a finitely computable classifier. If a component label is not observed, this construction still uses only the density of $\zeta$ on the actual marked coordinates; no latent component choice is made by the observer. The displayed objective is continuous at $s_j=0$ by its positive powers of $s_j$, and the feasible simplex is compact and nonempty. Its minimum is attained. Taking this minimum completes the proof.
\end{proof}

The mixture theorem permits different local dimensions and reliability exponents, including noninteger singular dimensions. It is more than a statement that strata may be supplied through an interface: it gives their coupled acquisition--memory optimization and proves both sides with the original masses. Proposition~\ref{prop:sublevel} remains available when an anisotropic component has no single power-law description.

\begin{proposition}[An acquired zero-loss atom]\label{prop:atom}
Suppose $\sigma=\lambda\delta_{(x_*,0)}+(1-\lambda)\sigma_0$, $0<\lambda<1$, where $\sigma_0$ has no mass at that atom. Then $\Q_m(b)=0$ for every $m\ge1$ if $b\lambda\ge1$. If $s=1-b\lambda>0$, then
\begin{equation}\label{eq:atom}
 s\Q_m^0\left(\frac{b(1-\lambda)}s\right)
 \le\Q_m(b)\le
 s\Q_{m-1}^0\left(\frac{b(1-\lambda)}s\right)\quad(m\ge2).
\end{equation}
For $m=1$ only the displayed lower applies; a valid upper fixes its center at $x_*$ and optimizes the remaining selected mass with that center. If $w>0$ almost surely on the remainder, then every finite-width profile is positive below threshold. Under pure relative transfer and the additional condition $w>0$ almost surely on the remainder, zero risk is impossible when $u_-b_-\lambda<1$. Without that additional condition, the displayed atom still supplies the sufficient zero-risk bound when $l_+b_+\lambda\ge1$ and the upper width is feasible, but other zero-loss atoms may lower the true threshold. If $w>0$ almost surely on the remainder, the exact matched threshold is $N=\kappa/\lambda$. The complete finite-width threshold without this restriction is Theorem~\ref{thm:zero-spectrum}.
\end{proposition}
\begin{proof}
Above threshold accept only the atom and decode $x_*$. Below threshold at least mass $s$ must come from the remainder. Scale any selected remainder measure down to mass $s$; it remains dominated by $b(1-\lambda)\sigma_0$ and nonnegative loss cannot increase. Normalize it to obtain the lower in \eqref{eq:atom}. For the upper retain all available atom mass $b\lambda$, allocate one center to it, and optimize the mass-$s$ remainder with $m-1$ centers. With one center there is no zero-center optimization, which explains the separate statement. If the remainder has $w>0$ almost surely, any positive selected probability has positive residual loss; attainment from Lemma~\ref{lem:attainment} makes the profile minimum strictly positive. Finally substitute the atom criterion into the two effective caps of \eqref{eq:relative}. Relative calibration preserves zero tasks but may change the required acquisition threshold by the stated factors; it does not preserve the numerical threshold without exact matching.
\end{proof}

\subsection*{Nonuniform defects and metastable preparation}
\begin{lemma}[Killed-potential occupation bound]\label{lem:drift}
Let a nonnegative predictable defect $d_i$ be charged at raw call $i$, and suppose, for every legal command before commitment, nonnegative integrable potentials obey
\[
 d_i+\E_*[V_i\one_{\{i<\tau\}}\mid\mathcal F_{i-1},a_i]
 \le V_{i-1}+\gamma\E_*[c_i\mid\mathcal F_{i-1},a_i].
\]
Here $V_0$ is integrable, $\gamma\ge0$, and after commitment the potential is killed. Then every budget-$N$ competitor satisfies $\E_*\sum_{i\le\tau}d_i\le\E_*V_0+\gamma N$. Potentials may depend on the hidden physical state; this is an analytic certificate, not uncharged observer memory.
\end{lemma}
\begin{proof}
Multiply by the predictable indicator of making call $i$, integrate and sum to a deterministic horizon $n$. All interior potentials cancel. The remaining final potential is nonnegative and may be dropped. Bound the expected partial cost by $N$, then use monotone convergence for the nonnegative defect sum. Randomized commands are handled by integrating the commandwise inequality. No optional stopping with unverified uniform integrability is needed.
\end{proof}

\begin{corollary}[Relative order with rare defective modes]\label{cor:metastable}
Suppose relative order holds on good raw histories, while visible conditional additive defects are bounded by the conditional expectations of $a_b\one_{\{D_{i-1}=1\}}$ and $e_b\one_{\{D_{i-1}=1\}}$ for a hidden mode $D\in\{0,1\}$. At every legal call, given the physical past, assume
\[
 \Prb_*(D_i=1\mid D_{i-1}=0,\cdot)\le u,
 \qquad\Prb_*(D_i=0\mid D_{i-1}=1,\cdot)\ge v>0.
\]
Let $\pi_0=\Prb_*(D_0=1)$ be the actual prepared mass and $c_i\ge c_*>0$. The central theorem applies with
\[
 A_N=\frac{a_b}{v}\pi_0+\frac{a_bu}{vc_*}N,
 \qquad E_N=\frac{e_b}{v}\pi_0+\frac{e_bu}{vc_*}N.
\]
These quantities may replace rowwise worst-case defects even if bad-mode errors are not small and the residence time $1/v$ is large.
\end{corollary}
\begin{proof}
Use $V(D)=(a_b/v)\one_{\{D=1\}}$. In a good mode its expected increase is at most $a_bu/v$; in a bad mode its expected decrease is at least $a_b$ and pays for the defect. Killing at commitment only improves the inequality. Since each raw call costs at least $c_*$, Lemma~\ref{lem:drift} gives the first bound. Counting only productive calls makes it smaller. Replace $a_b$ by $e_b$ for the task bound and use conditional expectation to pass from physical modes to the visible-history defects. This retains actual preparation and selection-biased residence, rather than substituting an invariant-mode distribution.
\end{proof}

\subsection{An explicit six-regime allocation diagram}\label{sec:strata-phase}
Here we solve the variational comparator in Theorem~\ref{thm:stratified}, rather than claim that its implicit comparison constants are one. Take two components with $(d_1,\alpha_1,\lambda_1)=(1,1,1/2)$ and $(d_2,\alpha_2,\lambda_2)=(2,2,1/2)$, and total dictionary width $m=4$. If $s$ is the selected mass from the first component, the comparator is
\[
 J_b(s)=\max\left\{\frac{4s^3}{b^2},\frac{s^2}{2b}\right\}
       +\max\left\{\frac{2(1-s)^2}{b},
                     \frac{(1-s)^{3/2}}{\sqrt{2b}}\right\},
\]
with $\max(0,1-b/2)\le s\le\min(1,b/2)$ and $b\ge1$. Its minimum describes the risk up to the inherited component-dependent constants. It is not an exact formula for every acquired law having these exponents.

These component hypotheses can arise from an actual bounded preparation: take the first location uniform on a unit segment in $\R^2$ and the second uniform on a unit square, with reliability variables having distribution functions $t$ and $t^2$, respectively. With a fresh fair sign $E$, audit $Z=Y+TEe_1$ and report $(Y,T)$. Then $f=Y$ and $w=T^2$ in a common compact convex carrier. The mixture has the stated acquired dimensions and exponents. The diagram still concerns the theorem's comparator rather than replacing its comparison constants by exact ones.

\begin{proposition}[Exact selected-mass phases for the comparator]\label{prop:strata-phases}
Put
\[
 b_1=6-2\sqrt6,\quad b_2=(16\sqrt6-24)/5,\quad
 b_3=8(\sqrt2-1),\quad b_4=32/5,\quad b_5=36/5.
\]
The minimizer is unique and continuous in $b$. It is
\begin{equation}\label{eq:strata-phases}
 s_*(b)=
 \begin{cases}
 1-b/2,&1\le b\le b_1,\\
 (\sqrt{b^2+12b}-b)/6,&b_1\le b\le b_2,\\
 1-b/8,&b_2\le b\le b_3,\\
 \text{the unique }s\in(0,1)\text{ with }128s^4=b^3(1-s),
       &b_3\le b\le b_4,\\
 b/8,&b_4\le b\le b_5,\\
 (\sqrt{81b^2+288b}-9b)/16,&b\ge b_5.
 \end{cases}
\end{equation}
In particular there are two genuine intervals where the minimizer lies at a nonsmooth crossover, not just isolated changes between differentiable stationary formulas. The initially decreasing allocation to the first component eventually increases to one.
\end{proposition}
\begin{proof}
Each summand is strictly convex on its nonnegative mass interval, and hence $J_b$ is strictly convex. Both components change from their location-limited term to their reliability-limited term when their own selected mass crosses $b/8$. On a differentiable interval the minimizing equation is equality of their marginal costs.

When both are reliability-limited, this equation is
$12s^2/b^2=4(1-s)/b$, or $3s^2=b(1-s)$. Its solution first enters the feasible interval at $s=1-b/2$, giving $b_1$, and reaches the second component's crossover $1-s=b/8$ at $b_2$. At that crossover the second marginal-cost interval is $[3/8,1/2]$. The first marginal cost $12(1-b/8)^2/b^2$ decreases through that interval as $b$ runs from $b_2$ to $b_3$. Thus $s=1-b/8$ is the minimizer throughout that interval.

With the first component reliability-limited and the second location-limited, marginal equality is
$12s^2/b^2=3\sqrt{1-s}/(2\sqrt{2b})$, equivalent to the displayed quartic. Its left side is strictly increasing in $s$ and its right side strictly decreasing. The root reaches $s=b/8$ at $b_4$. At this first-component crossover its own marginal-cost interval is $[1/8,3/16]$. The second marginal cost $3\sqrt{1-b/8}/(2\sqrt{2b})$ decreases from $3/16$ to $1/8$ between $b_4$ and $b_5$, giving the second nonsmooth interval. Thereafter both terms are location-limited; equality becomes $8s^2=9b(1-s)$, with the final displayed root. Substitution verifies the active conditions and the original mass caps in each interval. Strict convexity proves uniqueness; matching endpoint values proves continuity. The last equation implies $1-s_*(b)\sim8/(9b)$.
\end{proof}

\begin{figure}[htbp]
\centering
\setlength{\unitlength}{1pt}
\begin{picture}(410,210)(0,-14)
\put(38,22){\line(1,0){350}}
\put(38,22){\line(0,1){145}}
\put(0,182){$s_*(b)$}
\put(34,22.000){\line(1,0){4}}
\put(8,19.000){\makebox(23,8)[r]{$ 0 $}}
\put(34,93.000){\line(1,0){4}}
\put(8,90.000){\makebox(23,8)[r]{$ 0.5 $}}
\put(34,164.000){\line(1,0){4}}
\put(8,161.000){\makebox(23,8)[r]{$ 1 $}}
\put(38.000,18){\line(0,1){4}}
\put(30.000,3){\makebox(16,9){$ 1 $}}
\put(109.328,18){\line(0,1){4}}
\put(101.328,3){\makebox(16,9){$ 2 $}}
\put(203.619,18){\line(0,1){4}}
\put(195.619,3){\makebox(16,9){$ 5 $}}
\put(274.947,18){\line(0,1){4}}
\put(266.947,3){\makebox(16,9){$ 10 $}}
\put(388.000,18){\line(0,1){4}}
\put(380.000,3){\makebox(16,9){$ 30 $}}
\put(174,-11){$b$ (logarithmic scale)}
\thinlines
\put(47.903,22){\line(0,1){142}}
\put(40.903,170.000){\makebox(14,8){$b_1$}}
\thinlines
\put(152.360,22){\line(0,1){142}}
\put(145.360,178.000){\makebox(14,8){$b_2$}}
\thinlines
\put(161.287,22){\line(0,1){142}}
\put(154.287,170.000){\makebox(14,8){$b_3$}}
\thinlines
\put(229.022,22){\line(0,1){142}}
\put(222.022,178.000){\makebox(14,8){$b_4$}}
\thinlines
\put(241.143,22){\line(0,1){142}}
\put(234.143,170.000){\makebox(14,8){$b_5$}}
\thicklines
\qbezier(38.000,93.000)(40.734,91.062)(43.469,89.125)
\qbezier(43.469,89.125)(45.686,87.476)(47.903,85.828)
\qbezier(47.903,85.828)(48.420,85.942)(48.938,86.055)
\qbezier(48.938,86.055)(51.672,86.661)(54.406,87.267)
\qbezier(54.406,87.267)(57.141,87.878)(59.875,88.489)
\qbezier(59.875,88.489)(62.609,89.105)(65.344,89.720)
\qbezier(65.344,89.720)(68.078,90.341)(70.812,90.961)
\qbezier(70.812,90.961)(73.547,91.585)(76.281,92.210)
\qbezier(76.281,92.210)(79.016,92.838)(81.750,93.466)
\qbezier(81.750,93.466)(84.484,94.098)(87.219,94.730)
\qbezier(87.219,94.730)(89.953,95.365)(92.688,96.001)
\qbezier(92.688,96.001)(95.422,96.639)(98.156,97.277)
\qbezier(98.156,97.277)(100.891,97.918)(103.625,98.558)
\qbezier(103.625,98.558)(106.359,99.201)(109.094,99.844)
\qbezier(109.094,99.844)(111.828,100.488)(114.562,101.133)
\qbezier(114.562,101.133)(117.297,101.779)(120.031,102.424)
\qbezier(120.031,102.424)(122.766,103.071)(125.500,103.718)
\qbezier(125.500,103.718)(128.234,104.365)(130.969,105.013)
\qbezier(130.969,105.013)(133.703,105.660)(136.438,106.307)
\qbezier(136.438,106.307)(139.172,106.954)(141.906,107.601)
\qbezier(141.906,107.601)(144.641,108.247)(147.375,108.893)
\qbezier(147.375,108.893)(149.868,109.481)(152.360,110.069)
\qbezier(152.360,110.069)(152.602,109.942)(152.844,109.815)
\qbezier(152.844,109.815)(155.578,108.336)(158.312,106.858)
\qbezier(158.312,106.858)(159.800,106.020)(161.287,105.182)
\qbezier(161.287,105.182)(162.534,105.742)(163.781,106.301)
\qbezier(163.781,106.301)(166.516,107.538)(169.250,108.774)
\qbezier(169.250,108.774)(171.984,110.020)(174.719,111.266)
\qbezier(174.719,111.266)(177.453,112.519)(180.187,113.772)
\qbezier(180.187,113.772)(182.922,115.029)(185.656,116.286)
\qbezier(185.656,116.286)(188.391,117.543)(191.125,118.800)
\qbezier(191.125,118.800)(193.859,120.055)(196.594,121.309)
\qbezier(196.594,121.309)(199.328,122.556)(202.062,123.803)
\qbezier(202.062,123.803)(204.797,125.039)(207.531,126.275)
\qbezier(207.531,126.275)(210.266,127.496)(213.000,128.716)
\qbezier(213.000,128.716)(215.734,129.916)(218.469,131.116)
\qbezier(218.469,131.116)(221.203,132.292)(223.938,133.467)
\qbezier(223.938,133.467)(226.480,134.534)(229.022,135.600)
\qbezier(229.022,135.600)(229.214,135.812)(229.406,136.025)
\qbezier(229.406,136.025)(232.141,139.137)(234.875,142.248)
\qbezier(234.875,142.248)(237.609,145.530)(240.344,148.812)
\qbezier(240.344,148.812)(240.743,149.306)(241.143,149.800)
\qbezier(241.143,149.800)(243.478,150.060)(245.812,150.319)
\qbezier(245.812,150.319)(248.547,150.614)(251.281,150.908)
\qbezier(251.281,150.908)(254.016,151.192)(256.750,151.476)
\qbezier(256.750,151.476)(259.484,151.750)(262.219,152.024)
\qbezier(262.219,152.024)(264.953,152.287)(267.688,152.551)
\qbezier(267.688,152.551)(270.422,152.805)(273.156,153.059)
\qbezier(273.156,153.059)(275.891,153.303)(278.625,153.548)
\qbezier(278.625,153.548)(281.359,153.783)(284.094,154.018)
\qbezier(284.094,154.018)(286.828,154.244)(289.562,154.470)
\qbezier(289.562,154.470)(292.297,154.687)(295.031,154.904)
\qbezier(295.031,154.904)(297.766,155.112)(300.500,155.320)
\qbezier(300.500,155.320)(303.234,155.520)(305.969,155.720)
\qbezier(305.969,155.720)(308.703,155.912)(311.437,156.103)
\qbezier(311.437,156.103)(314.172,156.287)(316.906,156.471)
\qbezier(316.906,156.471)(319.641,156.647)(322.375,156.823)
\qbezier(322.375,156.823)(325.109,156.992)(327.844,157.161)
\qbezier(327.844,157.161)(330.578,157.322)(333.312,157.484)
\qbezier(333.312,157.484)(336.047,157.638)(338.781,157.793)
\qbezier(338.781,157.793)(341.516,157.940)(344.250,158.088)
\qbezier(344.250,158.088)(346.984,158.230)(349.719,158.371)
\qbezier(349.719,158.371)(352.453,158.506)(355.188,158.641)
\qbezier(355.188,158.641)(357.922,158.770)(360.656,158.899)
\qbezier(360.656,158.899)(363.391,159.022)(366.125,159.146)
\qbezier(366.125,159.146)(368.859,159.263)(371.594,159.381)
\qbezier(371.594,159.381)(374.328,159.494)(377.062,159.606)
\qbezier(377.062,159.606)(379.797,159.713)(382.531,159.820)
\qbezier(382.531,159.820)(385.266,159.923)(388.000,160.025)
\end{picture}

\caption{The exact minimizer of the displayed comparator, not a fitted risk curve. Vertical lines mark the five algebraic transition budgets. The two crossover intervals are part of the optimum; smoothing a maximum before differentiation would miss them. The plotted values are checked against independent bounded scalar minimization.}\label{fig:strata-phase}
\end{figure}

\section{Classical and quantum realizations from raw kernels}\label{sec:realizations}
The two realizations below are verified after the general theorem. Recurrent witness preparation, when present, belongs to the charged productive primitive; it is not a reset made freely available to other models. Neither supplies a premise used to prove it. Their conditional physical evolution between calls may be noncontracting, and their acquired laws may be singular. Both tasks score the post-command audit. A target which must remain unchanged is a different experiment; the nonintervening directional result is retained in the companion.

\subsection*{A common paid physical interface}
The gate port $G$ costs one per call and has success probability at least $p_0>0$, conditional on every physical and visible past. A success enables exactly one productive call at port $A$, also of cost one. Commitment is forbidden at $G$ and allowed after a productive report. Rejection at $A$ returns the protocol to $G$; it need not reset the hidden state. The two ports have disjoint physical command domains and no port tag is a free intercall input to an autonomous observer. Suppose the actual lower witness has gate success probability at most $p_1\le1$. Then
\begin{equation}\label{eq:gate-constants}
 \kappa_-=1+1/p_1,\qquad \kappa_+=1+1/p_0,
 \qquad B_0=B=2.
\end{equation}
Let $G_\tau$ and $A_\tau$ count the gate and productive calls through commitment. Every productive call consumes a gate success. A terminating path has no unconsumed success. Predictable counting gives $\E_*A_\tau\le p_1\E_*G_\tau$, so $\E_*C_\tau\ge\kappa_-\E_*A_\tau$. The probability of $j$ failures after any gate start is at most $(1-p_0)^j$, giving conditional expected gate count at most $1/p_0$ and whole-attempt cost at most $\kappa_+$. Both mandatory ports are visited before commitment, and Lemma~\ref{lem:roles} gives two live roles. Consequently the construction requires at least two event labels or three state labels when it has one output. If the witness has constant success $p_0=p_1$, then $\E_*C_\tau=\kappa\E_*A_\tau$ with $\kappa=1+1/p_0$. In that matched case budget feasibility is exactly $N\ge\kappa$ with the respective state thresholds. Even when $p_0=1$, the cost is two, not one.

\begin{theorem}[A relatively calibrated classical realization]\label{thm:classical}
Let a hidden symbol $E$ take values in a finite set of size $q$. A finite group $\mathcal G$ acts transitively on that set. The allowed productive commands apply $g\in\mathcal G$ and then, within the same call, a fixed positive normalized joint report--audit kernel $J_{gE}(dy,dz)$ with bounded audit in $H$. There is no unknown intervening operation between that command and its audit instrument. Evolution between port interactions is otherwise arbitrary and causal.

The genuinely prepared comparison has uniform input and joint law $K_0=q^{-1}\sum_eJ_e$. Its full report is its mark. Suppose the operating family contains a single witness whose conditional input weights, after every reached pre-command history, satisfy
\[
 (1-\eta_r)/q\le\Prb_*(E=e\mid h)\le(1+\eta_r)/q,
 \qquad 0\le\eta_r<1.
\]
Suppose one parameter-independent implementation chooses its command independently of the hidden current symbol conditional on its visible past, with conditional weights
\[
 (1-\eta_u)/|\mathcal G|\le\lambda_g\le(1+\eta_u)/|\mathcal G|,
 \qquad0\le\eta_u<1,
\]
uniformly over all models and surviving histories. Use the gate interface above. Then Theorem~\ref{thm:transfer} holds with identity exports, zero additive defects, and
\begin{equation}\label{eq:classical-relative}
 l_-=1-\eta_r,\quad u_-=1+\eta_r,
 \qquad l_+=1-\eta_u,\quad u_+=1+\eta_u.
\end{equation}
The comparison need not be an operating member. Fixed nonzero $\eta_r,\eta_u<1$ preserve the all-scale resource curve of Corollary~\ref{cor:uniform} whenever its acquired geometry holds.
\end{theorem}
\begin{proof}
A chosen permutation preserves the two componentwise bounds on the witness weights. Multiplying by each positive $J_e$ and summing gives the full command-conditioned joint-law order
$(1-\eta_r)K_0\le P_{h,g}\le(1+\eta_r)K_0$.
Take $P^0_{h,g}=K_0$ and the identity mark. Conditional moments given its full report are exactly \eqref{eq:derived}, so all lower certificates follow from Proposition~\ref{prop:matrix-order}. This holds for every chosen command and every competing policy, not merely for a randomizing policy.

For any fixed incoming symbol, uniform averaging over a transitive group gives the uniform distribution of its image: each image has the same number of preimages in the group. Averaging also over any incoming distribution preserves that fact. The actual componentwise command bounds and positivity therefore put the full joint upper-program law between $1-\eta_u$ and $1+\eta_u$ times $K_0$. Independence of the command from the hidden current input is needed for this conditional averaging; it is stated as a physical implementation property, not inferred from marginal coin frequencies. The upper accepts on the resulting report, not on a discarded group choice. Exports are the identity because the audit is the post-command audit. Group selection is within-call work, and only the gate or productive routing label or final decoder value persists. The gate proof establishes the same-witness cost lower, common cost upper and both live roles. The general theorem now applies.
\end{proof}

For a concrete acquired geometry take $q=2$, with identity and sign reversal. Independently draw $(Y,T)$ from a fixed law on $K\times[0,1]$, report these coordinates, and audit
\begin{equation}\label{eq:classical-mark}
 Z=Y+TEe_1
\end{equation}
in a sufficiently large compact convex carrier. The uniform preparation gives $f=Y,w=T^2$. Choose a $d$-regular location measure and $\Prb(T\le s)=s^\alpha$ to verify \eqref{eq:rectangle} directly; choose a finite mixture of such laws to verify Theorem~\ref{thm:stratified}; add a physically prepared point mass at $(x_*,0)$ to obtain Proposition~\ref{prop:atom}. A biased fresh witness and imperfect command probabilities satisfy \eqref{eq:classical-relative} with fixed nonzero constants. On $T=0$ the task is exactly deterministic even under these errors. On $T>0$ the two-point audit moment matrix has rank two, while at $T=0$ it has rank one. This is an actual changing-rank task measure, not a formal limiting parameter set.

The command group and its bit-level sampler are part of the program description and workspace account. A high-cardinality group does not provide free randomization or a free persistent seed. In the binary realization a single fresh sign coin suffices within a productive call. Conditional posterior continuity of the classical raw kernels is governed by Proposition~\ref{prop:continuity}, including likelihood zeros.

\begin{theorem}[A relatively calibrated noncommuting quantum realization]\label{thm:quantum}
Let the productive physical system have finite dimension $d\ge2$. Its available commands include the $d^2$ Weyl conjugations, followed within the same productive call by a fixed positive normalized quantum report--audit instrument $\mathcal J$, with classical joint law $\mathcal J(\rho)(dy,dz)$ and $z\in H$. The physical task is again the post-command audit. Put $\omega=I/d$ and $K_0=\mathcal J(\omega)$.

Suppose one actual witness obeys, conditionally before every selected productive command,
\[
 (1-\eta_r)\omega\preccurlyeq\rho\preccurlyeq(1+\eta_r)\omega,
 \qquad0\le\eta_r<1.
\]
Suppose the common implementation uses current Weyl weights between $(1-\eta_u)/d^2$ and $(1+\eta_u)/d^2$, independently of the current hidden input conditional on the past, with $0\le\eta_u<1$. With the same gate certificates, all conclusions of Theorem~\ref{thm:classical} hold with the constants \eqref{eq:classical-relative}. The operating witness need not be exactly maximally mixed. Incoming states for the upper may be pure, and post-report ranks may change on positive acquired mass. No common eigenbasis or normalized-filter contraction is assumed.
\end{theorem}
\begin{proof}
Let $W_{p,q}$ be the shift by $q$ followed by the phase $\exp(2\pi ipj/d)$ on the $j$th basis vector. On a matrix unit $|j\rangle\langle k|$, averaging the phase kills off-diagonal terms; averaging shifts makes the diagonal uniform. Hence
\[
 d^{-2}\sum_{p,q}W_{p,q}\rho W_{p,q}^*=I/d
\]
for every density matrix. This elementary randomization identity is standard \cite{AMTW}; we use its complete positive-instrument consequence, not a new state-randomization claim.

Each chosen conjugation fixes $\omega$ and preserves the witness Loewner order. Positivity of the full report--audit instrument therefore gives lower joint-law order relative to $K_0$ with factors $1\pm\eta_r$. Positivity means that the image of every positive operator is a positive classical measure, including both report and future audit. It consequently compares the full observable report and task, not just the report marginal.

For the upper, the componentwise weight bounds give
\[
 (1-\eta_u)\omega
 \preccurlyeq\sum_{p,q}\lambda_{p,q}W_{p,q}\rho W_{p,q}^*
 \preccurlyeq(1+\eta_u)\omega
\]
for every incoming state, including pure states. Apply the same positive instrument. Proposition~\ref{prop:matrix-order} supplies the task certificates, and identity export, counted live roles and the gate costs finish the verification. The lower witness has a relative input bound around $\omega$; we do not describe that hypothesis as absence of an input-rank condition. The absence of a floor concerns the post-report states and task strata, where the following example has genuine rank changes.
\end{proof}

Here is a raw qubit instrument giving the same acquired geometry as \eqref{eq:classical-mark} but noncommuting predictive states. Draw $(y,t)$ with any stated law $\nu_0$ on $K\times[0,1]$, choose an equatorial unit vector $n=(\cos\phi,\sin\phi,0)$ and $0\le\gamma(y,t)\le1$, and set
\[
 E_v=(I+v\gamma n\cdot\sigma)/2,\qquad L_v=\sqrt{E_v},\qquad v=\pm1.
\]
The report is $(y,t,v)$ and the unnormalized post-report state is $L_v\rho L_v$. At commitment measure $\sigma_z$, obtain $e=\pm1$, and score $Z=y+te e_1$. Under $\omega=I/2$, each $v$ has mass $1/2$ and the normalized posterior is $E_v$. Its $z$ component is zero, so the future audit sign is fair. Thus the reference report law is $\nu_0$ times a fair $v$, and $f=y,w=t^2$, derived directly from the instrument.

The eigenvalues of $E_v$ are $(1\pm\gamma)/2$. Taking $\gamma=1$ on a positive-mass set and $\gamma<1$ on another gives actual pure and full-rank posterior strata. Varying $\phi$ produces noncommuting effects. For a continuous compact realization one may take $\gamma(t)=1$ for $t\le1/2$ and $\gamma(t)=3/2-t$ for $t\ge1/2$, and a continuous nonconstant angle. The two-sided task order remains uniform across these strata; Proposition~\ref{prop:continuity} supplies actual-mass-weighted predictive-update and report continuity. A point mass at $t=0$ also gives a changing-rank task matrix and an exact zero-loss acquired component.

The qubit is part of the physical experiment, not a quantum memory register owned by the finite classical observer. The pulse and instrument form one productive call in this convention; their duration, finite sampling precision and arithmetic are separately charged. An interface requiring a separate pulse call must add its cost and any intervening routing state. Both realizations retain a strictly positive paid gate floor $p_0$.

\subsection*{Predictive realization rather than an assumed parameter manifold}
For the finite classical kernels and finite-dimensional quantum instruments whose entries depend continuously on compact report coordinates, the primitive continuity criterion in Proposition~\ref{prop:raw-quotient} supplies compact positive-evidence history patches. Their executable predictive quotient is therefore Borel by Theorem~\ref{thm:quotient}. Fixed continuous causal evolutions between calls are included and need not contract. For a completely arbitrary merely Borel hidden evolution, the risk theorem still uses the stated raw certificates, but we do not assert the compact-exhaustion quotient theorem without its regularity hypotheses. The unknown model's family of posteriors is never stored by the common upper program.

\section{Finite computation, calibration and causal composition}\label{sec:composition}
The inequalities proved so far concern Borel existence with a declared physical acquisition program. A finite source description and an executable precision certificate are additional mathematical data. We state exactly where they enter.

\begin{proposition}[Numerical implementation at a restored physical budget]\label{prop:numerics}
Suppose a chosen ideal density and dictionary at cap $\beta$ have a finite implementation whose acceptance probability, as a function of the current virtual report, differs pointwise from $h/\beta$ by at most $\eta\ge0$. Suppose its actual exported decoder values differ from the ideal exported values by at most $\delta_{\rm num}$, with no larger output width. These finite operations use the same legal acquisition program and counted live roles. Put $b=N/\kappa_+$ and
\[
 \beta=\frac{l_+b}{1+b(a_++\eta)}\ge1.
\]
If the chosen ideal pair has objective at most $\Q_m(\beta)+\xi$, the implemented program terminates almost surely at actual expected core cost at most $N$ and risk at most
\begin{equation}\label{eq:numerical-law}
 \left[
 \sqrt{\frac{u_+(1+b(a_++\eta))}{l_+}
                  (\Q_m(\beta)+\xi)+b(e_++L\eta)}
       +\delta_{\rm num}\right]^2.
\end{equation}
Here $\xi$ is a certified design/integration suboptimality, not a local optimizer's unverified gap. A uniform bound $T_{\mathrm{attempt}}$ on conditional expected physical and arithmetic time of a whole attempt yields expected time at most $bT_{\mathrm{attempt}}$. Program bits and scratch are separately those of the certified implementation.
\end{proposition}
\begin{proof}
Acceptance rounding reduces the conditional hazard by at most $\eta$. The chosen cap gives $l_+/\beta-a_+-\eta=1/b$. The same conditional-survival summation as in the main theorem yields at most $b$ attempts and cost at most $N$. Before perturbing decoder values, the one-attempt unnormalized loss changes by at most $L\eta$ under the actual joint law, because loss is bounded by $L$. Its other contribution is at most $u_+(\Q_m(\beta)+\xi)/\beta+e_+$. Summation gives the expression under the square root. At the actual committed report, the implemented and ideal exported values are within $\delta_{\rm num}$; Minkowski's inequality in the actual stopped $L^2$ law adds this distance to root-mean-square loss. It does not multiply the decoder error by the waiting time. The time claim is a separate predictable sum of the stated whole-attempt bound.
\end{proof}

In the preceding display $T_{\mathrm{attempt}}$ denotes a declared time quantity, not a free clock input to the observer. An implementation based on $m$ fixed centers in a $k$-dimensional bounded box, $p$-bit coordinate certificates and sequential squared-distance comparisons has sufficient read-only data $O(mkp)$ bits, within-call scratch $O(kp+\log m)$ bits and $O(mkp^2)$ bit operations per report, in addition to its primitive measurement and coordinate-evaluation costs. These are upper accounts, not optimal computational lower bounds. General Borel classifiers need not have such certificates. The complete dyadic implementation in the companion gives explicit rational thresholds, separate read-only/scratch counts and an independently proved profile lower.


\begin{corollary}[Finite-description matching upper under effective primitives]\label{cor:effective}
In the pure relative regular class, suppose additionally that the actual upper-program report laws are bounded above by $\bar u\nu$ on every surviving history. Assume effective compactness of the dictionary carrier and certified finite approximations of exported readouts inside $H$. Assume the calibrated report law has an effective sequence of finite physical partitions which approximates bounded finite decision vectors in $L^1(\nu)$, with certified integration of their masses and bounded task losses. Primitive classifiers and their finite costs are supplied as part of this effective interface. Then a finite-description upper with the same live and terminal role counts has risk at most a fixed multiple of $F(b,m)$ at actual expected cost at most $\kappa_+b$, for all feasible $b,m$. Every such construction has finite, separately charged program, scratch and execution accounts. No polynomial-time discovery or matching bit-complexity lower is asserted.
\end{corollary}
\begin{proof}
For $b\ge4/l_+$ choose the ideal cap $\beta=l_+b/2$. Its conditional acceptance is at least $2/b$. By Proposition~\ref{prop:regular}, its profile is at most a fixed multiple of $F(b,m)$. Represent its finite decision as functions $q_c(x)$, one for each dictionary value, whose sum is the acceptance probability. Approximate the vector on one finite physical partition by rational vectors in the subprobability simplex. The assumed effective approximation makes the sum of its $L^1(\nu)$ coordinate errors arbitrarily small, say at most $\gamma$. Domination of the actual report law bounds its actual conditional decision discrepancy by $\bar u\gamma$ on every history, and its unnormalized loss discrepancy by $L\bar u\gamma$. Approximate the actual finite exported values inside $H$ as well.

Choose $\gamma$ so small that $\bar u\gamma<1/(2b)$ and $bL\bar u\gamma\le F(b,m)$, and choose export error at most $\sqrt{F(b,m)}$. The conditional acceptance still exceeds $1/b$, so the expected cost is at most $\kappa_+b$. Summing loss and applying the stopped $L^2$ readout estimate gives a fixed multiple of $F(b,m)$. For $1\le b<4/l_+$, accept the first attempt and use one finitely described value in $H$. Bounded loss is a fixed multiple of the strictly positive lower bound for $b^{-2/\alpha}$ on this interval, uniformly in $m$.

The existence argument can be effected by finite enumeration using the supplied integration certificates and effective dictionary nets: the above choices leave strict mass and loss margins, so a candidate with certified acceptable intervals is eventually found. This need not be an efficient search. Current partition indices and rational randomization are charged within-call work; they do not create new intercall roles. All chosen descriptions and precision parameters are finite. Exact zero risk on an atom requires an exactly implementable atom classifier and readout; otherwise Proposition~\ref{prop:numerics}, not an unqualified finite-code zero claim, governs that component.
\end{proof}

\begin{proposition}[Resource-certified causal morphisms]\label{prop:composition}
A causal simulator must specify its parameter-independent strategy lift, legal command map, report transformation, internal state, clock and terminal task readout. Suppose a target program has $\ell$ live and $t$ terminal roles. A compatible simulator has at most $K$ continuing internal states, including every persistent phase, coin and clock, and at most $E$ terminal export classes sufficient for its final decoder. Then a source implementation uses at most $\ell K+tE$ state-stopped labels or $\max(\ell K,tE)$ event-stopped labels. Without a terminal sufficient-class certificate, retaining the whole simulator state gives the safe product bound.

If the complete stopped scored transcript, including mapped stopping time and premature commitment, is within total variation $\varepsilon$ of the target law uniformly over the declared models and policies, then the bounded common-task risk changes by at most $L\varepsilon$. Two compatible simulators multiply their continuing and terminal resource factors and add their deficiencies. Raw-call count and physical-time transfer require separate actual cost or killed-potential certificates. Reverse risk inequalities require a reverse simulator or an actual lower witness; terminal-law proximity alone does not provide them.
\end{proposition}
\begin{proof}
Before commitment retain the pair of target live role and continuing simulator state. After commitment retain the pair of target output role and sufficient export class. Compatibility supplies legal transitions and commands without an uncharged port or phase. Lemma~\ref{lem:roles} gives the two packing counts. Serial composition retains triples; rebracketing yields the products $K_1K_2$ and $E_1E_2$. For a loss in $[0,L]$, integration of its superlevel probabilities gives $|\E_P\ell-\E_Q\ell|\le L\|P-Q\|_{\TV}$. Total variation contracts under a common Markov kernel, so the triangle inequality adds complete-law deficiencies. Costs are generally unbounded functions of stopped paths and are not controlled by this bounded-loss argument; sum their certified increments or use Lemma~\ref{lem:drift} instead.
\end{proof}

Combining Propositions~\ref{prop:numerics} and \ref{prop:composition} adds $L\varepsilon$ to \eqref{eq:numerical-law} while charging the simulator's states and actual cost map. This composes approximation and simulation in one task metric. A bound on an allowed error is not, without attained alternatives, an unavoidable lower bound. The following construction supplies such alternatives on a single task.

\begin{theorem}[A same-task resource--calibration--deficiency law]\label{thm:joint}
Take any core class satisfying Theorem~\ref{thm:transfer} with identity exports and $B$ common live roles. Add one compulsory paid observation of an independent fair sign $E_0$, erased with probability $2\varepsilon$, $0\le\varepsilon\le1/2$, and a final paid audit of
\[
 (Z,E_0,s\delta_{\rm cal}),\qquad s\in\{-1,1\}.
\]
The fixed unknown sign $s$ affects no pre-audit law. The core, including its costs, has no access to $E_0$ beyond the initial visible tag. Use squared Euclidean loss on this product task, and charge these two additional calls. The core lower from \eqref{eq:central-lower}, evaluated with budget $N-2$ and safe projected dictionary width $M$, plus $2\varepsilon+\delta_{\rm cal}^2$, is a lower bound for either stopping interface.

Let $U(n,m)$ denote the upper expression in \eqref{eq:central-upper} for core budget $n$ and output width $m$, or its first-attempt alternative. For state stopping choose any $m\ge1$ with $1+3B+3m\le M$; for event stopping choose any $m\ge1$ with $\max(1+3B,3m)\le M$. Then an upper is
\begin{equation}\label{eq:joint-upper}
 U(N-2,m)+2\varepsilon+\delta_{\rm cal}^2.
\end{equation}
In the relative-order regular class, for $M$ large enough that these sufficient widths are a fixed fraction of $M$, this yields
\begin{equation}\label{eq:joint-law}
 R_{\mathrm{joint}}(N,M,\delta_{\rm cal},\varepsilon)
 \asymp F((N-2)/\kappa_-,M)+\delta_{\rm cal}^2+\varepsilon.
\end{equation}
The early sign channel has causal deficiency exactly $\varepsilon$ relative to the perfect-sign channel when the virtual sign must be supplied before any core call. The state counts are sufficient, not minimal for the enlarged experiment.
\end{theorem}
\begin{proof}
For the upper use one initial role, then three tagged copies of the core live roles for observed $+$, observed $-$ and erasure. Use three copies of its output dictionary. Decode the known sign or zero on erasure, and decode zero for the inaccessible last coordinate. This uses $1+3B$ live and $3m$ terminal roles, giving the stated widths by Lemma~\ref{lem:roles}. The common core upper is uniform after each tag. The other losses are exactly $2\varepsilon$ and $\delta_{\rm cal}^2$. The two compulsory calls give core budget $N-2$.

For the lower use the same actual core witness as in the central theorem. Conditioning on the initial tag does not change its full-report certificates; the initial sign and observer coins are additional independent randomization. The projected core decoder has at most $M$ values, so the actual-history proof gives its core lower. On erasure, induction over the core kernels and the stopping decisions leaves $E_0$ fair conditional on all available pre-audit information; its conditional squared prediction loss is at least one. Average the two indistinguishable fixed choices $s=\pm1$. For any last coordinate $r$, their mean loss is $\delta_{\rm cal}^2+r^2\ge\delta_{\rm cal}^2$. The core and erased-sign lower bounds hold under both of these identical pre-audit laws. Thus the three contributions hold simultaneously in one worst-case total task, rather than under incompatible separately optimized models. Regular geometry and fixed width ratios prove \eqref{eq:joint-law}.

On an erased branch every immediate virtual sign is independent of the true fair sign. The event that they agree has probability at most $1/2$, giving complete scored-law deficiency at least $\varepsilon$. Reporting the known sign otherwise and an independent fair sign on erasure gives a coupling with mismatch probability exactly $\varepsilon$. Run the same future core and audit in this coupling. The early-delivery requirement prevents waiting for the final audit to recover the sign, and proves the opposite inequality.
\end{proof}

\subsection*{A finite calibration acquisition with guarded stopping}
The reference and task-order certificates above are not inferred automatically from finite data. Here is one precise finite-sample certificate and its cost when the unknown quantities are the fixed probabilities of a physical command sampler. It also explains why high confidence alone cannot certify almost-sure stopping.

\begin{proposition}[Finite command calibration and a counted guard]\label{prop:pilot}
Suppose a physical sampler has a fixed unknown distribution $(\lambda_j)_{j=1}^J$ throughout calibration and deployment, its draws are independent, and deployment draws are independent of the current hidden input. The known target probabilities are $1/J$. Observe $n$ independent calibration draws at cost $c_{\mathrm{cal}}$ each, and let $\widehat\lambda_j$ be their frequencies. Fix $0<\zeta<1$, and choose a certified rational bound
\[
 r_n\ge\sqrt{\frac{\log(2J/\zeta)}{2n}}.
\]
Declare a relative certificate with parameter $\eta<1$ only when all intervals $[\widehat\lambda_j-r_n,\widehat\lambda_j+r_n]$ lie in $[(1-\eta)/J,(1+\eta)/J]$. The probability of a falsely declared certificate is at most $\zeta$.

On certificate failure, commit by a bounded-loss first-attempt fallback. On certificate success, deploy the corresponding upper of Theorem~\ref{thm:transfer}, with design hazard $1/b$, but force commitment with a fixed decoder after at most $T\ge1$ productive attempts. The whole-attempt cost bound $\kappa_+$ is required whether or not the certificate is true. If the true sampler lies in a tighter relative band $\eta_0<\eta$ and $2r_n\le(\eta-\eta_0)/J$, then the entire guarded procedure commits almost surely and has
\begin{align}
 \E C_{\mathrm{total}}&\le nc_{\mathrm{cal}}+\kappa_+(b+\zeta T),\label{eq:pilot-cost}\\
 R_{\mathrm{total}}&\le U+L\bigl(2\zeta+e^{-T/b}\bigr),\label{eq:pilot-risk}
\end{align}
where $U$ is the certified core upper. The calibration phase can be implemented with at most $(n+1)^J$ count-vector labels. Deployment uses a counted attempt index, multiplying its live-role bound by at most $T$; fallback and terminal roles are additional explicit roles. These resources may not be omitted from an autonomous budget.
\end{proposition}
\begin{proof}
For a Bernoulli variable $X$, the log moment-generating function of $X-\E X$ has second derivative at most $1/4$ under every exponential tilt, and value and derivative zero at the origin. Twice integrating gives $\log\E e^{t(X-\E X)}\le t^2/8$. Independence, the exponential Markov inequality and optimization in $t$ imply $\Prb(|\widehat\lambda_j-\lambda_j|>r)\le2e^{-2nr^2}$. A union bound proves that all confidence intervals contain the true probabilities except on an event of probability at most $\zeta$. Hence a false certificate has probability at most $\zeta$. If the tighter-band condition holds, this same simultaneous event puts each empirical interval within the declared wider band, so rejection of a true certificate also has probability at most $\zeta$.

On a true certificate, the conditional hazard proof bounds survival through $T$ attempts by $(1-1/b)^T\le e^{-T/b}$. The loss contributed by earlier acceptances is at most $U$; the forced last commitment adds at most $Le^{-T/b}$. Expected attempts on this branch are at most $b$. On a certificate failure the first-attempt fallback costs at most $\kappa_+$ and loses at most $L$. On a falsely accepted certificate the deterministic guard uses at most $T$ attempts and has loss at most $L$, even if the unguarded acceptance set has zero actual mass. Every attempt terminates almost surely by its finite whole-attempt cost bound. Combining these cases gives \eqref{eq:pilot-cost} and \eqref{eq:pilot-risk}, with the non-false cases bounded by $b\ge1$ attempts and the false-acceptance case by $\zeta T$ additional attempts.

A count vector stores every calibration total, including the number of draws already taken, and has at most $(n+1)^J$ possible values. At total $n$, the threshold comparison may be performed in charged scratch using a fixed certified finite arithmetic implementation; the resulting branch choice is encoded in the selected live role. The counts can then be discarded because the deployed dictionary is fixed for the nominal law and declared band, not fitted to the realized frequencies. The deployment guard requires its own retained index with $T$ values. This gives the asserted conservative finite-state accounts. Without the tighter-band assumption one still has the false-certificate probability bound and guarded termination, but no small rejection-risk bound; we do not omit that distinction.
\end{proof}

This result controls one genuinely acquired calibration quantity, not arbitrary unknown-kernel learning. It counts pilot observations, a physical fallback, the stopping guard and its memory. It prevents a rare incorrect calibration certificate from being hidden inside an unjustified robust almost-sure stopping assertion.

\appendix
\section{A predictive quotient constructed from positive-evidence histories}\label{sec:quotient}
The squared-task coordinates in \eqref{eq:derived} need not separate all future experiments. We therefore give the predictive realization independently of those coordinates, for a fixed prepared law with its declared legal test class. The theorem is a sufficient realization result and a detailed application of compact selection; predictive minimality itself has classical antecedents \cite{SC,BC}.

Let $\mathsf H$ be the standard Borel space of finite visible histories on which specified conditional versions are defined, embedded in a metrizable history carrier. Reports and resource records are included in the history. Actions belong to a finite set, with Borel legal domains. Histories with different legal continuations must be separated. Let $q_j(h)\in[0,1]$ be a countable determining list of conditional response coordinates for all finite open-loop future report words and bounded tests. Include legal-domain indicators and the coordinates necessary to determine report and cost-record laws. The list is closed in the following determining sense: the one-call prefixes of these future tests, localized by a countable measure-determining family of current-report functions, are determined by the list. This closure is a property of the declared executable tests, not a finite-moment assertion. Set $q(h)=(q_j(h))_{j\ge1}$ in $[0,1]^{\mathbb N}$.

The realization theorem in this appendix concerns one fixed prepared law, not an unknown-model quotient oracle. Its compact-fiber selection step is a concrete instance of classical measurable selection \cite{KRN65}; the explicit construction below checks its needed measurable domains instead of assuming a selector.

\begin{theorem}[Compact-exhausted executable predictive realization]\label{thm:quotient}
Suppose $\mathsf H=\bigcup_{n\ge1}C_n$ with $C_n$ compact in its metrizable carrier and every $q_j$ continuous on each $C_n$. Assume the specified response coordinates are the consistent conditional laws of the raw causal kernels and satisfy the determining and continuation conditions above. Then the actual image $S=q(\mathsf H)$ is Borel and has a Borel section $s:S\to\mathsf H$. Equality of $q$ is exactly equality of all declared finite executable future laws, including legal domains. The report and next-quotient law factor through $q$ as a Borel kernel. This factorization is modulo the common current-report null sets of each compared pair of histories; it is not a pointwise claim on impossible reports.

If another Borel statistic $T$ determines the response coordinates by Borel maps $q_j=g_j\circ T$, then $q=g\circ T$ measurably, where $g=(g_j)_{j\ge1}$ on the realized range. Thus $q$ is minimal for these executable tests. Equality extends to bounded stopped tests which terminate almost surely under both compared histories, by finite truncation. Unbounded costs and infinite-horizon limits require their separate integrability or convergence hypotheses.
\end{theorem}
\begin{proof}
The countable product is compact metrizable, and coordinate continuity makes $q$ continuous on every $C_n$. Thus $q(C_n)$ is compact and $S=\bigcup_nq(C_n)$ is Borel. A section can be constructed rather than assumed. Given $x\in S$, first choose the least $n$ for which $x\in q(C_n)$, a Borel decision. For this compact $C_n$, choose finite closed covers at each scale, refining each previously chosen compact cell, with diameters tending to zero. Empty cells are discarded. At each level choose the first child compact cell $D$ for which $x\in q(D)$. Such a child exists because its parent covers the nonempty fiber; the test is Borel because $q(D)$ is compact. The chosen cells are nested compact sets with diameters tending to zero and each intersects the fiber. Choose a fixed representative in each possible cell. Their sequence is Borel in $x$ and converges to a unique point of $C_n$. Continuity of $q$ on $C_n$ shows that this point maps to $x$. The countable piecewise construction in the least $n$ gives the required Borel section.

Suppose $q(h)=q(h')$. The legal-domain coordinates give the same permitted next commands. The determining current-report coordinates give the same next-report law $R_a$ for a common legal command $a$. For a future response coordinate $q_j$ and a member $g$ of the determining current-report function family, continuation consistency gives a conditional joint-test value
\[
 \int g(y)q_j(hay)R_a(dy).
\]
The prefix closure and equality of $q(h),q(h')$ imply equality of this integral with the expression using $h'ay$. Uniqueness of finite measures on a determining family implies $q_j(hay)=q_j(h'ay)$ for $R_a$-almost every $y$. Intersecting these sets over the countable coordinates gives $q(hay)=q(h'ay)$ almost surely. Concatenations outside the positive-evidence history domain may be completed arbitrarily, because they have zero report mass from the histories in question.

Define the quotient report/update kernel by running the raw row from $s(x)$ and pushing the continued history through $q$. The section, the raw kernel and concatenation map are Borel, so this is a Borel kernel. The preceding almost-sure equality shows that it agrees with the pushforward from every representative of $x$, not just from $s(x)$. Induction gives finite-prefix equality under a common quotient policy. More generally finite adaptive policies depending on the visible future history are handled by the same raw conditional continuation laws: with finitely many actions, sum over action words and multiply their open-loop report laws by the measurable feedback probabilities of the policy. Equality extends from bounded determining functions to Borel sets by the monotone-class argument. The converse follows since each listed response is itself a declared test coordinate.

For another statistic with the stated Borel response maps, the product map $g$ is Borel and satisfies $q=g\circ T$; the asserted minimality is precisely this factorization, not a finite-state claim. For a stopped test truncate at time $n$ with a distinguished surviving outcome. Its finite law is determined. Almost-sure termination under both histories makes both surviving probabilities tend to zero, so the bounded stopped law agrees. This limit does not control expected stopping cost without an additional bound.
\end{proof}

\begin{proposition}[Primitive positive-evidence patches]\label{prop:raw-quotient}
The continuity and exhaustion conditions of Theorem~\ref{thm:quotient} hold for the following raw class. Hidden-state and report coordinates lie in compact metric spaces, the action set is finite, and primitive transition--report kernels have bounded continuous densities with respect to fixed finite reference measures; the preparation is fixed. A finite-dimensional positive instrument whose matrix entries depend continuously on compact report coordinates is another such class. Finite physical phases are included as discrete coordinates, with continuous legal-domain indicators on the corresponding patches. Only histories of positive prepared likelihood are included.
\end{proposition}
\begin{proof}
Fix a finite action word and history length. Iterated integration of products of bounded continuous primitive densities gives a continuous nonnegative history likelihood $p(h)$. The same statement holds for the numerator of a conditional finite future test whose report test function is continuous; the integration domains are fixed compact spaces with finite measures. Uniform boundedness on each finite product and dominated convergence justify continuity. In the matrix case these quantities are traces of products of continuous instrument maps followed by finite integration, with the same conclusion.

The sets where $p(h)\ge1/n$ are compact and exhaust the positive-likelihood histories for this action word. On each, division by $p(h)$ is continuous and all conditional response coordinates are continuous. Countably many lengths, finitely many action words per length, and countably many $n$ give the compact exhaustion. On a compact metric report space there is a countable uniformly dense family of continuous functions; include their finite products and all finite action words. Their integrals determine finite report laws, and adjoining finite prefixes gives the required determining closure. Costs and legal phases must be included among the stipulated continuous or discrete observable coordinates. The probability of a zero-likelihood report history is zero under the prepared law. The construction does not claim that the likelihood has a uniform positive lower bound over the whole experiment.
\end{proof}

The proposition is about continuity of raw finite-word responses. It does not assume a pre-existing posterior manifold or a covariance geometry, and it does not give a finite algorithm for discovering all quotient classes. An unknown-model family of separate posterior laws is not automatically one executable posterior. Applying this appendix uniformly to a model family would require a common determining response construction and compatible null-set/regularity conditions; the finite-state upper in the main theorem instead obtains its uniformity from the explicitly checked task-order program and stores no such family.

\section{Sharp boundaries of the proved comparisons}\label{sec:boundaries}
The following witnesses isolate distinct hypotheses. They are mathematical counterexamples to specific uncharged inferences, not a replacement of the main theorem by an impossibility statement.

\begin{proposition}[Absolute task error can be amplified by selection]\label{prop:absolute-sharp}
For every $b\ge1$ and $0\le e\le1/b$, there are positive prepared and actual joint laws with identical report marginals and one-call absolute squared-task defect at most $e$, whose optimal stopped risks at expected call budget $b$ differ by $be$, even with one output value. The difference occurs in either direction. A fixed relative task order is therefore strictly stronger than a small absolute task defect near a zero-loss component.
\end{proposition}
\begin{proof}
Draw an independent uniform report $x\in[0,1]$ at every call and let $D$ have mass $1/b$. All audits have conditional mean zero. Off $D$, both laws have a fresh fair-sign audit of magnitude one. On $D$, the comparison has a fair sign of magnitude $\sqrt{be}$ and the actual law has audit zero. For every $q\in[0,1]$ and finite-image readout $r(x)$, their unnormalized squared losses differ by $be\int_Dq(x)dx\le e$. The comparison residual variance is everywhere at least $be$, so every committing strategy has risk at least $be$; selecting $D$ and outputting zero attains it with expected cost $b$. The actual law attains zero by the same procedure. Reversing the laws proves the other direction. A single event-stopped label can reject at itself and commit with a final event; a state-stopped realization uses an additional live label. On $D$ a task zero under one law is positive under the other, so no two-sided relative order with a positive lower factor can hold.
\end{proof}

A report-only comparison also fails: a constant mark can have the same marginal under two laws while an additional visible bit reveals a fair audit sign. A two-output decoder then has zero loss, although a profile formed from the constant mark alone has residual variance one. The full-report conditional moments in \eqref{eq:full-moments} exclude this omission.


There is also an all-program witness for the necessity of more than report balance. Take the three admitted models $\theta\in\{-1,0,1\}$, let every precommit report be the same constant, and let the unseen audit at commitment be $Z=\theta$. There is one legal productive command, costing one, and no command changes $\theta$. The actual comparison member $\theta=0$ has $\Q_m(b)=0$. Nevertheless the distribution of every committed readout $R$ is identical under $\theta=\pm1$, because all precommit laws, including stopping and coins, are identical. The average of their risks is $1+\E R^2\ge1$, and output zero attains worst-case risk one. This holds for every almost-surely terminating policy, at every feasible budget and every dictionary size. Report masses match exactly; the missing first task moment is the obstruction. It is not repaired by choosing a different causal rejection rule.

A live-state count alone cannot justify a terminal-width subtraction. In the two-port interface, use labels $0,1$ with live commands gate and productive respectively, and decoder values $-1,+1$. Gate failure continues at $0$, success at $1$; a rejected productive report returns to $0$. A productive commitment emits a final event with label $0$ or $1$ according to its chosen output. Two live and two output roles have used only two labels. The state-stopped model cannot share those roles because its halt map is intrinsic. This is the exact semantic repair retained from the earlier work.

Finally, bounded-law discrepancy is not a cost certificate. A source procedure can agree with a target terminal output with probability one while performing an independent preliminary waiting time of arbitrarily large finite mean, or of infinite mean. Terminal total variation is then zero and expected cost is uncontrolled. Our simulator definition includes the scored stopping information when its law is compared, and still requires a separate expected-cost certificate because time is unbounded. Likewise, neither fixed-time equivalence nor finite-dimensional positive matrix identities imply large-deviation equivalence, unrestricted optimal-control equivalence, or closure of unbounded operators.

\subsection*{What is and is not generalized}
The central theorem generalizes exact and absolute-calibrated causal profile transfer to a common affine task order and derives its all-scale relative and stratified consequences. The positive-matrix characterization concerns localized task order; the quotient construction concerns compact-exhausted continuous-response experiments; the pilot concerns a fixed independent command sampler. These results do not establish necessity of a common acquisition program for every causal experiment, nor the minimal possible simulator state under all implementations, nor matching lower bounds for program length, scratch and arithmetic time. Those stronger mathematical questions are distinct from the fully stated sufficient-class theorem proved here. The general theorem does not depend on a conclusion taken from either realization. The retained nonintervening sensing, occupation-flow, exact role and absolute-defect results remain substantive independent parts of the complete companion.

\section{The scope of the strengthened theory}\label{sec:current-scope}
The boundaries in Appendix~\ref{sec:boundaries} concern the inherited causal transfer theorem. The present synthesis theorems add necessary-and-sufficient conclusions on explicitly stated classes; they do not turn that sufficient transfer theorem into a classification of all standard Borel causal experiments.

The task-defect theorem is exact for a fixed selected loss family, its declared decision carrier, and one global finite decoder. Its occupation version holds for an actual occupation measure supplied by a competitor. Neither assertion manufactures that measure or an acquisition program. The scalar squared-task order is strictly weaker than the affine-square matrix order; the latter remains an exact characterization of its own larger test family. The Bregman comparison preserves squared-loss exponents only under the stated curvature bounds.

The finite-obstruction hierarchy characterizes one common full-history program for finitely branching instruments, bounded terminal scores and a strictly positive call-cost floor. It permits arbitrarily many operating models, but finite testing of an arbitrary uncountable family is not a decision procedure for its infinite synthesis problem. Compactness produces a policy, not a uniformly bounded number of live labels. Its finite dual concerns one specified horizon and finite model set. Fallback completion requires a genuinely legal common continuation and charges its cost and memory. The current supplement treats deadlocks, restricted preparation faces and the distinction between small survival mass and small expected cost tails.

The finite-memory certificate is intrinsic to a fixed finite-label observer on a declared time-homogeneous interface. Its projections and potential matrices are proof certificates, not stored observer variables. The finite algebraic decision conclusion requires finitely many models, algebraic primitive data and a semialgebraic decision carrier. It does not assert efficient synthesis, rational optimal probabilities, an exact finite-bit implementation of arbitrary algebraic randomization, or optimal program, scratch and arithmetic lower bounds. Exact initial-state uniformity does not replace a missing certificate at later roots of a history-dependent experiment. Returning from a balancing subroutine must be compiled at the current productive report before its prospective audit, as specified in the supplement.

The two-detector theorem has no balancing group and no free reset. It characterizes all integer classical live-memory budgets for two projectors, unit call costs, first-success commitment and no other physical command. The maximum principal-angle cosine is the worst spectral block of that model, not a universal rank or memory exponent. The initial-state saddle at unrestricted memory is proved from the physical instrument. The exact two-label minimax gap shows that a single preparation need not supply a sharp robust converse after imposing a nonconvex controller restriction. It does not imply such a gap for every memory-constrained experiment. Randomized initial labels are a separately stated interface, not an uncharged repair of a fixed-start observer.

All inherited valid results remain either in the native main inputs, the complete preceding companion, or their pinned independent realization lineages. In particular, the ordered-measurement local geometry and allocation theory, exact two-stage reset and message-discrimination theory, rank constraints and spectral variational values are not asserted to follow from the detector example. Conversely the current article does not abandon them because its general structural results have become central. The accompanying preservation and dependency maps identify their locations and prevent a change of research identity from being hidden in a reorganization.

Open mathematical goals remain: an intrinsic classification beyond finite branching; minimal memory under arbitrary resource-certified simulators; a broad nonintervening controlled-sensing application beyond the retained directional model; matching computational-resource lower bounds; and the independent unbounded-operator, large-deviation and wider A1/A2 pipeline targets. These are not hypotheses standing in for a missing proof of a theorem stated here. The present proofs establish their displayed conclusions with their displayed quantifiers. Finite diagnostics test particular algebraic identities, boundary cases and implementations; they are not formal proofs, independent expert assessment or an editorial decision on novelty and significance.



The report-resolved continuation theorem is exact for linear marked task equalities of a supplied time-homogeneous finite-return routine. It does not identify the reachable positive cone with every positive matrix on its support, nor replace a one-sided reachable cost bound by an unsupported operator inequality. The reset-instrument classification applies to failure branches which physically prepare known states and to first-success commitment. It includes arbitrary finite rank-one families, but not arbitrary quantum instruments or an arbitrary-rank projection without input-erasing failure. The zero-risk spectrum is exact for the acquired profile; its physical threshold has the same converse and implementation hypotheses as the exact transfer theorem. These distinctions separate the new structural results from the independent wider research goals.


The primitive finite-model route in Section~\ref{sec:primitive-risk} removes the reset and reference assumptions on its own declared interface. It supplies a common fallback, monotone two-sided approximations to the least feasible cost, and bounded-risk approximations at every strictly interior feasible budget with counted finite implementations. It does not replace the preparation-free all-initial-state quantifier by an arbitrary restricted preparation set, and it assumes finite models, finite globally legal commands/reports, time homogeneity and bounded positive call costs. It is not unknown-kernel statistical learning, a universal closed formula for the comparison profile, or optimal computational complexity. The six-regime diagram solves one explicit mixed-strata comparator with its original asymptotic qualification intact.


The compact-family route extends the primitive finite-model risk approximation to a continuously parameterized compact uncertainty class on a common finite globally legal physical interface. It controls model error and restores the same budget at every $N>N_{\min}$. Its algebraic specialization is a terminating approximate synthesis result for supplied exact primitive data, not learning from finite data, polynomial-time synthesis or a solution of the boundary finite-memory problem. The model stays fixed during a run. Exact geometric reference-profile realization, arbitrary unknown-model quotient construction and minimal computational resources remain separate mathematical questions. The examples showing boundary discontinuity and failure of vertex-only certification delimit the positive theorem rather than replace it.

\begin{thebibliography}{99}
\bibitem{Blackwell} D. Blackwell, \emph{Equivalent comparisons of experiments}, Ann. Math. Statist. \textbf{24} (1953), 265--272. DOI: 10.1214/aoms/1177729032.
\bibitem{CGM97} J. A. Cuesta-Albertos, A. Gordaliza and C. Matr\'an, \emph{Trimmed $k$-means: an attempt to robustify quantizers}, Ann. Statist. \textbf{25} (1997), 553--576. DOI: 10.1214/aos/1031833664.
\bibitem{GG99} L. A. Garc\'ia-Escudero and A. Gordaliza, \emph{Robustness properties of $k$ means and trimmed $k$ means}, J. Amer. Statist. Assoc. \textbf{94} (1999), 956--969. DOI: 10.1080/01621459.1999.10474200.
\bibitem{SC} C. R. Shalizi and J. P. Crutchfield, \emph{Computational mechanics: pattern and prediction, structure and simplicity}, J. Statist. Phys. \textbf{104} (2001), 817--879. DOI: 10.1023/A:1010388907793.
\bibitem{BC} N. Barnett and J. P. Crutchfield, \emph{Computational mechanics of input--output processes: structured transformations and the $\varepsilon$-transducer}, J. Statist. Phys. \textbf{161} (2015), 404--451; arXiv:1412.2690.
\bibitem{AMTW} A. Ambainis, M. Mosca, A. Tapp and R. de Wolf, \emph{Private quantum channels}, Proc. 41st Annual Symposium on Foundations of Computer Science (2000), 547--553. DOI: 10.1109/SFCS.2000.892142.
\bibitem{Buscemi} F. Buscemi, \emph{Comparison of quantum statistical models: equivalent conditions for sufficiency}, Comm. Math. Phys. \textbf{310} (2012), 625--647. DOI: 10.1007/s00220-012-1421-3.
\bibitem{BR14} N. Bertschinger and J. Rauh, \emph{The Blackwell relation defines no lattice}, Proc. IEEE International Symposium on Information Theory (2014), 2479--2483. DOI: 10.1109/ISIT.2014.6875280; arXiv:1401.3146.
\bibitem{Banerjee05} A. Banerjee, X. Guo and H. Wang, \emph{On the optimality of conditional expectation as a Bregman predictor}, IEEE Trans. Inform. Theory \textbf{51} (2005), 2664--2669. DOI: 10.1109/TIT.2005.850145.
\bibitem{Stengel96} B. von Stengel, \emph{Efficient computation of behavior strategies}, Games Econom. Behav. \textbf{14} (1996), 220--246. DOI: 10.1006/game.1996.0050.
\bibitem{Balder89} E. J. Balder, \emph{On compactness of the space of policies in stochastic dynamic programming}, Stochastic Process. Appl. \textbf{32} (1989), 141--150. DOI: 10.1016/0304-4149(89)90058-6.
\bibitem{Barry14} J. Barry, D. T. Barry and S. Aaronson, \emph{Quantum partially observable Markov decision processes}, Phys. Rev. A \textbf{90} (2014), 032311. DOI: 10.1103/PhysRevA.90.032311; arXiv:1406.2858.
\bibitem{Ying18} S. Ying and M. Ying, \emph{Reachability analysis of quantum Markov decision processes}, Inform. Comput. \textbf{263} (2018), 31--51. DOI: 10.1016/j.ic.2018.09.001; arXiv:1406.6146.
\bibitem{Saldi25} N. Saldi, S. Sanjari and S. Y\"uksel, \emph{Quantum Markov decision processes: general theory, approximations, and classes of policies}, SIAM J. Control Optim. \textbf{63} (2025), 1156--1183. DOI: 10.1137/24M1643864; arXiv:2402.14649.
\bibitem{Seidenberg54} A. Seidenberg, \emph{A new decision method for elementary algebra}, Ann. of Math. (2) \textbf{60} (1954), 365--374. DOI: 10.2307/1969640.
\bibitem{Chatterjee15} K. Chatterjee, M. Chmel\'ik, R. Gupta and A. Kanodia, \emph{Optimal cost almost-sure reachability in POMDPs}, Proc. AAAI Conference on Artificial Intelligence \textbf{29} (2015). DOI: 10.1609/aaai.v29i1.9683.
\bibitem{Andriushchenko23} R. Andriushchenko, A. Bork, M. \v{C}e\v{s}ka, S. Junges, J.-P. Katoen and M. Mac\'ak, \emph{Search and explore: symbiotic policy synthesis in POMDPs}, Computer Aided Verification (2023); arXiv:2305.14149v2.
\bibitem{Halmos69} P. R. Halmos, \emph{Two subspaces}, Trans. Amer. Math. Soc. \textbf{144} (1969), 381--389. DOI: 10.1090/S0002-9947-1969-0251519-5.
\bibitem{Kessler21} D. A. Kessler, E. Barkai and K. Ziegler, \emph{First-detection time of a quantum state under random probing}, Phys. Rev. A \textbf{103} (2021), 022222. DOI: 10.1103/PhysRevA.103.022222; arXiv:2012.01763.
\bibitem{DelVecchio26} G. Del Vecchio Del Vecchio and S. N. Majumdar, \emph{Optimal detection of quantum states via projective measurements}, J. Phys. A: Math. Theor. \textbf{59} (2026), 035001. DOI: 10.1088/1751-8121/ae34fe; arXiv:2509.08556.
\bibitem{Companion} Q. Qi, \emph{General Theta Foundations I: Acquired Geometry and Causal Resource Transfer}, complete technical companion on absolute calibration, exact stopping resources and nonintervening sensing, 2026. The complete retained R58 native source and its independently built companion accompany this manuscript; their identities are fixed in the source-provenance and preservation records.
\bibitem{Tzeng92} W.-G. Tzeng, \emph{A polynomial-time algorithm for the equivalence of probabilistic automata}, SIAM J. Comput. \textbf{21} (1992), 216--227. DOI: 10.1137/0221017.
\bibitem{LiYuYing12} Y. Li, N. Yu and M. Ying, \emph{Termination of nondeterministic quantum programs}, arXiv:1201.0891 (2012).
\bibitem{Fu24} J. Fu, H. Jiang, M. Xu, Y. Deng and Z.-B. Li, \emph{Algorithmic analysis of termination problems for nondeterministic quantum programs}, arXiv:2402.15827 (2024).

\bibitem{LiuVerdu18} J. Liu and S. Verd\'u, \emph{Rejection sampling and noncausal sampling under moment constraints}, Proc. IEEE International Symposium on Information Theory (2018), 1565--1569. DOI: 10.1109/ISIT.2018.8437857.
\bibitem{GocFlamich24} D. Goc and G. Flamich, \emph{On channel simulation with causal rejection samplers}, Proc. IEEE International Symposium on Information Theory (2024). DOI: 10.1109/ISIT57864.2024.10619339; examined text arXiv:2401.16579v2 (3 May 2024).
\bibitem{LiLiu26} Z. Li and T. Liu, \emph{Ranking experiments under sequential sampling}, arXiv:2608.19897v2 (28 August 2026). Version 1 appeared under the title \emph{The order of binary experiments under endogenous stopping}, with Z. Li as sole author; the version-2 statement is the one compared here.
\bibitem{Ying14pre} S. Ying and M. Ying, \emph{Reachability analysis of quantum Markov decision processes}, arXiv:1406.6146 (2014), Theorem 3.6; published journal version listed above as \cite{Ying18}.

\bibitem{KaraYuksel20} A. D. Kara and S. Y\"uksel, \emph{Robustness to incorrect system models in stochastic control}, SIAM J. Control Optim. \textbf{58} (2020), 1144--1182. DOI: 10.1137/18M1208058; examined text arXiv:1803.06046v7.
\bibitem{Demirci25} Y. E. Demirci, A. D. Kara and S. Y\"uksel, \emph{Sensitivity of filter kernels and robustness bounds to transition and measurement kernel perturbations in partially observable stochastic control}, arXiv:2508.10658v3 (2025).
\bibitem{Watrous18} J. Watrous, \emph{The Theory of Quantum Information}, Cambridge University Press, 2018, Chapter 3, Corollary 3.40.
\bibitem{KRN65} K. Kuratowski and C. Ryll-Nardzewski, \emph{A general theorem on selectors}, Bull. Acad. Polon. Sci. S\'er. Sci. Math. Astronom. Phys. \textbf{13} (1965), 397--403.
\end{thebibliography}

\end{document}
